\documentclass[11pt]{article}
\usepackage[margin=1in]{geometry}
\usepackage[utf8]{inputenc}
\usepackage[T1]{fontenc}
\usepackage{amsmath,amssymb,amsfonts,amsthm,mathtools}
\usepackage{booktabs}
\usepackage{microtype}
\usepackage{enumitem}
\usepackage{xcolor}
\usepackage{tikz}
\usepackage[colorlinks=true,linkcolor=blue!60!black,citecolor=blue!60!black,urlcolor=blue!60!black]{hyperref}
\usepackage[round,authoryear]{natbib}

\theoremstyle{plain}
\newtheorem{theorem}{Theorem}[section]
\newtheorem{lemma}[theorem]{Lemma}
\newtheorem{corollary}[theorem]{Corollary}
\newtheorem{proposition}[theorem]{Proposition}
\newtheorem{conjecture}[theorem]{Conjecture}
\theoremstyle{definition}
\newtheorem{definition}[theorem]{Definition}
\newtheorem{example}[theorem]{Worked example}
\newtheorem{problem}{Problem}[section]
\theoremstyle{remark}
\newtheorem{remark}[theorem]{Remark}

% exercise-solution environment used in the fragments
\newenvironment{solution}{\par\smallskip\noindent\textit{Solution.}\;}{\hfill$\square$\par\smallskip}

\newcommand{\R}{\mathbb{R}}
\newcommand{\Z}{\mathbb{Z}}
\newcommand{\E}{\mathbb{E}}
\newcommand{\one}{\mathbf 1}
\newcommand{\sign}{\operatorname{sign}}
\newcommand{\conv}{\operatorname{conv}}
\newcommand{\dist}{\operatorname{dist}}
\newcommand{\Lip}{\operatorname{Lip}}
\newcommand{\Orb}{\mathcal O}
\newcommand{\fdc}{f_{\mathrm{dc}}}
\newcommand{\Vinv}{V_{\mathrm{inv}}}
\newcommand{\Vac}{V_{\mathrm{ac}}}
\newcommand{\Sep}{\operatorname{Sep}}
\newcommand{\Piinv}{\Pi_{\mathrm{inv}}}

\title{\textbf{Shift-Invariance and Adversarial Robustness}\\[3pt]
\large A Collaborator Primer\\[4pt]
\normalsize background from first principles, the questions before the answers, and full proofs}
\author{Anass Al Ammiri}
\date{\today}

\begin{document}
\maketitle

\begin{center}\small\itshape
A standalone booklet for a new collaborator. It assumes only undergraduate linear algebra, real analysis, basic
calculus, and probability, and builds everything else from first principles. Part~I teaches the background, Part~II
poses the open problems (before any answers), and Part~III proves the results. Do the exercises; attempt the
Part~II problems before reading Part~III.
\end{center}

\medskip\hrule\medskip
\tableofcontents
\newpage

%==============================================================================
\section{Why this project exists}
\label{sec:why}
Convolutional networks are built to be (approximately) invariant to small spatial shifts, and for ordinary accuracy
this is a virtue. Whether it is a virtue for \emph{adversarial} robustness, the worst-case behaviour under tiny
crafted perturbations, is genuinely unsettled: published work points both ways. \citet{ge2021shift} prove a fully
shift-invariant classifier can be far less robust than a plain one; \citet{kamath2021} report a spatial-versus-adversarial
trade-off; \citet{galloway2019} blame a standard invariance-promoting ingredient (batch normalisation) for
vulnerability; while a line of architectural work \citep{zhang2019blurpool,sahagokhale2024} improves shift behaviour
and reports robustness \emph{gains}. These are not all wrong; they are measuring different things. The thread of this
project, and of this booklet, is that one scalar, the margin of the surviving invariant signal divided by its
sensitivity, reconciles them, and that a careful capacity-matched experiment under a strong attack tells the rest of
the story. The full chain that produced the project's questions, starting from a practical course that reproduced and
then broke the clean story, is the subject of Part~II; the rest of this section just fixes expectations.

\paragraph{How to read this.} Part~I is the background you need and nothing more, in dependency order, with a worked
example and an exercise per topic. Part~II is the harder skill: given the situation, generating and choosing the
research questions yourself, so it deliberately withholds the answers. Part~III gives the answers with full proofs,
each labelled honestly as classical, a rephrasing, or a genuine new contribution, and each tied back to the question
it settles.
\newpage

%==============================================================================
\part*{Part I: Background from first principles}
\addcontentsline{toc}{section}{Part I: Background from first principles}
This part assumes only linear algebra, real analysis, basic calculus, and probability. It builds, in order, the
robustness vocabulary, the invariance and Fourier machinery, and the learning theory the project rests on. Every
term is defined on first use, every central claim is proved in full, and each topic ends with a worked example and
an exercise. Read it in order; do the exercises.

% =====================================================================
%  PRIMER PART 1 — ROBUSTNESS CLUSTER
%  Three background sections, taught from scratch for a strong
%  undergraduate (linear algebra, real analysis, calculus, probability;
%  NO deep learning). Grounded in literature/pedagogy/ sources.
%
%  This is a FRAGMENT: no \documentclass, no \begin{document}.
%  The assembler (primer.tex) supplies the preamble, which already
%  defines: definition, theorem, lemma, proposition, corollary,
%  example (printed as "Worked example"), remark, and the macros
%  \R \Z \E \sign \Lip \one \conv \dist etc.  Exercises use
%  \paragraph{Exercise.}  (no custom environment needed).
%
%  Citation keys used appear in the comment block at the END of file.
% =====================================================================

\section{Adversarial robustness from scratch}
\label{sec:adv-from-scratch}

This section answers one question: \emph{what does it mean for a classifier to be robust, and why is the obvious thing (high test accuracy) not enough?} We follow the order that worked for the three undergraduates who got lost in the previous draft: first a surprising concrete demonstration, then just enough notation to talk about it, and only at the end the formal optimization picture. This is the order of the Kolter--Madry tutorial \citep{kolter2018adversarial}, which is the spine of this whole section.

\subsection{Demo first: a pig that the classifier thinks is a wombat}
\label{sub:demo}

Imagine you have a state-of-the-art image classifier, a fixed function that takes an image and outputs a label. You feed it a clear photo of a pig. It says ``pig'' with $99.6\%$ confidence. Good. Now you add to that photo a tiny amount of carefully chosen noise: so tiny that, pixel by pixel, no human eye can tell the new image apart from the old one. You feed in the new image. The same classifier now says ``wombat'' with $99.98\%$ confidence \citep{kolter2018adversarial}. With slightly different noise it says ``airliner.'' The pictures look identical to you; the classifier's answer flipped completely.

This is not a bug in one model or one photo. It happens to essentially every modern image classifier, on essentially every image, and the noise needed is astonishingly small. The phenomenon was discovered by \citet{szegedy2014intriguing}, who called these crafted inputs \emph{adversarial examples}. The rest of this section builds the language to say precisely what just happened and how we measure it.

\paragraph{Why should you care?} Two reasons, both from \citet{kolter2018adversarial}. First, the obvious security reason: if a system is deployed where someone has an incentive to fool it (spam filters, malware detectors, a camera that reads stop signs for a self-driving car), then ``works most of the time'' is not the right standard; we want ``works even against someone deliberately trying to break it.'' Second, the deeper reason: a model that can be flipped from ``pig'' to ``airliner'' by invisible noise clearly has \emph{not} learned the concept ``pig'' the way a human has. Adversarial examples expose that high test accuracy can hide a brittle, shallow understanding of the data.

\subsection{The minimum notation we need}
\label{sub:notation}

\begin{definition}[Model / hypothesis]
A \emph{model} (or \emph{hypothesis}) is a function $h_\theta : \mathcal X \to \R^k$. It takes an input $x$ from the input space $\mathcal X$ (for images, a long vector of pixel intensities) and outputs a vector of $k$ real numbers, one per class. The subscript $\theta$ collects all the model's tunable numbers, called \emph{parameters} (for a neural network these are all its weight matrices and biases). We \emph{train} a model by choosing $\theta$; once trained, $\theta$ is frozen and we study $x$.
\end{definition}

The $k$ output numbers are called \emph{logits}.

\begin{definition}[Logit, softmax, predicted class]
The components $h_\theta(x)_1,\dots,h_\theta(x)_k$ are the \emph{logits}: real numbers, possibly negative, one per class. The predicted class is the index of the largest logit, $\arg\max_i h_\theta(x)_i$. To turn logits into probabilities one applies the \emph{softmax} map $\sigma:\R^k\to\R^k$,
\[
  \sigma(z)_j \;=\; \frac{e^{z_j}}{\sum_{i=1}^{k} e^{z_i}},
\]
which is positive and sums to $1$. Bigger logit $\Rightarrow$ bigger probability; the $\arg\max$ is unchanged by softmax, so for deciding the label we can work with the raw logits and ignore softmax.
\end{definition}

\begin{definition}[Loss]
A \emph{loss function} $\ell(h_\theta(x),y)$ takes the model's logit vector and the true class index $y\in\{1,\dots,k\}$ and returns a single nonnegative number measuring how wrong the prediction is: small loss means the model put high probability on the true class. The standard choice is the \emph{cross-entropy} (or \emph{softmax}) loss
\[
  \ell(h_\theta(x),y) \;=\; -\log \sigma\!\big(h_\theta(x)\big)_y
  \;=\; -\,h_\theta(x)_y + \log\!\sum_{i=1}^k e^{h_\theta(x)_i},
\]
i.e.\ the negative log-probability assigned to the true class. In the pig example the loss of the correct ``pig'' label was about $0.0039$, which corresponds to probability $e^{-0.0039}\approx 0.996$ \citep{kolter2018adversarial}.
\end{definition}

\paragraph{Gradient with respect to the input.} The single technical tool that makes adversarial examples possible is the \emph{gradient of the loss with respect to the input}, written $\nabla_x \ell(h_\theta(x),y)$. Recall from calculus that for a scalar function $\ell$ of a vector $x\in\R^n$, the gradient $\nabla_x \ell$ is the vector of partial derivatives $\big(\partial \ell/\partial x_1,\dots,\partial \ell/\partial x_n\big)$; it points in the direction in input space along which the loss increases fastest, and its $i$-th entry tells you how much the loss changes if you nudge pixel $i$. Crucially, \emph{normal} training computes the gradient with respect to the parameters $\theta$ (to improve the model); here we keep $\theta$ fixed and differentiate with respect to the \emph{input} $x$ (to attack the model). Both are computed by the same mechanical procedure (backpropagation); for our purposes you may treat $\nabla_x \ell$ as a vector you can evaluate at any point.

\subsection{An adversarial example is an inner maximization}
\label{sub:inner-max}

Training minimizes the loss over $\theta$. An attacker does the opposite: keep the model fixed and \emph{change the input to make the loss large}. In the words of \citet{kolter2018adversarial}: ``instead of adjusting the image to minimize the loss\dots we're going to adjust the image to maximize the loss.''

We cannot let the attacker change the image arbitrarily, turning a pig into an actual photo of a wombat would not be impressive. So we restrict the change. Write the perturbation as $\delta$ and require it to lie in an allowed set $\Delta$:
\begin{equation}
  \label{eq:inner-max}
  \max_{\delta \in \Delta}\; \ell\!\big(h_\theta(x+\delta),\, y\big).
\end{equation}
Equation~\eqref{eq:inner-max} is the heart of the whole field. In words: \emph{over all allowed perturbations $\delta$, what is the largest loss the attacker can force, and which $\delta$ achieves it?} A solution $\delta^\star$ is an adversarial perturbation, and $x+\delta^\star$ is an adversarial example. The notation $\max_{\delta\in\Delta}$ means the maximum is taken over the variable $\delta$ ranging in $\Delta$; the result is a number (the worst-case loss), and the maximizer is the worst-case perturbation.

\subsection{Threat models: how big a change is ``allowed''?}
\label{sub:threat}

The set $\Delta$ is called the \emph{threat model}: it formalizes the attacker's power. The most common choice is a \emph{norm ball}: all perturbations whose size, measured by some norm $\|\cdot\|$, is at most a budget $\epsilon$. We fix conventions once: $\delta$ is the perturbation, $\epsilon$ is the radius, and
\[
  \Delta \;=\; \{\,\delta : \|\delta\| \le \epsilon\,\}.
\]
Two norms dominate the literature; both are worth defining carefully because students confuse them.

\begin{definition}[$\ell_\infty$ norm]
For a vector $z=(z_1,\dots,z_n)\in\R^n$,
\[
  \|z\|_\infty \;=\; \max_{1\le i\le n} |z_i| .
\]
It is the largest single coordinate in absolute value. The constraint $\|\delta\|_\infty\le\epsilon$ says \emph{every} pixel may move by at most $\epsilon$, independently. Example: if $\epsilon=0.1$ and an image has $784$ pixels, an $\ell_\infty$ attack may shift each of the $784$ pixels by up to $0.1$ at once.
\end{definition}

\begin{definition}[$\ell_2$ norm]
For $z\in\R^n$,
\[
  \|z\|_2 \;=\; \sqrt{\,z_1^2 + z_2^2 + \dots + z_n^2\,}\,,
\]
the ordinary Euclidean length. The constraint $\|\delta\|_2\le\epsilon$ caps the \emph{total} energy of the perturbation, so the attacker can spend a big change on a few pixels if it spends little elsewhere. Example: $\delta=(0.3,0,\dots,0)$ and $\delta'=(0.1,0.1,0.1,0,\dots)$ have very different ``shapes'' but the first has $\|\delta\|_2=0.3$ while the second has $\|\delta'\|_2=\sqrt{0.03}\approx0.173$.
\end{definition}

\paragraph{Why $\ell_\infty$ is the workhorse.} \citet{kolter2018adversarial} explain its appeal: for small $\epsilon$ an $\ell_\infty$ perturbation ``add[s] such a small component to each pixel\dots that they are visually indistinguishable from the original image.'' It is a clean, if crude, stand-in for ``looks the same to a human.'' The two norms are also numerically different in scale: the $\ell_2$ budget that matches an $\ell_\infty$ budget of $\epsilon=0.1$ on a $784$-pixel image is about $\tfrac{\sqrt{2\cdot784}}{\sqrt{\pi e}}\cdot 0.1 \approx 1.35$, i.e.\ much larger numerically, because in high dimensions a per-coordinate cap of $0.1$ over $784$ coordinates buys a lot of Euclidean length \citep{kolter2018adversarial}. This is the first hint that \emph{high dimension is the attacker's friend}, a theme we return to with FGSM in Section~\ref{sec:attacks}.

\subsection{Ordinary risk versus adversarial risk}
\label{sub:risk}

So far we attacked one input. To talk about a model's overall robustness we average over the data distribution $\mathcal D$ (the unknown distribution from which real inputs are drawn).

\begin{definition}[Risk and adversarial risk]
The ordinary \emph{risk} of a model is its expected loss,
\[
  R(h_\theta) \;=\; \E_{(x,y)\sim\mathcal D}\big[\ell(h_\theta(x),y)\big].
\]
The \emph{adversarial} (or \emph{robust}) \emph{risk} replaces each point's loss by the worst-case loss in its threat set:
\[
  R_{\mathrm{adv}}(h_\theta) \;=\; \E_{(x,y)\sim\mathcal D}\Big[\;\max_{\delta\in\Delta(x)} \ell\!\big(h_\theta(x+\delta),y\big)\Big].
\]
Here $\Delta(x)$ is allowed to depend on $x$ (e.g.\ so that $x+\delta$ stays a valid image in $[0,1]^n$). In practice $\mathcal D$ is unknown, so we estimate both by averaging over a finite held-out test set.
\end{definition}

Ordinary risk asks ``how often is the model right on typical inputs?''; adversarial risk asks ``how often is the model right \emph{even when each input is nudged adversarially within budget $\epsilon$}?'' A model can have tiny ordinary risk and huge adversarial risk, exactly the pig situation.

\subsection{Two ways to report robustness: accuracy at a radius vs.\ robust radius}
\label{sub:robacc-vs-radius}

There are two natural numbers to report, and keeping them apart is essential.

\begin{definition}[Robust accuracy at radius $\epsilon$]
Fix a radius $\epsilon$. A test point $x$ (true label $y$) is \emph{robustly correct at radius $\epsilon$} if the model classifies $x+\delta$ correctly for \emph{every} $\delta$ with $\|\delta\|\le\epsilon$. The \emph{robust accuracy at radius $\epsilon$} is the fraction of test points that are robustly correct:
\[
  \mathrm{RobAcc}(\epsilon)
  \;=\; \frac1N \sum_{x} \one\big[\text{$x$ is classified correctly for all $\|\delta\|\le\epsilon$}\big],
\]
where $N$ is the number of test points and $\one[\cdot]$ is $1$ when the bracket is true and $0$ otherwise. ``Robust error at radius $\epsilon$'' is one minus this. Read it as: \emph{the fraction of points that survive every attack of size up to $\epsilon$.}
\end{definition}

\begin{definition}[Robust radius]
The \emph{robust radius} of a point $x$ is the size of the \emph{smallest} perturbation that changes the model's prediction,
\[
  r(x) \;=\; \min\{\,\|\delta\| : h_\theta(x+\delta)\ \text{is misclassified}\,\}.
\]
It is the distance from $x$ to the nearest decision boundary, measured in the chosen norm. A large $r(x)$ means $x$ sits comfortably inside its class; $r(x)=0$ would mean $x$ is already on the boundary.
\end{definition}

These two are duals of each other. A point $x$ survives radius $\epsilon$ exactly when its robust radius exceeds $\epsilon$, so
\begin{equation}
  \label{eq:duality}
  \mathrm{RobAcc}(\epsilon) \;=\; \frac1N\sum_x \one\big[\,r(x) > \epsilon\,\big],
\end{equation}
i.e.\ \emph{the fraction of points whose robust radius is bigger than $\epsilon$}. Fixing $\epsilon$ and reading off the fraction is a single vertical slice through the function $\epsilon\mapsto\mathrm{RobAcc}(\epsilon)$; measuring $r(x)$ for every point recovers the whole curve. We will see in Section~\ref{sec:attacks} that the two families of attacks correspond exactly to these two views.

\subsection{Certificate (lower bound) versus attack (upper bound)}
\label{sub:cert-vs-attack}

We can never enumerate all $\delta$, so for any single point we work with bounds on $r(x)$, and there are two opposite kinds.

\begin{definition}[Attack and certificate]
An \emph{attack} exhibits a concrete adversarial perturbation $\delta$ with $\|\delta\|\le\epsilon$ that flips the label. It \emph{proves the model is not robust at radius $\|\delta\|$}, hence gives an \emph{upper bound} on the robust radius: $r(x)\le\|\delta\|$. A \emph{certificate} is a mathematical \emph{guarantee} that \emph{no} perturbation of size $\le\rho$ can flip the label; it gives a \emph{lower bound} $r(x)\ge\rho$, i.e.\ the model is provably safe out to radius $\rho$.
\end{definition}

The two bound $r(x)$ from opposite sides:
\[
  \underbrace{\rho}_{\text{certificate, lower bound}} \;\le\; r(x) \;\le\; \underbrace{\|\delta\|}_{\text{attack, upper bound}}.
\]
An attack can only ever say ``robustness is at most this''; a certificate can only ever say ``robustness is at least this.'' Beating one does not give the other. Sections~\ref{sec:attacks} and~\ref{sec:lipschitz} are precisely the two sides: attacks (upper bounds) and Lipschitz certificates (lower bounds).

\subsection{Worked example: the linear MNIST 0-vs-1 robustness story}
\label{sub:mnist01}

We now make the accuracy--robustness gap concrete and hand-checkable, following the linear-model chapter of \citet{kolter2018adversarial}. Consider a \emph{binary} linear classifier on the MNIST handwritten digits, restricted to the two classes ``0'' and ``1.'' Each image is a vector $x\in\R^{784}$ (a $28\times28$ grid flattened). The model is
\[
  h_{w,b}(x) \;=\; w^\top x + b, \qquad w\in\R^{784},\ b\in\R,
\]
predicting class $+1$ when $w^\top x + b > 0$ and class $-1$ otherwise (we use labels $y\in\{+1,-1\}$ here). The binary cross-entropy (logistic) loss is $\ell = L\big(y\,(w^\top x+b)\big)$ with $L(z)=\log(1+e^{-z})$, which is decreasing in $z$.

\begin{example}[The 0-vs-1 numbers]
\citet{kolter2018adversarial} train this model and report:
\begin{itemize}[nosep]
  \item \textbf{Clean.} Test error about $0.04\%$, the model gets all but one test digit right. Easy problem.
  \item \textbf{Attacked at $\epsilon=0.2$ ($\ell_\infty$).} Allowing each pixel to move by up to $0.2$, the test error jumps to $\mathbf{84.5\%}$. The same model that was essentially perfect is now wrong on five out of six images, under perturbations that do not fool a human.
  \item \textbf{Robustly trained, then attacked at $\epsilon=0.2$.} Retraining the model to optimize the \emph{adversarial} risk (Section~\ref{sub:risk}) drops the worst-case error to $\mathbf{2.5\%}$ robust error, at the cost of clean error rising from $0.04\%$ to about $0.3\%$.
\end{itemize}
Two lessons. (1) Standard training, even of a simple linear model on a trivial task, is wildly non-robust. (2) Robustness is buyable but not free: the $0.04\%\to0.3\%$ rise in clean error is a small instance of the \emph{accuracy--robustness trade-off} that recurs throughout the field \citep{kolter2018adversarial}.
\end{example}

Why is the linear case special and worth memorizing? Because for it the inner maximization~\eqref{eq:inner-max} can be solved \emph{exactly} in closed form, which the next subsection's exercise asks you to do. (For deep networks it cannot, which is why we need iterative attacks in Section~\ref{sec:attacks}.)

\subsection{Exercise: the worst-case $\ell_\infty$ perturbation of a linear model}
\label{sub:exercise-linear}

\paragraph{Exercise.} Let $h_{w,b}(x)=w^\top x + b$ be a two-class linear model that currently classifies a point $x$ correctly as class $+1$, so its \emph{score} $s(x)=w^\top x+b$ is positive. An attacker may add any $\delta$ with $\|\delta\|_\infty\le\epsilon$ and wants to \emph{lower} the score (toward the wrong side of the boundary). Show that the worst-case perturbation is
\[
  \delta^\star \;=\; -\,\epsilon\,\sign(w),
\]
where $\sign(w)$ is the vector of signs of the entries of $w$ ($+1,-1,0$ entrywise), and that the resulting score drop is exactly
\[
  s(x) - s(x+\delta^\star) \;=\; \epsilon\,\|w\|_1, \qquad \|w\|_1 = \sum_{i=1}^n |w_i|.
\]

\begin{solution}
The new score is $s(x+\delta) = w^\top(x+\delta)+b = (w^\top x + b) + w^\top\delta = s(x) + w^\top\delta$, so lowering the score means making $w^\top\delta$ as \emph{negative} as possible, subject to $\|\delta\|_\infty\le\epsilon$. Write the dot product coordinatewise:
\[
  w^\top\delta \;=\; \sum_{i=1}^n w_i\,\delta_i .
\]
Each term $w_i\delta_i$ is minimized independently because the constraint $|\delta_i|\le\epsilon$ couples nothing across coordinates. For a fixed $w_i$, the most negative value of $w_i\delta_i$ over $|\delta_i|\le\epsilon$ is $-\epsilon|w_i|$, attained by choosing $\delta_i=-\epsilon$ when $w_i>0$ and $\delta_i=+\epsilon$ when $w_i<0$, i.e.\ $\delta_i = -\epsilon\,\sign(w_i)$. Summing the per-coordinate minima,
\[
  \min_{\|\delta\|_\infty\le\epsilon} w^\top\delta \;=\; \sum_{i=1}^n \big(-\epsilon|w_i|\big) \;=\; -\epsilon\sum_{i=1}^n|w_i| \;=\; -\epsilon\,\|w\|_1 ,
\]
achieved at $\delta^\star=-\epsilon\,\sign(w)$. Therefore the score drops by exactly $\epsilon\,\|w\|_1$, as claimed.
\end{solution}

Three things to take away. (i) The optimal $\delta^\star$ does not depend on $x$ at all: the same perturbation attacks every input, which is why the MNIST attack was so cheap. (ii) The drop $\epsilon\|w\|_1$ grows with $\|w\|_1$, foreshadowing why \emph{controlling the norm of the weights} (Section~\ref{sec:lipschitz} and the margin theory in later parts) controls robustness. (iii) The quantity $\|w\|_1$ is the \emph{dual norm} of the $\ell_\infty$ ball: maximizing a linear function $w^\top\delta$ over $\|\delta\|_\infty\le\epsilon$ gives $\epsilon\|w\|_1$. The general fact is $\max_{\|\delta\|_p\le\epsilon} w^\top\delta = \epsilon\|w\|_q$ with $\tfrac1p+\tfrac1q=1$ \citep{kolter2018adversarial}; the $\ell_\infty$/$\ell_1$ pair is the case $p=\infty,q=1$, and we meet it again as FGSM's optimality and as Hein's certificate.

\subsection{Robust optimization and Danskin's theorem}
\label{sub:danskin}

Putting attacker and defender together gives the \emph{saddle-point} (min--max, or \emph{robust optimization}) problem of \citet{madry2018towards}, the single organizing equation of robust training:
\begin{equation}
  \label{eq:saddle}
  \min_{\theta}\; \E_{(x,y)\sim\mathcal D}\Big[\;\max_{\delta\in\Delta}\, \ell\!\big(h_\theta(x+\delta),y\big)\Big].
\end{equation}
The \emph{inner} maximization is the attacker (find the worst $\delta$); the \emph{outer} minimization is the defender (choose $\theta$ to make even the worst-case loss small). \citet{kolter2018adversarial} state the unifying slogan plainly: \emph{every adversarial attack and defense is a method for approximately solving the inner maximization and/or the outer minimization}; the right way to describe an attack is by its norm ball plus its optimization method, not by a brand name.

To train by gradient descent on~\eqref{eq:saddle} we must differentiate, with respect to $\theta$, a quantity that itself contains a maximization over $\delta$. The justification is Danskin's theorem; we state it and, as instructed, do not prove it.

\begin{theorem}[Danskin, stated without proof]
\label{thm:danskin}
Suppose the inner maximization in~\eqref{eq:saddle} is solved \emph{exactly}, with maximizer $\delta^\star=\delta^\star(\theta)$. Then the gradient with respect to $\theta$ of the value $\max_{\delta\in\Delta}\ell(h_\theta(x+\delta),y)$ equals the gradient of $\ell(h_\theta(x+\delta),y)$ evaluated at the fixed point $\delta=\delta^\star$:
\[
  \nabla_\theta\!\left[\max_{\delta\in\Delta}\ell\big(h_\theta(x+\delta),y\big)\right]
  \;=\; \nabla_\theta\,\ell\big(h_\theta(x+\delta^\star),y\big)\Big|_{\delta^\star\ \text{held fixed}} .
\]
\end{theorem}

In words: to take the gradient of the worst-case loss, first find the worst $\delta^\star$, then differentiate as if $\delta^\star$ were a constant. This is what licenses ``adversarial training'': repeatedly attack each batch, then take a normal gradient step on the attacked inputs. \citet{kolter2018adversarial} stress the catch: Danskin requires the inner max to be solved \emph{exactly}, which for deep networks it never is. In practice one solves it ``well enough'' (a strong attack such as PGD, Section~\ref{sec:attacks}) and accepts that the theorem holds only approximately; if the inner attack is weak, a stronger attack can later defeat the supposedly robust model.

% =====================================================================
\section{Lipschitz constants and robustness certificates}
\label{sec:lipschitz}

Section~\ref{sec:adv-from-scratch} promised certificates: \emph{provable lower bounds} on the robust radius, the guarantee that no attack within radius $\rho$ can change the label. This section delivers the most important one. The single idea is: \emph{if the model's output cannot change quickly as the input moves, and the input currently sits well inside its class, then a small input change cannot flip the decision.} ``Cannot change quickly'' is made precise by the Lipschitz constant; ``well inside its class'' by the margin. We build both from scratch; this material confused even the professor in the last draft, so we go slowly and keep the jargon minimal.

\subsection{Lipschitz continuity and the Lipschitz constant}
\label{sub:lip-def}

\begin{definition}[Lipschitz continuity]
A function $f:\R^n\to\R^m$ is \emph{Lipschitz continuous} if there is a constant $L\ge0$ such that
\[
  \|f(x)-f(y)\| \;\le\; L\,\|x-y\| \qquad \text{for all } x,y .
\]
The smallest such $L$ is the \emph{Lipschitz constant} of $f$, written $\Lip(f)$ \citep{scaman2018lipschitz}. (Until Section~\ref{sub:margin-cert} take both norms to be $\ell_2$.)
\end{definition}

The plain-English reading from the Wikipedia article \citep{wiki_lipschitz} is the cleanest: \emph{the absolute value of the slope of the line joining any two points on the graph is at most $L$.} The output can change by at most $L$ times the input change: $L$ is a speed limit on the function. Three examples fix the idea.

\begin{example}[Three one-dimensional functions]\leavevmode
\begin{itemize}[nosep]
  \item $f(x)=|x|$ is $1$-Lipschitz: by the reverse triangle inequality $\big||x|-|y|\big|\le|x-y|$, so $L=1$ works, and taking $y=0$ shows no smaller $L$ works. Note $f$ is not differentiable at $0$; Lipschitz does not require differentiability.
  \item $f(x)=cx$ (a line of slope $c$) has $\Lip(f)=|c|$: $|cx-cy|=|c|\,|x-y|$ with equality always.
  \item $f(x)=x^2$ is \emph{not} Lipschitz on all of $\R$: $|x^2-y^2|=|x+y|\,|x-y|$, and the factor $|x+y|$ is unbounded, so no single finite $L$ caps the slope everywhere \citep{wiki_lipschitz}. (It \emph{is} Lipschitz on any bounded interval; this is the difference between a global and a local constant, which matters below.)
\end{itemize}
\end{example}

For differentiable functions the constant is the supremum of the gradient norm. This is Rademacher's theorem applied to Lipschitz functions; we use it as a computational rule.

\begin{theorem}[Lipschitz constant via the gradient \citep{scaman2018lipschitz}]
\label{thm:lip-grad}
If $f:\R^n\to\R$ is (locally) Lipschitz, then it is differentiable almost everywhere and
\[
  \Lip(f) \;=\; \sup_{x} \|\nabla f(x)\| .
\]
For vector-valued $f:\R^n\to\R^m$ the same holds with $\|\nabla f(x)\|$ replaced by the operator norm $\|D_x f\|$ of the Jacobian (defined next).
\end{theorem}

This already connects to robustness: a small Lipschitz constant means the output cannot move far when the input is nudged, which is exactly the kind of stability an adversarial example tries to violate \citep{tsuzuku2018lipschitz}.

\subsection{The spectral norm: the Lipschitz constant of a linear map}
\label{sub:spectral}

The most important Lipschitz constant is for the simplest function, a linear map $f(x)=Wx$.

\begin{definition}[Operator / spectral norm]
For a matrix $W\in\R^{m\times n}$, the \emph{operator norm} (also \emph{spectral norm}) is the largest factor by which $W$ can stretch a unit vector:
\[
  \|W\|_2 \;=\; \sup_{\|z\|_2 = 1} \|Wz\|_2 .
\]
\end{definition}

\begin{theorem}[Spectral norm $=$ Lipschitz constant of $x\mapsto Wx$]
\label{thm:spectral-lip}
For $f(x)=Wx$, $\ \Lip(f)=\|W\|_2=\sigma_{\max}(W)=\sqrt{\lambda_{\max}(W^\top W)}$, where $\sigma_{\max}$ is the largest singular value of $W$ and $\lambda_{\max}$ the largest eigenvalue.
\end{theorem}

\begin{proof}
We compute $\sup_{\|z\|=1}\|Wz\|$ and identify it with the largest singular value via the Rayleigh quotient. Each step on its own line.
\begin{enumerate}[label=(\arabic*),nosep]
  \item For any $x,y$, linearity gives $f(x)-f(y)=Wx-Wy=W(x-y)$.
  \item Hence $\|f(x)-f(y)\| = \|W(x-y)\|$, and setting $z=x-y$ the Lipschitz constant is $\Lip(f)=\sup_{z\neq0}\|Wz\|/\|z\| = \sup_{\|z\|=1}\|Wz\|=\|W\|_2$. So the Lipschitz constant \emph{is} the operator norm; it remains to evaluate it.
  \item Squaring is monotone on nonnegatives, so $\sup_{\|z\|=1}\|Wz\|$ is the square root of $\sup_{\|z\|=1}\|Wz\|^2$.
  \item Expand the squared norm using the definition of the inner product: $\|Wz\|^2 = (Wz)^\top(Wz) = z^\top W^\top W z$.
  \item The matrix $A:=W^\top W$ is symmetric ($A^\top=(W^\top W)^\top = W^\top W = A$) and positive semidefinite ($z^\top A z = \|Wz\|^2\ge0$). By the spectral theorem it has an orthonormal eigenbasis $v_1,\dots,v_n$ with real eigenvalues $\lambda_1\ge\dots\ge\lambda_n\ge0$.
  \item Write a unit vector $z=\sum_i c_i v_i$ with $\sum_i c_i^2=1$. Then $z^\top A z = \sum_i \lambda_i c_i^2$, a weighted average of the $\lambda_i$ with weights $c_i^2$ summing to $1$.
  \item A weighted average is maximized by putting all weight on the largest value: $\sum_i \lambda_i c_i^2 \le \lambda_1$, with equality when $c_1=1$ (i.e.\ $z=v_1$). Hence $\sup_{\|z\|=1} z^\top A z = \lambda_1 = \lambda_{\max}(W^\top W)$.
  \item Taking the square root from step (3): $\|W\|_2 = \sqrt{\lambda_{\max}(W^\top W)} = \sigma_{\max}(W)$, since singular values of $W$ are square roots of eigenvalues of $W^\top W$ by definition.
\end{enumerate}
\end{proof}

\begin{example}[Spectral norms by hand]\leavevmode
\begin{itemize}[nosep]
  \item $W=\begin{psmallmatrix}3&0\\0&1\end{psmallmatrix}$. Then $W^\top W=\begin{psmallmatrix}9&0\\0&1\end{psmallmatrix}$, eigenvalues $9$ and $1$, so $\|W\|_2=\sqrt9=3$. (For a diagonal matrix the spectral norm is just the largest $|$entry$|$.)
  \item $W=\begin{psmallmatrix}1&1\\0&1\end{psmallmatrix}$. Then $W^\top W=\begin{psmallmatrix}1&1\\1&2\end{psmallmatrix}$, with eigenvalues solving $\lambda^2-3\lambda+1=0$, i.e.\ $\lambda=\tfrac{3\pm\sqrt5}{2}$. The spectral norm is $\sqrt{\tfrac{3+\sqrt5}{2}}\approx1.618$ (the golden ratio). Notice the entries are all $0$ or $1$ yet the stretch exceeds $1$: off-diagonal coupling matters.
\end{itemize}
\end{example}

\subsection{Activations are 1-Lipschitz, and constants multiply across layers}
\label{sub:composition}

A neural network alternates linear maps with simple nonlinear functions called \emph{activations}, applied entrywise. The most common is the ReLU (rectified linear unit) $\mathrm{ReLU}(z)=\max\{0,z\}$. We need two facts.

\begin{lemma}[ReLU is $1$-Lipschitz]
\label{lem:relu}
$|\mathrm{ReLU}(a)-\mathrm{ReLU}(b)|\le|a-b|$ for all real $a,b$, so $\Lip(\mathrm{ReLU})=1$. Applied entrywise to a vector, the map is still $1$-Lipschitz in $\ell_2$.
\end{lemma}

\begin{proof}
Check the four sign cases for $a,b$ (each line is one case).
\begin{itemize}[nosep]
  \item $a\ge0,\ b\ge0$: $|\mathrm{ReLU}(a)-\mathrm{ReLU}(b)|=|a-b|$.
  \item $a\le0,\ b\le0$: both map to $0$, difference $0\le|a-b|$.
  \item $a\ge0,\ b\le0$: $|\mathrm{ReLU}(a)-\mathrm{ReLU}(b)|=|a-0|=a\le a-b=|a-b|$ (since $b\le0$).
  \item $a\le0,\ b\ge0$: symmetric to the previous case, gives $b\le|a-b|$.
\end{itemize}
In every case the inequality holds, and equality occurs (e.g.\ both positive), so the constant $1$ is tight. Applied entrywise, $\|\mathrm{ReLU}(u)-\mathrm{ReLU}(v)\|_2^2=\sum_i(\mathrm{ReLU}(u_i)-\mathrm{ReLU}(v_i))^2\le\sum_i(u_i-v_i)^2=\|u-v\|_2^2$.
\end{proof}

Other activations have known constants too: the sigmoid is $\tfrac14$-Lipschitz and $\tanh$ is $1$-Lipschitz \citep{tsuzuku2018lipschitz}. Now the composition rule, which lets us bound a whole network.

\begin{theorem}[Composition / product over layers \citep{szegedy2014intriguing,tsuzuku2018lipschitz}]
\label{thm:composition}
If $g$ has Lipschitz constant $L_g$ and $f$ has Lipschitz constant $L_f$, then their composition $f\circ g$ (apply $g$, then $f$) satisfies $\Lip(f\circ g)\le L_f\,L_g$. Consequently, for a $d$-layer network $F=f_d\circ\cdots\circ f_1$,
\[
  \Lip(F) \;\le\; \prod_{k=1}^{d} \Lip(f_k).
\]
\end{theorem}

\begin{proof}
Apply the two Lipschitz bounds in turn, inner first. For any $x_1,x_2$:
\begin{align*}
  \|f(g(x_1))-f(g(x_2))\|
   &\le L_f\,\|g(x_1)-g(x_2)\| &&\text{(}f\text{ is }L_f\text{-Lipschitz, on inputs }g(x_1),g(x_2)\text{)}\\
   &\le L_f\,\big(L_g\,\|x_1-x_2\|\big) &&\text{(}g\text{ is }L_g\text{-Lipschitz)}\\
   &= (L_f L_g)\,\|x_1-x_2\| .
\end{align*}
So $L_fL_g$ is a valid Lipschitz constant for $f\circ g$, hence $\Lip(f\circ g)\le L_fL_g$. Iterating over the $d$ layers gives the product. The bound is an \emph{inequality}, not equality, because the worst-stretch direction of one layer need not align with that of the next; this looseness is the subject of Section~\ref{sub:honest}.
\end{proof}

For a standard network, each linear layer $f_k(z)=W_kz+b_k$ has $\Lip(f_k)=\|W_k\|_2$ (the bias $b_k$ shifts the output but does not change the constant, since it cancels in $f_k(x_1)-f_k(x_2)$), and each ReLU layer has constant $1$. So
\[
  \Lip(F) \;\le\; \prod_{k} \|W_k\|_2 ,
\]
the product of the spectral norms of the weight matrices, the original bound of \citet{szegedy2014intriguing}.

\subsection{The prediction margin}
\label{sub:margin}

The Lipschitz constant says how fast the logits can move. To turn that into a guarantee we also need to know how much room the current prediction has before the top logit loses to a runner-up.

\begin{definition}[Prediction margin \citep{tsuzuku2018lipschitz}]
Let $F(x)\in\R^k$ be the logit vector and let $t$ be the predicted (top) class at $x$. The \emph{prediction margin} is
\[
  M_{F,x} \;=\; F(x)_t \;-\; \max_{i\neq t} F(x)_i ,
\]
the gap between the winning logit and the best competitor. It is positive exactly when $x$ is classified as $t$, and a larger margin means the decision is more decisive.
\end{definition}

\subsection{The Lipschitz--margin certificate}
\label{sub:margin-cert}

Here is the payoff: a provable lower bound on the robust radius, due to \citet{tsuzuku2018lipschitz}. It is the certificate the paper builds on.

\begin{proposition}[Lipschitz--margin certificate, Tsuzuku Prop.~1 \citep{tsuzuku2018lipschitz}]
\label{prop:tsuzuku}
Let $F:\R^n\to\R^k$ have $\ell_2$-Lipschitz constant $L$ (so $\|F(x)-F(x+\delta)\|_2\le L\|\delta\|_2$ for any perturbation $\delta$) and let $M_{F,x}$ be the prediction margin at $x$. If
\[
  M_{F,x} \;\ge\; \sqrt2\,L\,\|\delta\|_2 ,
\]
then the predicted class does not change at $x+\delta$. Equivalently, the certified $\ell_2$ robust radius is
\[
  r(x) \;\ge\; \frac{M_{F,x}}{\sqrt2\,L} .
\]
\end{proposition}

We prove it in full, slowly. The one nonobvious step (where the $\sqrt2$ comes from) is a Cauchy--Schwarz inequality on two coordinates, so we first record a small lemma.

\begin{lemma}[The ``max of the rest'' is $1$-Lipschitz \citep{tsuzuku2018lipschitz}]
\label{lem:maxrest}
For real vectors $a$ and $b$, $\ \big|\max_{i\neq t} a_i - \max_{i\neq t} b_i\big| \le \max_{i\neq t}|a_i-b_i|.$
\end{lemma}

\begin{proof}
Without loss of generality assume $\max_{i\neq t} a_i \ge \max_{i\neq t} b_i$ (otherwise swap $a,b$). Let $j=\arg\max_{i\neq t} a_i$. Then
\begin{align*}
  \Big|\max_{i\neq t} a_i - \max_{i\neq t} b_i\Big|
   &= \max_{i\neq t} a_i - \max_{i\neq t} b_i &&\text{(the quantity is nonnegative by assumption)}\\
   &= a_j - \max_{i\neq t} b_i &&\text{(}j\text{ attains the max for }a\text{)}\\
   &\le a_j - b_j &&\text{(}\max_{i\neq t} b_i \ge b_j\text{, so subtracting it is }\le\text{ subtracting }b_j\text{)}\\
   &\le \max_{i\neq t}|a_i-b_i| . &&\text{(}a_j-b_j\le|a_j-b_j|\le\text{ the max)}
\end{align*}
\end{proof}

\begin{proof}[Proof of Proposition~\ref{prop:tsuzuku}]
It suffices to show the margin at the perturbed point stays nonnegative; in fact we show it drops by at most $\sqrt2\,L\|\delta\|_2$:
\begin{equation}
  \label{eq:cert-target}
  M_{F,x+\delta} \;=\; F(x+\delta)_t - \max_{i\neq t}F(x+\delta)_i
  \;\ge\; M_{F,x} - \sqrt2\,L\,\|\delta\|_2 .
\end{equation}
Granting~\eqref{eq:cert-target}, if $M_{F,x}\ge\sqrt2 L\|\delta\|_2$ then $M_{F,x+\delta}\ge0$, so class $t$ still wins. Now prove~\eqref{eq:cert-target}. Add and subtract the unperturbed logits:
\begin{align*}
  M_{F,x+\delta}
   &= F(x+\delta)_t - \max_{i\neq t}F(x+\delta)_i \\
   &= \underbrace{\big(F(x)_t - \max_{i\neq t}F(x)_i\big)}_{=\,M_{F,x}}
      + \underbrace{\big(F(x+\delta)_t - F(x)_t\big)}_{=:a_1}
      - \underbrace{\big(\max_{i\neq t}F(x+\delta)_i - \max_{i\neq t}F(x)_i\big)}_{=:a_2'} .
\end{align*}
Bound the two change terms from below by their negative absolute values:
\[
  M_{F,x+\delta} \;\ge\; M_{F,x} - |a_1| - |a_2'| .
\]
For $a_1$, $\ |a_1| = |F(x+\delta)_t - F(x)_t|$ is the change in a single coordinate of the logit vector. For $a_2'$, Lemma~\ref{lem:maxrest} (with $a=F(x+\delta)$, $b=F(x)$) gives $|a_2'|\le\max_{i\neq t}|F(x+\delta)_i - F(x)_i|$, again the absolute change of a \emph{single} (other) coordinate. Call the change of the top logit $\Delta_t$ and the change of the worst competitor $\Delta_j$; both are entries of the vector $F(x+\delta)-F(x)$, so
\[
  \sqrt{\Delta_t^2 + \Delta_j^2} \;\le\; \big\|F(x+\delta)-F(x)\big\|_2 \;\le\; L\,\|\delta\|_2 ,
\]
the last step being the definition of the Lipschitz constant $L$. We must bound $|a_1|+|a_2'| = |\Delta_t|+|\Delta_j|$, the sum of two absolute values, by something times $\sqrt{\Delta_t^2+\Delta_j^2}$. This is where $\sqrt2$ enters, by Cauchy--Schwarz on the two-vector $(|\Delta_t|,|\Delta_j|)$ against $(1,1)$:
\[
  |\Delta_t| + |\Delta_j|
  \;=\; (1,1)\cdot(|\Delta_t|,|\Delta_j|)
  \;\le\; \big\|(1,1)\big\|_2\,\big\|(\Delta_t,\Delta_j)\big\|_2
  \;=\; \sqrt2\,\sqrt{\Delta_t^2+\Delta_j^2}\,.
\]
Chaining the bounds:
\[
  |a_1|+|a_2'| \;\le\; \sqrt2\,\sqrt{\Delta_t^2+\Delta_j^2} \;\le\; \sqrt2\,L\,\|\delta\|_2 .
\]
Substituting into $M_{F,x+\delta}\ge M_{F,x}-|a_1|-|a_2'|$ gives~\eqref{eq:cert-target}, completing the proof.
\end{proof}

\paragraph{Reading the certificate.} The radius $r(x)\ge M_{F,x}/(\sqrt2\,L)$ says exactly what robustness should: \emph{decisive predictions (big margin) and gentle models (small $L$) are certifiably safe over a large ball.} Both knobs appear, and the certificate is the lower-bound (defender's) counterpart to the upper-bound attacks of Section~\ref{sec:attacks}.

\subsection{Worked examples: one-layer certificates}
\label{sub:cert-examples}

\begin{example}[The exact two-class radius \citep{hein2017formal}]
For a \emph{linear} two-class model the $\sqrt2$ disappears and the certificate becomes exact, which makes a clean sanity check. Let the two class weight vectors differ by $w_c-w_j=(2,0)$, evaluated at $x=(2,5)$. The decision is governed by the scalar $g(x)=\langle w_c-w_j,\,x\rangle = 2\cdot2 + 0\cdot5 = 4$ (the margin), and $\|w_c-w_j\|_2 = 2$. Because here the margin is a single scalar linear function (not a max over many classes), the certified $\ell_2$ radius is exactly $\dfrac{4}{2}=2$. The general exact formula of \citet{hein2017formal} for a linear classifier is
\[
  r(x) \;=\; \min_{j\neq c}\frac{\langle w_c-w_j,\,x\rangle}{\|w_c-w_j\|_2}\,,
\]
the smallest distance to any pairwise boundary, itself a clean instance of the dual-norm computation from Section~\ref{sub:exercise-linear}.
\end{example}

\begin{example}[One-hidden-layer network]
Take $F(x)=W_2\,\mathrm{ReLU}(W_1 x)$ with $\|W_1\|_2=3$ and $\|W_2\|_2=2$. By Theorem~\ref{thm:composition} and Lemma~\ref{lem:relu}, the global Lipschitz constant is bounded by $L\le 3\cdot1\cdot2 = 6$. If the prediction margin at some $x$ is $M_{F,x}=12$, the Tsuzuku certificate gives
\[
  r(x) \;\ge\; \frac{M_{F,x}}{\sqrt2\,L} \;=\; \frac{12}{\sqrt2\cdot6} \;=\; \frac{12}{6\sqrt2} \;=\; \sqrt2 \;\approx\; 1.414 .
\]
So we have \emph{proved} that no $\ell_2$ perturbation of size up to $1.414$ can change this prediction: no attack, however clever, can do better at this point.
\end{example}

\subsection{Honest caveat: the bound is loose, and the exact constant is NP-hard}
\label{sub:honest}

A certificate is only useful if you are honest about its slack. Two points, both from the sources.

First, the product bound $\Lip(F)\le\prod_k\|W_k\|_2$ is an \emph{over-estimate} of the true Lipschitz constant, often badly so. The reason was already visible in the proof of Theorem~\ref{thm:composition}: it assumes every layer is stretched maximally in the same direction, which never happens. A larger $L$ than the truth makes the certified radius $M/(\sqrt2 L)$ \emph{smaller} than it could be, so the guarantee is conservative: safe, but weak. \citet{hein2017formal} report that \emph{local} Lipschitz estimates (valid in a neighborhood of $x$) can be up to about $8\times$ smaller than the global product bound, yielding correspondingly stronger certificates.

Second, you might hope to just compute the exact $\Lip(F)$ and avoid the slack. You cannot, in general.

\begin{theorem}[Exact Lipschitz computation is NP-hard \citep{scaman2018lipschitz}]
\label{thm:nphard}
Computing the exact Lipschitz constant of even a two-layer ReLU network is NP-hard. Hence every practical certificate uses an \emph{upper bound} on $L$ (such as the product of spectral norms), and the resulting robust radius is a \emph{lower} bound on the true robust radius.
\end{theorem}

This is the right mental model: \citet{szegedy2014intriguing} put it as ``small Lipschitz bounds forbid attacks, but large bounds do not imply attacks exist.'' A small certified $L$ guarantees safety; a large $L$ (or a loose bound) simply means we failed to prove safety, not that the model is necessarily fragile.

\subsection{Exercise: a full two-layer certificate}
\label{sub:cert-exercise}

\paragraph{Exercise.} Let $F(x)=W_2\,\mathrm{ReLU}(W_1 x)$ with
\[
  W_1=\begin{psmallmatrix}2&0\\0&1\end{psmallmatrix},\qquad
  W_2=\begin{psmallmatrix}1&1\\-1&1\end{psmallmatrix}.
\]
(a) Compute $\|W_1\|_2$ and $\|W_2\|_2$ (eigenvalues of $W^\top W$). (b) Give the global Lipschitz bound $L$ for $F$. (c) At $x=(1,1)$, compute $F(x)$, the predicted class, and the prediction margin $M_{F,x}$. (d) State the Tsuzuku certified radius $M_{F,x}/(\sqrt2\,L)$, and explain in one sentence why it is a conservative \emph{lower} bound on the true robust radius.

\begin{solution}
(a) $W_1$ is diagonal, so $\|W_1\|_2 = \max\{2,1\}=2$. For $W_2$, $W_2^\top W_2=\begin{psmallmatrix}2&0\\0&2\end{psmallmatrix}=2I$, so both eigenvalues are $2$ and $\|W_2\|_2=\sqrt2$. (Indeed $W_2/\sqrt2$ is a rotation; $W_2$ scales every direction by $\sqrt2$.)

(b) By Theorem~\ref{thm:composition} and ReLU being $1$-Lipschitz, $L \le \|W_1\|_2\cdot1\cdot\|W_2\|_2 = 2\cdot\sqrt2 = 2\sqrt2 \approx 2.828$.

(c) $W_1 x = (2,1)$, both positive, so $\mathrm{ReLU}(W_1x)=(2,1)$. Then
$F(x)=W_2(2,1)^\top = (2+1,\ -2+1) = (3,-1)$.
The top logit is the first ($3$), so the predicted class is $1$, and $M_{F,x}=3-(-1)=4$.

(d) Certified radius $= \dfrac{M_{F,x}}{\sqrt2\,L} = \dfrac{4}{\sqrt2\cdot 2\sqrt2} = \dfrac{4}{4} = 1.$ It is a \emph{lower} bound because $L=2\sqrt2$ is itself an over-estimate of the true Lipschitz constant (the product bound assumes both layers stretch in the same direction, Theorem~\ref{thm:composition}; and computing the exact constant is NP-hard, Theorem~\ref{thm:nphard}), so the real safe radius is at least $1$ and likely larger.
\end{solution}

% =====================================================================
\section{Attacks in detail, and how robustness is measured}
\label{sec:attacks}

We can now make Section~\ref{sec:adv-from-scratch}'s ``attack = upper bound'' concrete. This section gives the standard attacks, each with: definition, intuition, the algorithm, a hand-computable example, and a short exercise. We follow the unifying taxonomy of \citet{kolter2018adversarial,madry2018towards}.

\subsection{The organizing idea: norm ball $+$ optimizer}
\label{sub:taxonomy}

\citet{kolter2018adversarial} insist on one classification: \emph{every attack is (1) a norm ball plus (2) a method for optimizing over it.} This splits attacks into two families, which match the two robustness numbers of Section~\ref{sub:robacc-vs-radius}.
\begin{itemize}[nosep]
  \item \textbf{Fixed-$\epsilon$ (norm-bounded).} Solve $\max_{\|\delta\|\le\epsilon}\ell(h(x+\delta),y)$: ``can this point be broken within budget $\epsilon$?'' Answering it for many points gives $\mathrm{RobAcc}(\epsilon)$, one vertical slice of the robustness curve. Members: FGSM, PGD, AutoAttack.
  \item \textbf{Minimum-norm.} Solve $\min\|\delta\|$ such that $x+\delta$ is misclassified: ``how far is the boundary?'' Answering it gives the robust radius $r(x)$, the whole curve. Members: C\&W, DDN.
\end{itemize}
They are duals via Equation~\eqref{eq:duality}: $\mathrm{RobAcc}(\epsilon)=\tfrac1N\sum_x\one[r(x)>\epsilon]$. A fixed-$\epsilon$ attack reads the curve at one $\epsilon$; a minimum-norm attack measures $r(x)$ directly.

\subsection{FGSM: the Fast Gradient Sign Method}
\label{sub:fgsm}

\paragraph{Definition.} The \emph{Fast Gradient Sign Method} of \citet{goodfellow2015explaining} takes one step from $x$ in the direction that, to first order, increases the loss the most under an $\ell_\infty$ budget:
\begin{equation}
  \label{eq:fgsm}
  x' \;=\; x + \epsilon\,\sign\!\big(\nabla_x \ell(h_\theta(x),y)\big).
\end{equation}
One gradient evaluation, one step. The $\sign$ is applied entrywise.

\paragraph{Intuition: the linearity hypothesis.} Why $\sign$ of the gradient, and why $\ell_\infty$? \citet{goodfellow2015explaining} argue that adversarial examples come from models being \emph{too linear in high dimension}, not too nonlinear. Linearize the model near $x$: a linear unit responds to $x' = x+\delta$ as $w^\top x' = w^\top x + w^\top\delta$. By the exercise of Section~\ref{sub:exercise-linear}, maximizing the perturbation term $w^\top\delta$ over $\|\delta\|_\infty\le\epsilon$ gives $\delta=\epsilon\,\sign(w)$ and value $\epsilon\|w\|_1$. Now the punchline: if $w$ has $n$ entries of average magnitude $m$, then $\|w\|_1\approx mn$, so the change to the output grows \emph{linearly with the dimension $n$}, while the per-pixel change $\|\delta\|_\infty$ stays fixed at $\epsilon$ \citep{goodfellow2015explaining}. In high dimension many imperceptible per-pixel nudges, all aligned with the sign of the weights, add up to a large swing in the output. The famous demonstration: panda $+\ 0.007\cdot\sign(\nabla)$ classified as ``gibbon'' with $99.3\%$ confidence \citep{goodfellow2015explaining}.

\paragraph{Exactness for logistic regression.} For a two-class logistic model FGSM is not an approximation; it is the \emph{exact} worst case. The reason is the identity $w^\top\sign(w)=\|w\|_1$ used in Section~\ref{sub:exercise-linear}: the sign step realizes the dual-norm optimum exactly when the model is linear \citep{goodfellow2015explaining,kolter2018adversarial}.

\paragraph{Algorithm.}
\begin{enumerate}[label=\textbf{\arabic*.},nosep]
  \item Compute the gradient $g = \nabla_x \ell(h_\theta(x),y)$ by backpropagation.
  \item Form the signed step $\delta = \epsilon\,\sign(g)$.
  \item Return $x' = x+\delta$ (clip to the valid input range if needed).
\end{enumerate}

\begin{example}[FGSM on a 2D logistic model]
Let $w=(2,-1)$, $b=0$, true label $y=+1$, and a point $x$ with positive score $s(x)=w^\top x$. The gradient of the logistic loss with respect to $x$ points along $-yw=-w$ (loss decreases as the score $y\,w^\top x$ grows), so the loss-increasing sign direction is $-\sign(w)=(-1,+1)$, giving the attack $\delta^\star = -\epsilon\,\sign(w) = (-\epsilon,\,+\epsilon)$. The score change is
\[
  w^\top\delta^\star = (2)(-\epsilon) + (-1)(+\epsilon) = -3\epsilon = -\epsilon\,\|w\|_1, \qquad \|w\|_1=|2|+|-1|=3.
\]
The margin drops by $\epsilon\|w\|_1=3\epsilon$, matching Section~\ref{sub:exercise-linear} exactly. Generalize to $w\in\R^n$ with average $|w_i|=m$: the drop is $\epsilon\|w\|_1\approx \epsilon m n$, which grows with $n$, the high-dimension punchline, derived by hand.
\end{example}

\paragraph{Exercise.} For $w=(1,1,1,1)$, $b=0$, $\epsilon=0.1$, true label $+1$: write the FGSM perturbation $\delta^\star$, compute the margin drop $\epsilon\|w\|_1$, and compare to the same computation for $w=(1,1,\dots,1)\in\R^{100}$. \emph{(Answer: $\delta^\star=(-0.1,-0.1,-0.1,-0.1)$, drop $0.4$; in $\R^{100}$ the drop is $0.1\cdot100=10$, with the same per-pixel size $0.1$, illustrating linear growth in dimension.)}

\subsection{PGD: Projected Gradient Descent}
\label{sub:pgd}

\paragraph{Definition.} FGSM takes one big step and trusts the linear approximation over the whole ball. \emph{Projected Gradient Descent} \citep{madry2018towards} takes many small steps, projecting back into the $\epsilon$-ball after each, and starting from a random point inside the ball. It is the empirical ``ultimate first-order adversary.''

\paragraph{The saddle point.} PGD is the inner solver for the robust-optimization problem~\eqref{eq:saddle},
\[
  \min_\theta\ \E\Big[\max_{\delta\in\Delta}\ \ell(\theta,x+\delta,y)\Big],
  \qquad \Delta=\{\delta:\|\delta\|_\infty\le\epsilon\}:
\]
the inner $\max$ is the attack, the outer $\min$ the defense \citep{madry2018towards}.

\paragraph{Iteration.} Starting from a \emph{uniformly random} $x^0$ in the $\epsilon$-ball around $x$, repeat
\begin{equation}
  \label{eq:pgd}
  x^{t+1} \;=\; \mathrm{Proj}_{x+\Delta}\Big(\,x^t + \alpha\,\sign\!\big(\nabla_x \ell(h_\theta(x^t),y)\big)\Big),
\end{equation}
where $\alpha>0$ is the step size and $\mathrm{Proj}_{x+\Delta}$ projects back into the threat set. For the $\ell_\infty$ ball the projection is just \emph{clipping}: any coordinate of $\delta=x^{t+1}-x$ that exceeds $\epsilon$ in absolute value is set to $\pm\epsilon$ \citep{kolter2018adversarial}.

\paragraph{Step-then-project, and why a random start.} Each iteration first takes a signed-gradient step (like a mini-FGSM), which may leave the box, then clips back to the nearest face or corner of the box, the ``step then project'' picture. FGSM is exactly one such step with no random start and $\alpha=\epsilon$. The random start matters because PGD can get stuck at local maxima of the (non-concave) inner loss; restarting from different random points inside the ball and keeping the worst found is a cheap way to escape some of them \citep{kolter2018adversarial}. \citet{madry2018towards} observe empirically that the loss values at many random-restart local maxima are \emph{well concentrated}, which is the evidence behind calling PGD the ``ultimate first-order adversary.''

\paragraph{Normalized steepest descent and the step-size rule.} Why $\sign(\nabla)$ rather than $\nabla$ itself? At $\delta=0$ the raw gradient is tiny (about $10^{-6}$ per pixel in the tutorial's CNN), so plain gradient ascent would need an absurd step size ($\sim10^4$) and then overshoot \citep{kolter2018adversarial}. The fix is \emph{normalized steepest descent} under the chosen norm: choose the update direction $v$ maximizing $v^\top\nabla$ subject to a norm cap on $v$. For $\ell_\infty$ this gives $v=\sign(\nabla)$ (the sign step); for $\ell_2$ it gives $v=\nabla/\|\nabla\|_2$ (the normalized gradient) \citep{kolter2018adversarial}. Practical rule of thumb \citep{kolter2018adversarial}: take $\alpha$ a small fraction of $\epsilon$, and the number of iterations a small multiple of $\epsilon/\alpha$.

\paragraph{Algorithm.}
\begin{enumerate}[label=\textbf{\arabic*.},nosep]
  \item Initialize $\delta$ uniformly at random in $\{\delta:\|\delta\|_\infty\le\epsilon\}$.
  \item For $t=1,\dots,T$: compute $g=\nabla_x\ell(h_\theta(x+\delta),y)$; update $\delta \leftarrow \delta + \alpha\,\sign(g)$; clip each coordinate of $\delta$ to $[-\epsilon,\epsilon]$.
  \item (Optional) repeat from random restarts; keep the $\delta$ of highest loss.
\end{enumerate}

\begin{example}[Step-then-project by hand]
Let the model be the 2D linear $s(x)=w^\top x$, $w=(2,-1)$, $y=+1$, $\epsilon=1$, $\alpha=2$, starting from $\delta^0=(0,0)$. The loss-increasing direction is $-\sign(w)=(-1,+1)$. \textbf{Step:} $\delta = (0,0) + 2\cdot(-1,+1) = (-2,+2)$. \textbf{Project} onto $\|\delta\|_\infty\le1$: clip each coordinate to $[-1,1]$, giving $\delta=(-1,+1)$, a corner of the box, which is exactly the FGSM solution $-\epsilon\,\sign(w)$. So for a linear model PGD reaches the FGSM corner in one step-and-clip; the value of many small steps shows up only for genuinely nonlinear models, where the gradient direction changes as you move.
\end{example}

\paragraph{Exercise.} Explain, in two sentences, where FGSM and PGD each sit on the accuracy-vs-radius curve of Section~\ref{sub:robacc-vs-radius}, and why PGD at a fixed $\epsilon$ never reports \emph{higher} robust accuracy than the true value. \emph{(Answer: both are fixed-$\epsilon$ attacks, so each estimates a single point $\mathrm{RobAcc}(\epsilon)$; since any successful $\delta$ they find with $\|\delta\|\le\epsilon$ proves $r(x)\le\epsilon$, an attack can only \emph{lower} the reported robust accuracy toward the truth, i.e.\ it is an upper bound on robustness (Section~\ref{sub:cert-vs-attack}).)}

\subsection{C\&W: a minimum-norm attack}
\label{sub:cw}

\paragraph{Definition and intuition.} The Carlini--Wagner attack \citep{carlini2017towards} switches families: instead of fixing $\epsilon$, it asks for the \emph{smallest} perturbation that misclassifies, decoupling ``be small'' from ``be misclassified.'' It minimizes
\begin{equation}
  \label{eq:cw}
  \min_{\delta}\; \|\delta\|_p \;+\; c\cdot f(x+\delta)
  \qquad\text{subject to } x+\delta\in[0,1]^n,
\end{equation}
where $f$ is a surrogate that is $\le0$ exactly when the example is misclassified, and the trade-off constant $c>0$ is found by \emph{binary search}: small $c$ favors a tiny but possibly non-adversarial $\delta$; large $c$ forces misclassification at the cost of a larger $\delta$. \citet{carlini2017towards} pick the smallest $c$ for which the attack still succeeds.

\paragraph{The surrogate $f$ (the ``$f_6$'' margin objective).} Among seven candidates they measured, the best is the logit-difference margin
\[
  f(x') \;=\; \max\Big(\max_{i\neq t} Z(x')_i - Z(x')_t,\; -\kappa\Big),
\]
where $Z(\cdot)$ are the \emph{logits} (not softmax probabilities), $t$ the target class, and $\kappa\ge0$ a confidence parameter. This $f$ is $\le0$ exactly when the target logit beats all others, so ``$f\le0$ $\Leftrightarrow$ misclassified.'' Using logits rather than softmax avoids vanishing gradients, the bug that broke an earlier defense (defensive distillation), and \citet{carlini2017towards} found cross-entropy the \emph{worst} surrogate and this logit-difference the best, because it keeps the ``be small'' and ``be misclassified'' terms balanced across values of $c$.

\paragraph{The change-of-variables (box) trick.} The constraint $x+\delta\in[0,1]^n$ is awkward. C\&W eliminate it by writing each coordinate as
\[
  \delta_i \;=\; \tfrac12\big(\tanh(w_i)+1\big) - x_i ,
\]
and optimizing freely over the new variable $w$. Since $\tanh\in(-1,1)$, the value $x_i+\delta_i=\tfrac12(\tanh(w_i)+1)$ is automatically in $(0,1)$, so any $w$ gives a valid image and a smooth, unconstrained problem solvable by a standard optimizer (Adam).

\paragraph{Algorithm.}
\begin{enumerate}[label=\textbf{\arabic*.},nosep]
  \item Choose a value of $c$ (binary search over a log-spaced range).
  \item Optimize Eq.~\eqref{eq:cw} over $w$ (with the $\tanh$ substitution) to convergence, recording the best adversarial $\delta$.
  \item Update $c$: increase if no adversarial example was found, decrease if one was; repeat.
  \item Return the smallest-norm adversarial $\delta$ over all $c$.
\end{enumerate}

\begin{example}[Reading $\|\delta\|^2 + c\,f$ for small vs.\ large $c$]
Take a 1D two-class model where the runner-up-minus-target logit gap, as a function of the scalar perturbation $\delta$, is $\,g(\delta)=1-\delta\,$ (so $f_6=\max(1-\delta,-\kappa)$ with $\kappa=0$ flips to ``misclassified'' once $\delta\ge1$). The objective is $J(\delta)=\delta^2 + c\,\max(1-\delta,0)$. For very small $c$ (say $c=0.1$): on $\delta<1$, $J=\delta^2+0.1(1-\delta)$ is minimized near $\delta\approx0.05$, which is \emph{not} adversarial ($\delta<1$): too small a $c$ buys a tiny but useless $\delta$. For large $c$ (say $c=10$): the penalty dominates until $\delta\ge1$, pushing the minimizer to $\delta=1$, the smallest adversarial value. Sweeping $c$ traces from ``too small to flip'' to ``flips with minimal norm,'' which is exactly what the binary search automates; \citet{carlini2017towards} show this as their Figure~2, where the attack rarely succeeds for $c<0.1$ and always succeeds for $c>1$.
\end{example}

\paragraph{Exercise.} With $J(\delta)=\delta^2+c\max(1-\delta,0)$ as above, find the threshold value of $c$ at which the unconstrained minimizer first becomes adversarial ($\delta^\star\ge1$). \emph{(Hint: on $\delta<1$ the derivative is $2\delta-c$; the interior minimizer is $\delta=c/2$, which reaches $1$ when $c=2$. For $c\ge2$ the minimizer sits at the kink $\delta=1$.)}

\subsection{DDN: Decoupled Direction and Norm}
\label{sub:ddn}

\paragraph{Definition and intuition.} C\&W is slow: a binary search over $c$, each with thousands of iterations. The Decoupled Direction and Norm attack \citep{rony2019decoupling} reaches the same minimum-norm goal in about $100$ iterations, cheap enough to run \emph{inside} adversarial training. The idea in the name: decouple the \emph{direction} of the perturbation (the unit-normalized gradient) from its \emph{norm} (a simple up/down rule plus a projection onto an $\epsilon$-sphere). There is no penalty constant $c$ at all.

\paragraph{The norm update.} At each step DDN checks whether the current point is already adversarial. If it is, the example is ``too easy,'' so \emph{shrink} the radius: $\epsilon\leftarrow(1-\gamma)\epsilon$. If it is not adversarial, \emph{grow} the radius: $\epsilon\leftarrow(1+\gamma)\epsilon$. The candidate is then placed on the sphere of the current radius, $\tilde x = x + \epsilon\,\delta/\|\delta\|_2$. The radius oscillates around the true boundary distance and converges to the minimum norm $r(x)$ \citep{rony2019decoupling}.

\paragraph{Algorithm \citep{rony2019decoupling}.}
\begin{enumerate}[label=\textbf{\arabic*.},nosep]
  \item Initialize $\delta=0$, $\tilde x = x$, $\epsilon=\epsilon_0$.
  \item For $k=1,\dots,K$: compute the gradient $g$ of the loss; take a normalized step $g\leftarrow\alpha\,g/\|g\|_2$ and update $\delta\leftarrow\delta+g$.
  \item If $\tilde x$ is adversarial, set $\epsilon\leftarrow(1-\gamma)\epsilon$; else $\epsilon\leftarrow(1+\gamma)\epsilon$.
  \item Project: $\tilde x \leftarrow x + \epsilon\,\delta/\|\delta\|_2$, then clip to the valid range.
  \item Return the lowest-norm $\tilde x$ found that is adversarial.
\end{enumerate}

\begin{example}[Oscillating radius by hand]
Take a classifier whose boundary is the vertical line at distance $1$ from $x$ along the attack direction, and set $\epsilon_0=1$, $\gamma=0.2$, with the step always pointing straight at the boundary. Iterate the radius rule, checking ``adversarial?'' $=$ ``radius $\ge1$?'':
\begin{itemize}[nosep]
  \item Start $\epsilon=1.0$: on the boundary, treat as adversarial $\to$ shrink to $1.0\cdot0.8 = 0.8$.
  \item $\epsilon=0.8$: not adversarial ($0.8<1$) $\to$ grow to $0.8\cdot1.2 = 0.96$.
  \item $\epsilon=0.96$: not adversarial $\to$ grow to $0.96\cdot1.2 = 1.152$.
  \item $\epsilon=1.152$: adversarial $\to$ shrink to $1.152\cdot0.8 = 0.9216$.
\end{itemize}
The radius bounces $1.0,0.8,0.96,1.152,0.9216,\dots$ around the true distance $1$, the swings narrowing as $\gamma$ effectively shrinks the relative gap; the converged value is the robust radius $r(x)=1$. Contrast with C\&W, which would binary-search a penalty constant $c$ to find the same number much more slowly \citep{rony2019decoupling}.
\end{example}

\paragraph{Exercise.} With the same setup but $\gamma=0.1$ and $\epsilon_0=1.5$, write the first four radii and state the limit. \emph{(Answer: $1.5\to1.35\to1.215\to1.0935\to\dots$, all $>1$ so all shrink, converging down to $r(x)=1$.)}

\subsection{AutoAttack: a standardized ensemble}
\label{sub:autoattack}

\paragraph{What it is.} A recurring failure in the field is reporting robustness using a weak attack: the model looks robust only because the attack was bad (a problem called \emph{gradient masking}). AutoAttack \citep{croce2020reliable} is the community's response: a \emph{parameter-free} ensemble of four diverse attacks, run together as a single standardized robustness test, so that no hand-tuning can inflate the result.

\paragraph{The four components.} We describe the ensemble at the level of its four members \citep{croce2020reliable}:
\begin{enumerate}[label=\textbf{\arabic*.},nosep]
  \item \textbf{APGD-CE}, Auto-PGD with the cross-entropy loss. Auto-PGD is PGD with an \emph{adaptive} step size (no step-size tuning) and a momentum term; it spends a budget of iterations exploring then exploiting, halving the step when progress stalls.
  \item \textbf{APGD-DLR} (also a targeted variant, APGD-T), Auto-PGD with the \emph{Difference-of-Logits-Ratio} loss. The DLR loss is used because cross-entropy is sensitive to a mere rescaling of the logits (which can cause gradient masking), whereas DLR is scale-invariant.
  \item \textbf{FAB}, a minimum-norm attack (the Fast Adaptive Boundary attack).
  \item \textbf{Square}, a \emph{black-box}, gradient-free attack (random search over square-shaped perturbations), included so that the ensemble still works even when gradients are uninformative.
\end{enumerate}
A point counts as robust under AutoAttack only if it survives \emph{all four}. The finer internals of FAB and Square are beyond what a primer needs; the four-member structure is the part that matters for understanding why the ensemble is hard to fool.

\paragraph{Why this is the right way to measure.} AutoAttack ties the section together: it is a fixed-$\epsilon$ \emph{evaluation} that combines white-box first-order attacks (the APGD pair), a minimum-norm attack (FAB), and a gradient-free attack (Square), precisely so that the reported $\mathrm{RobAcc}(\epsilon)$ is a trustworthy \emph{upper bound} on the model's true robustness (Section~\ref{sub:cert-vs-attack}) rather than an artifact of a weak optimizer.

\paragraph{Exercise.} A paper reports $90\%$ robust accuracy using only FGSM. Using the taxonomy of Section~\ref{sub:taxonomy} and the certificate/attack duality of Section~\ref{sub:cert-vs-attack}, explain why an AutoAttack re-evaluation can only \emph{lower} this number, never raise it, and why that makes AutoAttack a more honest measurement. \emph{(Answer: every attack gives an upper bound on robustness; a stronger or more diverse attack can find perturbations FGSM missed, decreasing the surviving fraction toward the truth, but it can never make a broken point survive, so the number can only fall.)}

% =====================================================================
%  CITATION KEYS USED (for the bibliography assembler)
% =====================================================================
% Keys that ALREADY EXIST in paper/report/references.bib:
%   goodfellow2015explaining   (FGSM — Goodfellow, Shlens, Szegedy 2015)
%   madry2018towards           (PGD / saddle point — Madry et al. 2018)
%   rony2019decoupling         (DDN — Rony et al. 2019)
%   croce2020reliable          (AutoAttack — Croce & Hein 2020)
%   hein2017formal             (Cross-Lipschitz / exact linear radius — Hein & Andriushchenko 2017)
%   tsuzuku2018lipschitz       (Lipschitz-Margin certificate, Prop. 1 — Tsuzuku, Sato, Sugiyama 2018)
%
% Keys this fragment USES that are NOT YET in references.bib — PLEASE ADD:
%   szegedy2014intriguing      (Szegedy et al., "Intriguing Properties of Neural Networks",
%                               ICLR 2014; arXiv:1312.6199 — discovery of adversarial examples +
%                               origin of the product-of-spectral-norms Lipschitz bound, Sec 4.3)
%   carlini2017towards         (Carlini & Wagner, "Towards Evaluating the Robustness of Neural
%                               Networks", IEEE S&P 2017; arXiv:1608.04644 — C&W attack)
%   scaman2018lipschitz        (Scaman & Virmaux, "Lipschitz regularity of deep neural networks:
%                               analysis and efficient estimation", NeurIPS 2018; arXiv:1805.10965
%                               — L = sup gradient norm; NP-hardness; AutoLip)
%   kolter2018adversarial      (Kolter & Madry, "Adversarial Robustness: Theory and Practice",
%                               NeurIPS 2018 tutorial, https://adversarial-ml-tutorial.org —
%                               demo-first arc, MNIST 0/1 numbers, norm-ball+optimizer taxonomy,
%                               Danskin discussion, dual-norm linear-attack derivation)
%   wiki_lipschitz             (Wikipedia, "Lipschitz continuity" — secant-slope picture,
%                               |x| is 1-Lipschitz, x^2 is not. Replace with a textbook cite,
%                               e.g. Rudin, if a Wikipedia citation is undesired.)
%
% Honesty notes for the assembler:
%  - AutoAttack internals are stated ONLY at the four-component level, per the DIGEST caveat that
%    the available PDF was read for the component list only (Sec 8 caveat).
%  - Danskin's theorem is STATED, not proved, as instructed.
%  - The Tsuzuku sqrt(2) proof and Lemma (max-of-rest 1-Lipschitz) are faithful to Tsuzuku 2018
%    Appendix A (verified against the PDF). The Hein exact-linear radius is Hein-Andriushchenko
%    2017 Thm 2.1 / the linear corollary (verified against the PDF).
% =====================================================================

% =====================================================================
%  PRIMER PART I -- THE INVARIANCE CLUSTER
%  Four background sections for a strong undergraduate:
%    (1) orthogonal projectors,
%    (2) the discrete Fourier transform and circulant/shift structure,
%    (3) group actions, invariance/equivariance, group averaging,
%    (4) shift-invariant descriptors: power spectrum and bispectrum.
%
%  This fragment supplies NO preamble: the assembler provides
%  amsmath/amsthm and the theorem-like environments
%  (definition, theorem, lemma, corollary, proposition, example,
%   problem, remark) and the macros \R \Z \one \sign \conv \Orb.
%
%  Notation conventions are FIXED here, once, to the engineering/forward
%  DFT convention of the project digest (see the "Notation" box below):
%      omega = e^{-2 pi i / N},  F_{jk} = omega^{jk},  DFT(x) = F x.
%  All worked numbers were checked by hand and by computer.
% =====================================================================

\section{Orthogonal projectors: the one linear-algebra tool everything rests on}
\label{sec:projectors}

Almost every later idea in this primer turns out to be the same gesture: take an
object, throw away the part of it that some symmetry can move around, and keep only
the part the symmetry leaves alone. The mathematical name for ``keep only the part a
subspace can see'' is \emph{orthogonal projection}. We start here so that the rest of
the document can lean on one clean theorem instead of re-deriving it three times.

We work in a real or complex inner-product space $V$ (think of $V = \R^n$ or $V =
\mathbb{C}^n$ with the usual dot product). The \emph{inner product} of two vectors is
written $\langle u, v\rangle$; in $\R^n$ it is $\langle u,v\rangle = \sum_k u_k v_k$,
and in $\mathbb{C}^n$ it is $\langle u,v\rangle = \sum_k u_k \overline{v_k}$, where
$\overline{c}$ is the \emph{complex conjugate} of $c$ (if $c = a + bi$ then
$\overline{c} = a - bi$). The \emph{norm} is $\|v\| = \sqrt{\langle v,v\rangle}$, and
two vectors are \emph{orthogonal}, written $u \perp v$, when $\langle u,v\rangle = 0$.
Throughout, an \emph{operator} means a linear map $P : V \to V$, which we may identify
with a square matrix once a basis is fixed.

\subsection{Two adjectives: idempotent and self-adjoint}

\begin{definition}[Idempotent]
\label{def:idempotent}
An operator $P : V \to V$ is \emph{idempotent} if $P^2 = P$, i.e.\ $P(Pv) = Pv$ for
every $v \in V$. In words: applying $P$ twice does exactly what applying it once does.
\end{definition}

The intuition is that $P$ \emph{lands} every vector somewhere, and once a vector has
already landed, $P$ leaves it alone. The set of places vectors can land is the
\emph{range} (or image) $\operatorname{range}(P) = \{Pv : v \in V\}$; the set of
vectors $P$ sends to $0$ is the \emph{kernel} (or null space) $\ker(P) = \{v : Pv =
0\}$. Idempotence says precisely that $P$ acts as the identity on its own range: if
$w = Pv$ is in the range, then $Pw = P(Pv) = Pv = w$.

\begin{example}[A $2\times 2$ idempotent that is \emph{not} a clean projection]
\label{ex:oblique}
Take, for a fixed number $\alpha \neq 0$,
\[
  P \;=\; \begin{bmatrix} 0 & 0 \\ \alpha & 1 \end{bmatrix}.
\]
A direct multiplication gives
\[
  P^2 \;=\; \begin{bmatrix} 0 & 0 \\ \alpha & 1 \end{bmatrix}
            \begin{bmatrix} 0 & 0 \\ \alpha & 1 \end{bmatrix}
        \;=\; \begin{bmatrix} 0\cdot 0 + 0\cdot\alpha & 0\cdot 0 + 0\cdot 1 \\
                              \alpha\cdot 0 + 1\cdot\alpha & \alpha\cdot 0 + 1\cdot 1
              \end{bmatrix}
        \;=\; \begin{bmatrix} 0 & 0 \\ \alpha & 1 \end{bmatrix} \;=\; P,
\]
so $P$ is idempotent: it \emph{is} a projection. But it is a \emph{skewed} one. Its
range is the line spanned by $(0,1)$ and its kernel is the line spanned by
$(1,-\alpha)$ (check: $P\,(1,-\alpha)^\top = (0, \alpha - \alpha)^\top = 0$). These two
lines are \emph{not} perpendicular when $\alpha \neq 0$. Such a $P$ is called an
\emph{oblique} projection. This example, due to Wikipedia's projection article
\citep{wiki_projection}, is the cautionary tale of the section: idempotence alone does
not make a projection ``drop straight down.'' One more property is needed.
\end{example}

\begin{definition}[Self-adjoint]
\label{def:selfadjoint}
An operator $P : V \to V$ is \emph{self-adjoint} if
\[
  \langle Pu, v \rangle \;=\; \langle u, Pv \rangle \qquad \text{for all } u,v \in V.
\]
In matrix terms this is $P = P^\top$ for real $V$ (the matrix equals its transpose)
and $P = P^\ast$ for complex $V$, where $P^\ast = \overline{P}^{\top}$ is the
\emph{conjugate transpose} (transpose, then conjugate every entry). Self-adjoint real
matrices are also called \emph{symmetric}.
\end{definition}

\begin{example}[The clean version of the same idea]
\label{ex:lineproj}
Let $a = (1,1)^\top$ and form
\[
  P \;=\; \frac{a\,a^\top}{a^\top a}
        \;=\; \frac{1}{2}\begin{bmatrix} 1 & 1 \\ 1 & 1 \end{bmatrix}.
\]
Then $P = P^\top$ (it is visibly symmetric), and $P^2 = \tfrac14
\left[\begin{smallmatrix} 2 & 2 \\ 2 & 2 \end{smallmatrix}\right] = P$. So $P$ is
\emph{both} idempotent and self-adjoint. Geometrically it projects every vector
straight down onto the line through $a$: for $x = (3,1)^\top$ we get $Px = (2,2)^\top$,
and the leftover piece $x - Px = (1,-1)^\top$ is orthogonal to $a$ (their dot product
is $1 - 1 = 0$). That is the behaviour we want, and the next theorem says it always
happens \citep{projectors_tamu,wiki_projection}.
\end{example}

\subsection{The central theorem}

\begin{theorem}[Idempotent {\boldmath$+$} self-adjoint {\boldmath$\Rightarrow$}
orthogonal projection]
\label{thm:orthoproj}
Let $P : V \to V$ be idempotent ($P^2 = P$) and self-adjoint ($P = P^\ast$). Then $P$
is the \emph{orthogonal} projection onto its range $W = \operatorname{range}(P)$:
for every $v$, the vector $Pv$ lies in $W$, the leftover $v - Pv$ is orthogonal to
all of $W$, and $Pv$ is the unique closest point of $W$ to $v$.
\end{theorem}

\begin{proof}
The whole proof rests on showing one orthogonality statement; everything else follows
from it. We follow Garrett's argument \citep{projectors_garrett}.

\emph{Step 1: the complementary operator.} Write $I$ for the identity and set $Q = I -
P$. Then $Q$ measures the leftover: $Qv = v - Pv$. Note $Q$ is also self-adjoint,
because $Q^\ast = (I - P)^\ast = I^\ast - P^\ast = I - P = Q$, using $P^\ast = P$.

\emph{Step 2: the key product vanishes.} Compute, for the operator $QP = (I-P)P$,
\begin{align*}
  QP \;&=\; (I - P)P \\
     \;&=\; IP - P^2 \\
     \;&=\; P - P \qquad\text{(here we used idempotence $P^2 = P$)} \\
     \;&=\; 0 .
\end{align*}
So $QP = 0$ as an operator: $Q$ annihilates everything $P$ produces.

\emph{Step 3: range of $P$ is orthogonal to range of $Q$.} Take any $v, w \in V$. We
test the inner product of an arbitrary range vector $Pv$ against an arbitrary leftover
vector $Qw$:
\begin{align*}
  \langle Pv,\; Qw \rangle
    \;&=\; \langle Pv,\; (I-P)w \rangle \\
    \;&=\; \langle (I-P)^\ast Pv,\; w \rangle
        &&\text{(move the operator to the other slot via the adjoint)} \\
    \;&=\; \langle (I-P) Pv,\; w \rangle
        &&\text{(since $(I-P)^\ast = I-P$ by Step 1)} \\
    \;&=\; \langle QP\, v,\; w \rangle \\
    \;&=\; \langle 0,\; w \rangle
        &&\text{(since $QP = 0$ by Step 2)} \\
    \;&=\; 0 .
\end{align*}
Because $v$ and $w$ were arbitrary, \emph{every} vector $Pv$ in the range of $P$ is
orthogonal to \emph{every} leftover $Qw = w - Pw$. In particular, taking $w$ to range
over all of $V$, the leftover $v - Pv$ is orthogonal to the entire range $W$.

\emph{Step 4: closest point (Pythagoras).} Fix $v$ and let $w$ be any point of $W$.
Since $Pv \in W$ and $w \in W$, also $Pv - w \in W$, hence by Step~3 it is orthogonal
to the leftover $v - Pv$. The Pythagorean law for orthogonal vectors
\citep{fourier_rutgers} then gives
\[
  \|v - w\|^2 \;=\; \|\,(v - Pv) + (Pv - w)\,\|^2
              \;=\; \|v - Pv\|^2 + \|Pv - w\|^2
              \;\ge\; \|v - Pv\|^2,
\]
with equality if and only if $\|Pv - w\|^2 = 0$, i.e.\ $w = Pv$. Therefore $Pv$ is the
unique point of $W$ nearest to $v$. This is exactly what ``$P$ projects $v$
orthogonally onto $W$'' means.
\end{proof}

\begin{remark}
The converse is also true and equally short: an orthogonal projection is automatically
self-adjoint \citep{projectors_garrett,wiki_projection}. We will only need the
direction proved above. The single fact you must carry forward is the slogan
\[
  \boxed{\ \text{idempotent} + \text{self-adjoint} \;\Longrightarrow\;
          \text{orthogonal projection (the closest-point map)}.\ }
\]
\end{remark}

\begin{corollary}[Eigenvalues of a projector are $0$ or $1$]
\label{cor:eig01}
If $P^2 = P$ and $Pv = \lambda v$ with $v \neq 0$, then $\lambda v = Pv = P^2 v =
\lambda(Pv) = \lambda^2 v$, so $\lambda = \lambda^2$, forcing $\lambda \in \{0,1\}$
\citep{wiki_projection}. The eigenvectors with $\lambda = 1$ span the range (what $P$
keeps); those with $\lambda = 0$ span the kernel (what $P$ discards).
\end{corollary}

\subsection{Spectral norm: how much an operator can stretch}

We also need one number that measures the largest stretching a linear map can do,
because later it controls how a shift or a perturbation can grow.

\begin{definition}[Spectral / operator norm]
\label{def:specnorm}
For a matrix $W$ the \emph{spectral norm} (also operator norm) is the largest factor by
which $W$ can lengthen a unit vector:
\[
  \|W\|_2 \;=\; \max_{\|z\| = 1} \|Wz\| .
\]
Equivalently $\|W\|_2 = \sigma_{\max}(W)$, the largest \emph{singular value} of $W$,
which equals $\sqrt{\lambda_{\max}(W^\ast W)}$, the square root of the largest
eigenvalue of $W^\ast W$.
\end{definition}

The equality $\|W\|_2 = \sqrt{\lambda_{\max}(W^\ast W)}$ comes from maximizing the
\emph{Rayleigh quotient}: $\|Wz\|^2 = \langle Wz, Wz\rangle = \langle W^\ast W z,
z\rangle$, and over unit vectors this is largest exactly at the top eigenvector of the
self-adjoint matrix $W^\ast W$, where it equals $\lambda_{\max}(W^\ast W)$.

\begin{example}
For $W = \left[\begin{smallmatrix} 3 & 0 \\ 0 & 1 \end{smallmatrix}\right]$ we have
$W^\ast W = \left[\begin{smallmatrix} 9 & 0 \\ 0 & 1 \end{smallmatrix}\right]$, whose
eigenvalues are $9$ and $1$, so $\|W\|_2 = \sqrt 9 = 3$. A useful sanity check on the
whole section: an \emph{orthogonal} projection never stretches, $\|P\|_2 \le 1$,
because $\|Pv\| \le \|v\|$ is just the inequality $\|v-Pv\|^2 \ge 0$ rearranged through
the Pythagoras identity of Step~4.
\end{example}

\begin{problem}
\label{prob:proj}
Build $P = a a^\top / (a^\top a)$ for $a = (2,1)^\top$, so $P = \tfrac15
\left[\begin{smallmatrix} 4 & 2 \\ 2 & 1 \end{smallmatrix}\right]$. Verify $P^2 = P$
and $P = P^\top$, and check that $P$ annihilates $(-1,2)^\top$ (which is orthogonal to
$a$). Then take the oblique $P' = \left[\begin{smallmatrix} 0 & 0 \\ 1 & 1
\end{smallmatrix}\right]$: confirm $P'^2 = P'$ but $P' \neq P'^\top$, find its range
and kernel, and observe they are not perpendicular. The moral: Theorem~\ref{thm:orthoproj}'s
two hypotheses are both load-bearing.
\end{problem}


% =====================================================================
\section{The discrete Fourier transform and the structure of shifts}
\label{sec:dft}

A previous version of this primer asserted, with no explanation, that ``shift operators
are diagonalized by the DFT.'' That sentence is true and it is the keystone of the
whole subject, but it is useless until every word in it is defined and the claim is
proved. This section does that, building up from scratch, following the
``discover the transform, don't postulate it'' approach of Bamieh
\citep{fourier_bamieh}, with the concrete computations of the MIT~18.06 notes
\citep{fourier_mit1806} and the rigorous theorem--proof spine of the Rutgers notes
\citep{fourier_rutgers}.

\paragraph{Notation, fixed once.}
We work with vectors of length $N$, indexed $0,1,\dots,N-1$, and read every index
\emph{modulo $N$} (so index $-1$ means $N-1$, index $N$ means $0$, and so on; we think
of the entries as sitting on a circle). The basic constant is the \emph{primitive
$N$-th root of unity}
\[
  \omega \;=\; e^{-2\pi i / N},
\]
chosen with a minus sign to match the engineering / forward DFT convention used throughout this
primer. It satisfies $\omega^N = 1$ but $\omega^m \neq 1$ for $0 < m < N$; its powers
$1, \omega, \omega^2, \dots, \omega^{N-1}$ are the $N$ solutions of $z^N = 1$, equally
spaced points on the unit circle. The \emph{discrete Fourier transform} (DFT) matrix
$F$ has entries
\[
  F_{jk} \;=\; \omega^{jk}, \qquad 0 \le j,k \le N-1,
  \qquad\text{and we write } \widehat{x} = \mathrm{DFT}(x) = F x .
\]
Two facts about $\omega$ we will use constantly: $\overline{\omega} = e^{+2\pi i/N} =
\omega^{-1}$ (conjugating flips the sign of the exponent), and the \emph{finite
geometric series}
\begin{equation}
\label{eq:geomsum}
  \sum_{m=0}^{N-1} \omega^{rm} \;=\;
  \begin{cases}
    N, & \text{if } r \equiv 0 \pmod N, \\[2pt]
    \dfrac{1 - \omega^{rN}}{1 - \omega^{r}} = \dfrac{1-1}{1-\omega^r} = 0, &
    \text{otherwise,}
  \end{cases}
\end{equation}
where the first case is ``add $N$ copies of $1$'' and the second uses $\omega^{rN} =
(\omega^N)^r = 1$ while $\omega^r \neq 1$. Equation~\eqref{eq:geomsum} is the engine
behind both Parseval and the diagonalization.

\subsection{What ``diagonalize'' means, and the shift operator}

\begin{definition}[Eigenvector, eigenvalue, diagonalize]
\label{def:diag}
A nonzero vector $v$ is an \emph{eigenvector} of a matrix $A$ with \emph{eigenvalue}
$\lambda$ if $A v = \lambda v$: the matrix only stretches $v$, never tilts it. To
\emph{diagonalize} $A$ is to find an invertible matrix $V$ (its columns a full set of
eigenvectors $v_1,\dots,v_N$) so that $V^{-1} A V = D$ is diagonal, with the
eigenvalues $\lambda_1,\dots,\lambda_N$ on the diagonal. In that eigenvector basis the
matrix acts by simple coordinatewise scaling. A family of matrices is
\emph{simultaneously diagonalized} by one $V$ if the same eigenvector basis works for
all of them at once.
\end{definition}

\begin{definition}[Cyclic shift]
\label{def:shift}
The \emph{cyclic (right) shift} $S$ moves every entry one step forward around the
circle:
\[
  (S x)_k \;=\; x_{k-1} \quad (\text{indices mod } N),
  \qquad\text{e.g.}\quad S\,(x_0,x_1,x_2,x_3) = (x_3,x_0,x_1,x_2).
\]
As a matrix $S$ is a \emph{permutation matrix} (a single $1$ in each row and column).
It is the most fundamental object of this whole section; we now find its eigenvectors.
\end{definition}

\subsection{The shift is diagonalized by the DFT}

\begin{theorem}[Eigenvectors of the shift]
\label{thm:shifteig}
Let $h_j$ denote the $j$-th column of the DFT matrix $F$, that is
\[
  h_j \;=\; \bigl(\omega^{0\cdot j},\ \omega^{1\cdot j},\ \omega^{2\cdot j},\ \dots,\
            \omega^{(N-1)j}\bigr)^{\!\top}
        \;=\; \bigl(1,\ \omega^{j},\ \omega^{2j},\ \dots,\ \omega^{(N-1)j}\bigr)^{\!\top}.
\]
Then each $h_j$ is an eigenvector of the cyclic shift $S$:
\[
  S\, h_j \;=\; \omega^{-j}\, h_j \qquad (j = 0,1,\dots,N-1),
\]
so the eigenvalues of $S$ are the $N$ roots of unity $\{\omega^{-j}\} = \{1, \omega^{-1},
\dots, \omega^{-(N-1)}\}$, the same set as $\{1, \omega, \dots, \omega^{N-1}\}$ in a
different order.
\end{theorem}

\begin{proof}
We chase one entry. The $k$-th entry of $h_j$ is $(h_j)_k = \omega^{kj}$. Apply the
definition of the shift, $(S\,h_j)_k = (h_j)_{k-1}$, and compute, for each $k$,
\begin{align*}
  (S\, h_j)_k \;&=\; (h_j)_{k-1} \\
              \;&=\; \omega^{(k-1)j}
                  &&\text{(definition of $h_j$ at index $k-1$)} \\
              \;&=\; \omega^{kj}\,\omega^{-j}
                  &&\text{(law of exponents $\omega^{(k-1)j} = \omega^{kj}\cdot\omega^{-j}$)} \\
              \;&=\; \omega^{-j}\,(h_j)_k .
\end{align*}
This holds for every entry $k$ (the wrap-around at $k=0$ is automatic because the
exponent is read mod $N$ and $\omega^{N} = 1$). Therefore $S\,h_j = \omega^{-j} h_j$ as
vectors: $h_j$ is an eigenvector with eigenvalue $\omega^{-j}$. The $N$ eigenvalues
$\omega^{0}, \omega^{-1}, \dots, \omega^{-(N-1)}$ are distinct, so the $N$ eigenvectors
$h_0,\dots,h_{N-1}$ are linearly independent and form a basis: the columns of $F$
diagonalize $S$.
\end{proof}

\begin{remark}[A note on conventions, so you are never confused later]
With the \emph{forward} convention $\omega = e^{-2\pi i/N}$ adopted here, the shift
eigenvalue attached to column $h_j$ is $\omega^{-j}$. Some textbooks
(e.g.\ \citep{fourier_mit1806,fourier_rutgers}) use the opposite sign $\omega =
e^{+2\pi i/N}$ and obtain $S h_j = \omega^{j} h_j$. The eigenvectors are the same lines,
the eigenvalues are the same set of roots of unity; only the bookkeeping label changes.
We will always state which convention is in force.
\end{remark}

\subsection{Circulant matrices: every shift-invariant linear map}

\begin{definition}[Circulant matrix and circular convolution]
\label{def:circulant}
A matrix $C$ is \emph{circulant} if each column is the cyclic shift of the previous one;
equivalently its entries depend only on $k - l \bmod N$, written $(C_a)_{kl} =
a_{k-l}$, where $a = (a_0,\dots,a_{N-1})$ is its first column. The action of a circulant
is exactly a \emph{circular convolution}: writing $a \conv x$ for the convolution,
\[
  (C_a x)_k \;=\; \sum_{l=0}^{N-1} a_{k-l}\, x_l \;=\; (a \conv x)_k \qquad
  (\text{indices mod } N).
\]
A convolution ``slides one signal across another and sums the overlaps.'' The shift
itself is the circulant whose first column is $(0,1,0,\dots,0)$.
\end{definition}

The reason circulants matter is that they are \emph{precisely} the matrices that
commute with the shift, and ``commutes with the shift'' is the algebraic spelling of
``treats every position the same way'' --- the defining feature of a translation-blind
operation \citep{fourier_bamieh,groups_bronstein}.

\begin{lemma}[Circulants are polynomials in the shift]
\label{lem:polyS}
A matrix $C$ is circulant if and only if it is a polynomial in $S$,
\[
  C \;=\; c_0 I + c_1 S + c_2 S^2 + \dots + c_{N-1} S^{N-1},
\]
where $(c_0,\dots,c_{N-1})$ is the first column of $C$. In particular every circulant
commutes with $S$, and any two circulants commute with each other.
\end{lemma}

\begin{proof}
Note first that $S^p$ is the circulant that shifts by $p$ places, $(S^p x)_k = x_{k-p}$.
Hence a polynomial $\sum_p c_p S^p$ sends $x$ to the vector with $k$-th entry $\sum_p
c_p x_{k-p}$, which is the circular convolution $(c \conv x)_k$ --- a circulant with
first column $c$. Conversely, given a circulant $C$ with first column $c$, the operator
$\sum_p c_p S^p$ has the same first column ($S^p e_0 = e_p$, so $\sum_p c_p S^p$ applied
to $e_0$ returns $c$) and, being shift-invariant, the same every column; so the two
matrices agree. Finally $S$ commutes with any polynomial in $S$, and two polynomials in
$S$ commute with each other because powers of $S$ do
\citep{fourier_rutgers}.
\end{proof}

\begin{theorem}[The DFT simultaneously diagonalizes all circulants;\\
the convolution theorem]
\label{thm:circdiag}
Let $C_a$ be the circulant with first column $a$, equivalently $C_a x = a \conv x$.
Then every column $h_j$ of $F$ is an eigenvector of $C_a$, and the eigenvalue is the
$j$-th DFT coefficient of $a$:
\[
  C_a\, h_j \;=\; \widehat{a}_j\, h_j,
  \qquad\text{where } \widehat{a}_j = (Fa)_j = \sum_{l=0}^{N-1} a_l\, \omega^{jl}.
\]
Consequently the DFT turns convolution into entrywise multiplication --- the
\emph{convolution theorem}:
\[
  \boxed{\ \mathrm{DFT}(a \conv x) \;=\; \mathrm{DFT}(a)\,\odot\,\mathrm{DFT}(x)\ }
  \qquad (\odot = \text{entrywise product}).
\]
\end{theorem}

\begin{proof}
\emph{Eigenvalues.} Write $C_a = \sum_p a_p S^p$ as in Lemma~\ref{lem:polyS}. Apply it
to the shift eigenvector $h_j$ and use $S h_j = \omega^{-j} h_j$
(Theorem~\ref{thm:shifteig}), hence $S^p h_j = \omega^{-pj} h_j$:
\[
  C_a\, h_j \;=\; \sum_{p=0}^{N-1} a_p\, S^p h_j
            \;=\; \sum_{p=0}^{N-1} a_p\, \omega^{-pj}\, h_j
            \;=\; \Bigl(\underbrace{\sum_{p=0}^{N-1} a_p\, \omega^{-pj}}_{=\,(Fa)_{-j}=\widehat a_{-j}}\Bigr)\, h_j .
\]
The bracket is a single number, so $h_j$ is an eigenvector. (With our forward
convention the eigenvalue on column $h_j$ is $\widehat a_{-j}$; the whole \emph{set} of
eigenvalues is $\{\widehat a_0, \dots, \widehat a_{N-1}\}$, just relabelled. The clean
``eigenvalue $=$ DFT of the first column'' statement holds verbatim with the
opposite-sign convention \citep{fourier_bamieh}; we keep ours and note the relabelling.)
Because the $h_j$ form a basis, $C_a$ is diagonal in that basis: this is what
``simultaneously diagonalized'' means, since the same columns of $F$ work for
\emph{every} choice of $a$.

\emph{Convolution theorem (direct index chase, following
\citep{fourier_rutgers,wiki_convolution}).} We compute the $k$-th DFT coefficient of
$a \conv x$ from the definitions:
\begin{align*}
  \mathrm{DFT}(a \conv x)_k
    \;&=\; \sum_{j=0}^{N-1} (a \conv x)_j\, \omega^{kj} \\
    \;&=\; \sum_{j=0}^{N-1} \Bigl(\sum_{l=0}^{N-1} a_{j-l}\, x_l\Bigr)\, \omega^{kj}
        &&\text{(definition of convolution)} \\
    \;&=\; \sum_{l=0}^{N-1} x_l \sum_{j=0}^{N-1} a_{j-l}\, \omega^{kj}
        &&\text{(swap the two finite sums)} \\
    \;&=\; \sum_{l=0}^{N-1} x_l \sum_{m=0}^{N-1} a_{m}\, \omega^{k(m+l)}
        &&\text{(set $m = j - l$; still a full period)} \\
    \;&=\; \sum_{l=0}^{N-1} x_l\, \omega^{kl} \sum_{m=0}^{N-1} a_{m}\, \omega^{km}
        &&\text{(factor $\omega^{k(m+l)} = \omega^{km}\,\omega^{kl}$)} \\
    \;&=\; \Bigl(\sum_{l} x_l \omega^{kl}\Bigr)\Bigl(\sum_{m} a_m \omega^{km}\Bigr)
        \;=\; \widehat{x}_k\,\widehat{a}_k .
\end{align*}
That is exactly the entrywise product $\widehat a \odot \widehat x$ at index $k$. The
substitution $m = j - l$ is legitimate because both $a$ and the exponential are
$N$-periodic in their index, so summing $j$ over one full period is the same as summing
$m$ over one full period.
\end{proof}

\subsection{Parseval: the DFT preserves energy (up to a factor of $N$)}

\begin{theorem}[Parseval / Plancherel]
\label{thm:parseval}
With the forward (unnormalized) DFT, for every vector $x \in \mathbb{C}^N$,
\[
  \|\mathrm{DFT}(x)\|^2 \;=\; N\,\|x\|^2,
  \qquad\text{equivalently}\qquad \|x\|^2 = \tfrac1N\,\|\mathrm{DFT}(x)\|^2 .
\]
(With the \emph{unitary} normalization $\tfrac{1}{\sqrt N}F$ it reads $\|x\| =
\|\mathrm{DFT}(x)\|$ exactly.)
\end{theorem}

\begin{proof}
The proof is the orthogonality of the columns of $F$, which is exactly the geometric
series \eqref{eq:geomsum}. Take two columns $h_j, h_k$ of $F$ and form their inner
product (recall $\langle u,v\rangle = \sum_m u_m \overline{v_m}$, and
$\overline{\omega^{km}} = \omega^{-km}$):
\begin{align*}
  \langle h_j,\, h_k\rangle
    \;&=\; \sum_{m=0}^{N-1} \omega^{jm}\,\overline{\omega^{km}}
        \;=\; \sum_{m=0}^{N-1} \omega^{jm}\,\omega^{-km}
        \;=\; \sum_{m=0}^{N-1} \omega^{(j-k)m}
        \;=\;
        \begin{cases} N, & j = k, \\ 0, & j \neq k, \end{cases}
\end{align*}
where the last step is \eqref{eq:geomsum} with $r = j-k$. In matrix form this says
$F^\ast F = N I$, i.e.\ the columns are orthogonal each of squared length $N$. Now
expand the DFT in this basis. Writing $\widehat x = Fx$,
\[
  \|\widehat x\|^2
    \;=\; \langle Fx,\, Fx\rangle
    \;=\; \langle F^\ast F x,\, x\rangle
    \;=\; \langle N x,\, x\rangle
    \;=\; N\,\|x\|^2 ,
\]
where we moved $F$ to the other slot of the inner product through its adjoint $F^\ast$
and then used $F^\ast F = N I$. Dividing by $N$ gives the second form.
\end{proof}

\subsection{Worked examples at {\boldmath$N=4$} (do these by hand once)}

With $N = 4$ the root of unity is $\omega = e^{-2\pi i/4} = -i$, and the forward DFT
matrix is
\[
  F_4 \;=\;
  \begin{bmatrix}
    1 & 1 & 1 & 1 \\
    1 & -i & -1 & i \\
    1 & -1 & 1 & -1 \\
    1 & i & -1 & -i
  \end{bmatrix}.
\]
(The entry $F_{jk} = (-i)^{jk}$; the MIT~18.06 notes \citep{fourier_mit1806} print the
complex-conjugate version because they use the $+$ sign convention.)

\begin{example}[Convolution two ways]
Let $a = (1,1,0,0)$ and $x = (1,2,3,4)$. Direct circular convolution: $(a\conv x)_k =
\sum_l a_{k-l} x_l$ with $a_0 = a_1 = 1$, $a_2 = a_3 = 0$, so $(a\conv x)_k = x_k +
x_{k-1}$:
\[
  a \conv x = (x_0 + x_3,\ x_1 + x_0,\ x_2 + x_1,\ x_3 + x_2)
            = (1+4,\ 2+1,\ 3+2,\ 4+3) = (5,3,5,7).
\]
Via the transform: $\widehat a = F_4 a = (2,\,1-i,\,0,\,1+i)$ and $\widehat x = F_4 x =
(10,\,-2+2i,\,-2,\,-2-2i)$. Their entrywise product is $\widehat a \odot \widehat x =
(20,\,(1-i)(-2+2i),\,0,\,(1+i)(-2-2i)) = (20,\,4i,\,0,\,-4i)$, hmm let us simplify:
$(1-i)(-2+2i) = -2+2i+2i-2i^2 = -2+4i+2 = 4i$ and similarly the last entry is $-4i$.
Applying the inverse DFT, $\tfrac14 F_4^\ast(20,4i,0,-4i)$, returns $(5,3,5,7)$,
matching the direct computation \citep{fourier_rutgers}. The two routes agree exactly,
which is the convolution theorem at work.
\end{example}

\begin{example}[Parseval check]
For $x = (1,2,-1,0)$ we have $\|x\|^2 = 1 + 4 + 1 + 0 = 6$. Its DFT is $\widehat x = F_4
x = (2,\,2-2i,\,-2,\,2+2i)$, so $\|\widehat x\|^2 = |2|^2 + |2-2i|^2 + |-2|^2 +
|2+2i|^2 = 4 + 8 + 4 + 8 = 24 = 4 \cdot 6 = N\|x\|^2$, exactly as
Theorem~\ref{thm:parseval} predicts.
\end{example}

\begin{problem}
\label{prob:dft}
Show by hand that the $4\times 4$ cyclic shift $S$ has eigenvalues equal to the $4$th
roots of unity $\{1, i, -1, -i\}$, with eigenvectors the columns of $F_4$. Then verify
Parseval for $x = (1,2,-1,0)$ (the example above) and for one vector of your own
choosing. Finally, confirm the convolution theorem on $a = (1,1,0,0)$, $x = (1,2,3,4)$
by reproducing both routes in the previous example.
\end{problem}


% =====================================================================
\section{Symmetry made precise: group actions, orbits, invariance,
and group averaging}
\label{sec:groups}

The DFT story was about \emph{one} symmetry, the cyclic shift. To talk about
shift-invariance and adversarial robustness in general we need the vocabulary that
makes ``a symmetry acting on data'' precise. The right language is that of group
actions, and the right tool is the operation that ``averages a function over a
symmetry'' --- the classical Reynolds operator. We follow the geometric deep learning
treatment of Bronstein, Bruna, Cohen and Veli\v{c}kovi\'c \citep{groups_bronstein} and
the Wikipedia articles on group actions and equivariant maps
\citep{wiki_groupaction,wiki_equivariant}, downgrading their integrals over abstract
groups to finite sums an undergraduate can compute.

\subsection{Groups, actions, orbits}

\begin{definition}[Group]
\label{def:group}
A \emph{group} $(G,\cdot)$ is a set $G$ with a multiplication $\cdot$ that is
associative, has an \emph{identity} element $e$ (with $e g = g e = g$ for all $g$), and
gives every $g$ an \emph{inverse} $g^{-1}$ (with $g g^{-1} = g^{-1} g = e$). The group is
\emph{abelian} (commutative) if $g h = h g$ always. The running example is the
\emph{cyclic group} $\Z_N = \{0,1,\dots,N-1\}$ under addition mod $N$, which is abelian
and is generated by the single element $1$ (every element is a repeated sum of $1$'s).
\end{definition}

\begin{definition}[Group action]
\label{def:action}
A (left) \emph{action} of a group $G$ on a set $X$ is a rule $G \times X \to X$, written
$(g,x) \mapsto g \cdot x$, satisfying two axioms \citep{wiki_groupaction}:
\[
  \text{(identity)}\quad e \cdot x = x, \qquad\qquad
  \text{(compatibility)}\quad g \cdot (h \cdot x) = (g h) \cdot x .
\]
The identity axiom says ``doing nothing does nothing''; the compatibility axiom says
``doing $h$ then $g$ is the same as doing the product $g h$.'' When $X$ is a vector
space and each $g$ acts by a linear (in fact invertible) map, the action is called a
\emph{representation}; we write that map as a matrix $\rho(g)$, and compatibility reads
$\rho(g)\rho(h) = \rho(gh)$.
\end{definition}

\begin{example}[The cyclic shift action, the one we care about]
\label{ex:shiftaction}
Let $G = \Z_N$ act on signals $x \in \mathbb{C}^N$ by shifting: the group element $u$
acts by $\rho(u) = S^u$, so $(\,u \cdot x\,)_k = x_{k - u}$ (shift the signal forward by
$u$). The identity $u = 0$ does nothing; compatibility $S^u S^v = S^{u+v}$ holds because
shifting by $u$ then by $v$ shifts by $u + v$. This is the action that turns
``$\Z_N$-symmetry'' into ``the shift operators of Section~\ref{sec:dft}.''
\end{example}

\begin{definition}[Orbit]
\label{def:orbit}
The \emph{orbit} of a point $x$ is everything reachable from $x$ by the group,
\[
  \Orb(x) \;=\; G \cdot x \;=\; \{\, g \cdot x : g \in G \,\}.
\]
Orbits \emph{partition} $X$: every point is in exactly one orbit, because the relation
``$y = g\cdot x$ for some $g$'' is reflexive ($e$), symmetric ($g^{-1}$) and transitive
(compatibility). An orbit is the set of all configurations that the symmetry regards as
``the same up to the symmetry.''
\end{definition}

\begin{example}[Orbits of the shift on $\R^4$]
\label{ex:orbits4}
Under $\Z_4$ acting by shifts on $\R^4$: the orbit of $(1,0,0,0)$ is the four cyclic
shifts $\{(1,0,0,0),(0,1,0,0),(0,0,1,0),(0,0,0,1)\}$ --- the four standard basis
vectors. The orbit of the constant vector $(1,1,1,1)$ is just itself: shifting it
changes nothing, so it is a \emph{fixed point}. Already you can feel the
shift-invariance theme: a function that should not care about position must give the
same value everywhere along an orbit.
\end{example}

\subsection{Invariant versus equivariant}

This distinction is the single most important conceptual point of the section, and the
old primer never made it.

\begin{definition}[Invariant and equivariant functions]
\label{def:inveq}
Let $G$ act on a domain $X$ and (for equivariance) also on a codomain $Y$. A function
$f : X \to Y$ is
\[
  \textbf{invariant} \quad\text{if}\quad f(g\cdot x) = f(x) \ \ \forall g,
  \qquad\qquad
  \textbf{equivariant} \quad\text{if}\quad f(g\cdot x) = g\cdot f(x) \ \ \forall g .
\]
An invariant function gives the same output no matter how the input is transformed: it
\emph{collapses each orbit to a single value}. An equivariant function lets the
transformation pass through predictably: transform the input and the output transforms
the same way. Invariance is the special case of equivariance where the action on the
codomain is trivial \citep{wiki_equivariant}.
\end{definition}

\begin{example}[The textbook pictures]
\label{ex:areacentroid}
For triangles under the rigid motions of the plane (translations, rotations,
reflections): the \emph{area} is invariant --- moving a triangle does not change its
area. The \emph{centroid} is equivariant --- move the triangle and its centroid moves
to exactly the image of the old centroid \citep{wiki_equivariant}. In statistics, the
sample \emph{mean} is equivariant under rescaling of the data (change units, the mean
changes units the same way), while the \emph{median} is equivariant under \emph{any}
monotone relabelling. In machine learning the canonical pair is: image
\emph{classification} is shift-invariant (a cat is a cat wherever it sits), whereas
\emph{segmentation} is shift-equivariant (shift the image and the predicted mask shifts
identically) \citep{groups_bronstein}.
\end{example}

\begin{remark}[A warning the old primer should have given]
Convolutional layers are shift-\emph{equivariant}, not invariant; this is exactly the
statement $S\,C_a x = C_a\,S x$ for a circulant $C_a$, which you already proved when you
showed circulants commute with $S$ (Lemma~\ref{lem:polyS}). Invariance is only obtained
at the \emph{end}, by a pooling/averaging step that collapses the orbit. Many
signal-processing texts loosely call equivariance ``shift-invariance''; do not be
fooled \citep{groups_bronstein}.
\end{remark}

\subsection{The Reynolds (group-averaging) operator}

How does one \emph{manufacture} an invariant function from an arbitrary one? Average
over the group. This idea is classical (Reynolds; Weyl and Hilbert in invariant theory)
and is the symmetry layer's workhorse \citep{groups_bronstein}.

\begin{definition}[Group average / Reynolds operator]
\label{def:reynolds}
For a finite group $G$ acting on a vector space, define the \emph{group-averaging
operator} $R$ on functions (or, when $G$ acts linearly, on vectors) by
\[
  (R f)(x) \;=\; \frac{1}{|G|} \sum_{g \in G} f(g \cdot x),
  \qquad\text{or for the linear action on vectors:}\quad
  R \;=\; \frac{1}{|G|} \sum_{g \in G} \rho(g).
\]
For a finite group this finite sum is the exact counterpart of Bronstein et al.'s
continuous average $\tfrac{1}{\mu(G)} \int_G f(g\cdot x)\,d\mu(g)$ over the group's
Haar measure \citep{groups_bronstein}: for a finite group the Haar measure just counts
elements.
\end{definition}

\begin{theorem}[The group average is invariant, idempotent, and (for an orthogonal
action) an orthogonal projection onto the invariants]
\label{thm:reynolds}
Let $G$ be finite and act linearly. Then:
\begin{enumerate}[label=\textup{(\roman*)}]
  \item $R$ is \emph{invariant}: $R\,\rho(h) = R$ for every $h \in G$, i.e.\ $(Rf)(h\cdot
        x) = (Rf)(x)$.
  \item $R$ is \emph{idempotent}: $R^2 = R$. Its range is exactly the invariant
        subspace $V_{\mathrm{inv}} = \{v : \rho(g)v = v\ \forall g\}$.
  \item If every $\rho(g)$ is \emph{orthogonal} (or unitary), $\rho(g)^\ast =
        \rho(g)^{-1}$, then $R$ is also self-adjoint, hence by
        Theorem~\ref{thm:orthoproj} the \emph{orthogonal} projection onto
        $V_{\mathrm{inv}}$ --- the closest invariant vector to any given $v$.
\end{enumerate}
\end{theorem}

\begin{proof}
\emph{(i) Invariance, by re-indexing the group.} Fix $h \in G$ and compute, using
compatibility $\rho(g)\rho(h) = \rho(gh)$:
\begin{align*}
  R\,\rho(h)
    \;&=\; \frac{1}{|G|}\sum_{g \in G} \rho(g)\rho(h)
    \;=\; \frac{1}{|G|}\sum_{g \in G} \rho(g h)
    \;=\; \frac{1}{|G|}\sum_{g' \in G} \rho(g')
    \;=\; R .
\end{align*}
The middle step is the only subtlety: as $g$ runs once through all of $G$, the product
$g' = g h$ also runs \emph{once} through all of $G$ (right multiplication by the fixed
$h$ is a bijection of $G$, with inverse right multiplication by $h^{-1}$). So we are
summing the same $|G|$ terms in a different order. Hence $R$ is unchanged by any group
element on its input.

\emph{(ii) Idempotence and range.} First, $R$ fixes every invariant vector: if
$\rho(g)v = v$ for all $g$, then $Rv = \tfrac{1}{|G|}\sum_g \rho(g) v =
\tfrac{1}{|G|}\sum_g v = v$. Second, every output $Rv$ \emph{is} invariant: by part
(i), $\rho(h)\,Rv = R v$ would follow from $\rho(h) R = R$; let us verify that form too,
\[
  \rho(h) R = \frac{1}{|G|}\sum_g \rho(h)\rho(g) = \frac{1}{|G|}\sum_g \rho(hg)
            = \frac{1}{|G|}\sum_{g'} \rho(g') = R ,
\]
the same re-indexing as in (i) but with \emph{left} multiplication. So $Rv \in
V_{\mathrm{inv}}$ for every $v$, and $R$ fixes $V_{\mathrm{inv}}$ pointwise; therefore
$R(Rv) = Rv$, i.e.\ $R^2 = R$, and the range of $R$ is exactly $V_{\mathrm{inv}}$.

\emph{(iii) Self-adjointness for an orthogonal action.} Compute the adjoint of the sum,
using $(\rho(g))^\ast = \rho(g)^{-1} = \rho(g^{-1})$:
\[
  R^\ast = \frac{1}{|G|}\sum_g \rho(g)^\ast
         = \frac{1}{|G|}\sum_g \rho(g^{-1})
         = \frac{1}{|G|}\sum_{g'} \rho(g')
         = R ,
\]
where again $g' = g^{-1}$ ranges over all of $G$ as $g$ does (inversion is a bijection
of $G$). So $R = R^\ast$. Now $R$ is idempotent (part ii) and self-adjoint, so
Theorem~\ref{thm:orthoproj} applies: $R$ is the orthogonal projection onto its range
$V_{\mathrm{inv}}$, the unique nearest invariant vector. This is the rigorous content of
the slogan ``averaging over the symmetry returns the closest symmetric thing.''
\end{proof}

\subsection{The cyclic case, and why linear invariants are too weak}

\begin{example}[Shift-averaging returns the mean]
\label{ex:shiftavg}
Apply Theorem~\ref{thm:reynolds} to $G = \Z_N$ acting on $\mathbb{C}^N$ by shifts,
$\rho(u) = S^u$. The group average is
\[
  (R x)_k = \frac{1}{N}\sum_{u=0}^{N-1} (S^u x)_k
          = \frac{1}{N}\sum_{u=0}^{N-1} x_{k-u}
          = \frac{1}{N}\sum_{j=0}^{N-1} x_j
          = \text{mean}(x),
\]
the constant vector all of whose entries equal the average of $x$. Each $S^u$ is a
permutation matrix, hence orthogonal, so by part~(iii) this $R$ is the orthogonal
projection onto the one-dimensional invariant subspace spanned by $(1,1,\dots,1)$ ---
which is exactly the orbit of fixed points from Example~\ref{ex:orbits4}. Connecting to
Section~\ref{sec:dft}: $R = \tfrac1N\sum_u S^u$ is a circulant, and by
Theorem~\ref{thm:circdiag} its eigenvalues are the DFT of its first column
$\tfrac1N(1,1,\dots,1)$, which is $(1,0,0,\dots,0)$ --- it keeps only the ``zero
frequency'' (the mean) and kills everything else.
\end{example}

The example carries a moral that motivates the entire next section. The \emph{only}
linear shift-invariant feature of a signal is its mean (any other shift-invariant linear
functional must be a multiple of the sum, by the same averaging argument). Bronstein et
al.\ make the same point vividly: for images under translation, a linear invariant can
see nothing but the \emph{average colour} \citep{groups_bronstein}. To get richer
invariants --- ones that can actually tell two signals apart while ignoring their
position --- we must go \emph{nonlinear}. That is precisely what the power spectrum and
bispectrum do.

\begin{problem}
\label{prob:reynolds}
Show that under $\Z_N$ the group average of any vector equals the constant vector whose
entries are the mean; verify directly that $R^2 = R$ on a length-$3$ example; and argue
that the only shift-invariant \emph{linear} functionals $x \mapsto \langle w, x\rangle$
are the multiples of the sum $\sum_k x_k$ (hint: invariance forces $S^\ast w = w$, so
$w$ is itself shift-invariant). Conclude that nontrivial shift-invariants must be
nonlinear.
\end{problem}


% =====================================================================
\section{Shift-invariant descriptors: the power spectrum and the bispectrum}
\label{sec:invariants}

We now build the two nonlinear shift-invariants that the project actually uses. Both
are computed from the DFT, both throw away absolute position, and the difference between
them --- one discards phase, the other keeps relative phase --- is the punchline. The
sources are the MIT~6.011 notes on power spectral density \citep{invariants_mit6011},
de~Buyl's Wiener--Khinchin derivation \citep{invariants_wienerkhinchin}, Kakarala's
introduction to bispectral invariants \citep{invariants_kakarala}, and Chen, Zehni and
Zhao's discrete treatment \citep{invariants_chen}.

\subsection{The DFT shift theorem: shifting only rotates phase}

\begin{lemma}[Shift theorem]
\label{lem:shifttheorem}
Let $R_\tau$ shift a signal by $\tau$, $(R_\tau x)_n = x_{n-\tau}$. Then
\[
  \mathrm{DFT}(R_\tau x)_k \;=\; \widehat{x}_k\, \omega^{k\tau}
                          \;=\; \widehat{x}_k\, e^{-2\pi i k \tau / N}.
\]
A shift leaves the \emph{magnitude} $|\widehat x_k|$ of every Fourier coefficient
untouched and only multiplies it by a \emph{unit-modulus phase} $e^{-2\pi i k\tau/N}$
(a complex number of absolute value $1$).
\end{lemma}

\begin{proof}
Compute from the definition, substituting $m = n - \tau$ and using periodicity:
\begin{align*}
  \mathrm{DFT}(R_\tau x)_k
    \;&=\; \sum_{n=0}^{N-1} (R_\tau x)_n\, \omega^{kn}
    \;=\; \sum_{n=0}^{N-1} x_{n-\tau}\, \omega^{kn}
    \;=\; \sum_{m=0}^{N-1} x_{m}\, \omega^{k(m+\tau)} \\
    \;&=\; \omega^{k\tau}\sum_{m=0}^{N-1} x_m\,\omega^{km}
    \;=\; \omega^{k\tau}\,\widehat x_k. \qedhere
\end{align*}
\end{proof}

Everything below is a consequence of one observation: \emph{position lives entirely in
the phase}. Any quantity built from the $\widehat x_k$ that cannot detect a
phase-rotation $e^{-2\pi i k\tau/N}$ will be shift-invariant.

\subsection{Power spectrum and Wiener--Khinchin}

\begin{definition}[Autocorrelation and power spectrum]
\label{def:powerspec}
The \emph{circular autocorrelation} of $x$ is $r_\tau = \sum_{t=0}^{N-1} x_t\,
x_{t+\tau}$ (for real signals; conjugate the second factor for complex ones), measuring
how much $x$ resembles its own shift by $\tau$. The \emph{power spectrum} is the squared
magnitude of the DFT,
\[
  P_k \;=\; |\widehat x_k|^2 \;=\; \widehat x_k\,\overline{\widehat x_k}.
\]
It is real, nonnegative, and (for real $x$) symmetric \citep{invariants_mit6011}.
\end{definition}

\begin{theorem}[Power spectrum is shift-invariant]
\label{thm:psinv}
For every shift $\tau$, $\ |\mathrm{DFT}(R_\tau x)_k|^2 = |\widehat x_k|^2$.
\end{theorem}

\begin{proof}
Using the shift theorem (Lemma~\ref{lem:shifttheorem}) and $|e^{-2\pi i k\tau/N}| = 1$:
\[
  |\mathrm{DFT}(R_\tau x)_k|^2
    = \bigl|\,\widehat x_k\, e^{-2\pi i k\tau/N}\,\bigr|^2
    = |\widehat x_k|^2 \cdot \bigl|e^{-2\pi i k\tau/N}\bigr|^2
    = |\widehat x_k|^2 \cdot 1
    = |\widehat x_k|^2 .
\]
The phase, which is where the position information sat, is annihilated by the
absolute-value square. Hence the power spectrum is the same for $x$ and every shift of
$x$: it is a shift-invariant descriptor.
\end{proof}

\begin{theorem}[Wiener--Khinchin: power spectrum $=$ transform of autocorrelation]
\label{thm:wk}
The power spectrum is the DFT of the autocorrelation: $\widehat r = P$, equivalently the
autocorrelation is the inverse DFT of the power spectrum.
\end{theorem}

\begin{proof}[Proof skeleton (continuous version, following de~Buyl
\citep{invariants_wienerkhinchin})]
Start from the autocorrelation
\[
  C(\tau) = \lim_{T\to\infty} \tfrac1T \int_0^T x(t)\, x(t+\tau)\,dt
\]
and replace each $x$ by its inverse Fourier transform. After exchanging
the order of integration one is left with a time-average of $e^{i(\omega+\omega')t}$:
that average is $0$ unless $\omega + \omega' = 0$ (the oscillations cancel to order
$1/T$) and equals $1$ when $\omega + \omega' = 0$. Keeping only the surviving
$\omega' = -\omega$ terms and using realness $\widetilde x(-\omega) =
\overline{\widetilde x(\omega)}$ gives
\[
  C(\tau) = \frac{1}{2\pi}\int e^{i\omega\tau}\,\widetilde x(\omega)\,
            \overline{\widetilde x(\omega)}\,d\omega
          = \frac{1}{2\pi}\int e^{i\omega\tau}\,|\widetilde x(\omega)|^2\,d\omega,
\]
i.e.\ the autocorrelation is the inverse transform of $|\widetilde x|^2$, the power
spectrum. The discrete version replaces the integrals by the length-$N$ DFT; the random
(noisy) version is handled by MIT~6.011's periodogram argument
\citep{invariants_mit6011}.
\end{proof}

Why does this matter? Because it shows the power spectrum is a \emph{complete summary of
the second-order structure} of a signal --- but \emph{only} the second order. By
discarding all phase it cannot tell apart two signals whose Fourier magnitudes agree;
Kakarala's Figure~1 shows two genuinely different shapes with identical power spectra,
the power spectrum ``fooled'' \citep{invariants_kakarala}. To recover the lost
discriminative power without giving up shift-invariance, we keep \emph{relative} phase.

\subsection{The bispectrum: a shift-invariant that remembers relative phase}

\begin{definition}[Triple correlation and bispectrum]
\label{def:bispec}
The \emph{triple correlation} is the autocorrelation with one extra shift,
$a_3(s_1,s_2) = \sum_t \overline{x_t}\, x_{t+s_1}\, x_{t+s_2}$. Its DFT is the
\emph{bispectrum}, which in frequency variables is the triple product
\[
  B_{k_1,k_2} \;=\; \widehat x_{k_1}\,\widehat x_{k_2}\,\overline{\widehat x_{k_1+k_2}}
\]
\citep{invariants_kakarala,wiki_bispectrum}. Chen--Zehni--Zhao write the equivalent form
\[
  B_{k_1,k_2} \;=\; \widehat x_{k_1}\,\overline{\widehat x_{k_2}}\,\widehat x_{k_2-k_1}
\]
\citep{invariants_chen}; the two differ only by a relabelling of the frequency axes. Setting
$k_2 = 0$ recovers the power spectrum: $B_{k,0} = \widehat x_k\,\widehat x_0\,
\overline{\widehat x_k} = \widehat x_0\,|\widehat x_k|^2$, so the bispectrum
\emph{contains} the power spectrum \citep{invariants_kakarala}.
\end{definition}

\begin{theorem}[Bispectrum is shift-invariant, by explicit phase cancellation]
\label{thm:bispecinv}
For every shift $\tau$, $\ B^{(R_\tau x)}_{k_1,k_2} = B^{(x)}_{k_1,k_2}$.
\end{theorem}

\begin{proof}
Apply the shift theorem (Lemma~\ref{lem:shifttheorem}) to each of the three Fourier
coefficients in the triple product. Write $\zeta = e^{-2\pi i \tau/N}$, so a shift
multiplies $\widehat x_k$ by $\zeta^{k}$, and \emph{conjugating} $\widehat x_{k_1+k_2}$
multiplies by $\overline{\zeta^{\,k_1+k_2}} = \zeta^{-(k_1+k_2)}$:
\begin{align*}
  B^{(R_\tau x)}_{k_1,k_2}
    \;&=\; \mathrm{DFT}(R_\tau x)_{k_1}\;
           \mathrm{DFT}(R_\tau x)_{k_2}\;
           \overline{\mathrm{DFT}(R_\tau x)_{k_1+k_2}} \\
    \;&=\; \bigl(\zeta^{k_1}\widehat x_{k_1}\bigr)\,
           \bigl(\zeta^{k_2}\widehat x_{k_2}\bigr)\,
           \overline{\bigl(\zeta^{k_1+k_2}\widehat x_{k_1+k_2}\bigr)}
        &&\text{(shift theorem on each factor)} \\
    \;&=\; \zeta^{k_1}\,\zeta^{k_2}\,\zeta^{-(k_1+k_2)}\;
           \widehat x_{k_1}\,\widehat x_{k_2}\,\overline{\widehat x_{k_1+k_2}}
        &&\text{(conjugation flips the phase)} \\
    \;&=\; \zeta^{\,k_1 + k_2 - (k_1 + k_2)}\;
           \widehat x_{k_1}\,\widehat x_{k_2}\,\overline{\widehat x_{k_1+k_2}} \\
    \;&=\; \zeta^{0}\; B^{(x)}_{k_1,k_2}
        \;=\; B^{(x)}_{k_1,k_2} .
\end{align*}
The three shift-phases cancel exactly: $k_1 + k_2 - (k_1+k_2) = 0$, so $\zeta^0 = 1$.
That cancellation is the whole point --- it is why the bispectrum is shift-invariant
\citep{invariants_kakarala,invariants_chen}. Unlike the power spectrum, though, the
bispectrum did \emph{not} have to take an absolute value to achieve invariance: it kept
the \emph{relative} phase among the three coefficients (the part of the phase that a
global shift cannot touch), and that retained phase is exactly the extra information
that lets it separate signals the power spectrum confuses.
\end{proof}

\subsection{Worked examples at {\boldmath$N=4$}}

\begin{example}[A shift leaves $|\mathrm{DFT}|^2$ unchanged]
\label{ex:psn4}
Take $x = (1,0,0,0)$. Then $\widehat x = F_4 x = (1,1,1,1)$, so $P = (1,1,1,1)$. Now
shift to $x' = (0,1,0,0)$: $\widehat{x'} = (1,-i,-1,i)$, and $|\widehat{x'}|^2 =
(1,1,1,1)$ again --- identical power spectrum, even though every phase changed.
\end{example}

\begin{example}[Discrete Wiener--Khinchin by hand]
\label{ex:wkn4}
Take $x = (2,1,0,1)$. Its circular autocorrelation $r_\tau = \sum_t x_t x_{t+\tau}$ is
$r = (6,4,2,4)$ (e.g.\ $r_0 = 4+1+0+1 = 6$). Its DFT is $\widehat r = F_4 r =
(16,4,0,4)$. Independently, $\widehat x = F_4 x = (4,2,0,2)$, so the power spectrum is
$|\widehat x|^2 = (16,4,0,4)$. The two agree: $\widehat r = |\widehat x|^2$, exactly
Theorem~\ref{thm:wk}.
\end{example}

\begin{example}[Same power spectrum, different bispectrum --- the disambiguation]
\label{ex:bispecn4}
Build two real length-$4$ signals
\[
  x_A = (1.25,\ 0.25,\ 0.25,\ 0.25), \qquad
  x_B = (0.75,\ -0.25,\ 0.75,\ 0.75).
\]
Their DFTs are $\widehat{x}_A = (2,1,1,1)$ and $\widehat{x}_B = (2,\,i,\,1,\,-i)$, so
both have the \emph{identical} magnitude spectrum $|\widehat x| = (2,1,1,1)$ and hence
the identical power spectrum $(4,1,1,1)$. The power spectrum cannot tell them apart. The
bispectrum can: with $B_{k_1,k_2} = \widehat x_{k_1}\widehat x_{k_2}
\overline{\widehat x_{k_1+k_2}}$,
\[
  B^{(A)}_{1,1} = \widehat{x}_{A,1}\,\widehat{x}_{A,1}\,\overline{\widehat{x}_{A,2}}
               = 1\cdot 1 \cdot 1 = +1,
  \qquad
  B^{(B)}_{1,1} = \widehat{x}_{B,1}\,\widehat{x}_{B,1}\,\overline{\widehat{x}_{B,2}}
               = i\cdot i \cdot 1 = -1 .
\]
The two differ ($+1$ versus $-1$), entirely because of the relative phase on the first
coefficient. And since both are shift-invariant, this distinction survives every shift:
shifting $x_B$ to $(0.75,0.75,-0.25,0.75)$ leaves $|\widehat x|^2 = (4,1,1,1)$ and
$B_{1,1} = -1$ unchanged. This is the discrete analogue of Kakarala's ``two shapes the
power spectrum confuses but the bispectrum separates'' \citep{invariants_kakarala}.
\end{example}

\begin{problem}
\label{prob:bispec}
For $x = (2,1,0,1)$: (a) verify the discrete Wiener--Khinchin identity $\widehat r =
|\widehat x|^2$ as in Example~\ref{ex:wkn4}; (b) shift $x$ by $2$ and confirm
$|\widehat x|^2$ is unchanged while the individual phases rotate; (c) construct a second
signal with the same $|\widehat x|$ but a flipped phase on one coefficient, and show
that some bispectrum entry $B_{k_1,k_2}$ differs --- so the bispectrum distinguishes
what the power spectrum cannot, even up to shift.
\end{problem}

\paragraph{Where this leaves us.} The four sections fit together as one thread. A
shift-invariant feature is what you get by averaging away the symmetry
(Section~\ref{sec:groups}); doing it \emph{linearly} keeps only the mean, so useful
invariants are nonlinear (Section~\ref{sec:invariants}); the natural nonlinear ones live
in the Fourier domain, where shifting is just a phase rotation (Section~\ref{sec:dft}),
and the whole averaging-is-projection story rests on the one projector theorem of
Section~\ref{sec:projectors}. The power spectrum is the phase-blind invariant; the
bispectrum is the phase-retaining one; both are exactly what the project measures when
it asks whether a network has learned to ignore position.

% =====================================================================
% \citep KEYS USED IN THIS FILE (for the assembler's bibliography):
%   Section 1 (projectors):
%     projectors_garrett   -- Garrett, UMN, "Projectors, projection maps,
%                             orthogonal projections" (central theorem proof)
%     projectors_tamu      -- TAMU MATH423 lecture 36 (numeric examples)
%     wiki_projection      -- Wikipedia, "Projection (linear algebra)"
%                             (oblique counterexample, eigenvalues in {0,1})
%     fourier_rutgers      -- Goodman, Rutgers Math357 (Pythagorean law,
%                             inner-product / Parseval machinery)
%   Section 2 (DFT / circulants):
%     fourier_bamieh       -- Bamieh, "Discovering Transforms" (arXiv:1805.05533)
%     fourier_mit1806      -- Johnson, MIT 18.06 circulant-matrices notes
%     fourier_rutgers      -- Goodman, Rutgers Math357 (diagonalization,
%                             convolution theorem, Parseval proofs)
%     wiki_convolution     -- Wikipedia, "Convolution theorem"
%     groups_bronstein     -- Bronstein-Bruna-Cohen-Velickovic 2021
%   Section 3 (groups / actions / Reynolds):
%     wiki_groupaction     -- Wikipedia, "Group action"
%     wiki_equivariant     -- Wikipedia, "Equivariant map"
%     groups_bronstein     -- Bronstein-Bruna-Cohen-Velickovic 2021 (GDL book)
%       (also attribute classically: Reynolds; Weyl/Hilbert invariant theory;
%        Cohen-Welling 2016; Zaheer et al. 2017 -- prose attributions only)
%   Section 4 (power spectrum / bispectrum):
%     invariants_mit6011        -- MIT 6.011 power spectral density notes
%     invariants_wienerkhinchin -- de Buyl, Wiener-Khinchin derivation
%     invariants_kakarala       -- Kakarala 2012, bispectral Fourier descriptors
%     invariants_chen           -- Chen, Zehni & Zhao 2018 (the file mislabeled
%                                  "bendory"); genuine Bendory et al. is its ref
%     wiki_bispectrum           -- Wikipedia, "Bispectrum"
%       (also attribute classically: Wiener; Khinchin; Einstein 1914;
%        Tukey/Brillinger polyspectra -- prose attributions only)
%
%  FULL \citep KEY LIST:
%    projectors_garrett, projectors_tamu, wiki_projection,
%    fourier_bamieh, fourier_mit1806, fourier_rutgers, wiki_convolution,
%    groups_bronstein, wiki_groupaction, wiki_equivariant,
%    invariants_mit6011, invariants_wienerkhinchin, invariants_kakarala,
%    invariants_chen, wiki_bispectrum
% =====================================================================

% =====================================================================
% Part 1 -- Learning-theory cluster for the shift-invariance / robustness primer.
% Three sections + closing bridge box:
%   1. Margins and the max-margin classifier (SVM), convex-hull view.
%   2. What gradient descent actually learns -- implicit bias.
%   3. NTK and the lazy-vs-rich dichotomy.
% This file is meant to be \input into a host document that already loads
% amsmath/amsthm and defines the theorem environments
%   theorem, lemma, corollary, proposition, definition, example, remark
% (exactly as in primer.tex). No \documentclass / \begin{document} here.
% Every concept is attributed with \citep{...}; the key list is at the end.
% =====================================================================

\section{Margins and the max-margin classifier}
\label{sec:margins}

Everything in this primer eventually comes back to a single geometric
quantity: \emph{how far is a data point from the surface where the
classifier changes its mind?} That distance is called the \emph{margin},
and a classifier that makes it as large as possible is a
\emph{max-margin classifier}, the simplest of which is the
(hard-margin) \emph{support vector machine}, or SVM
\citep{boser1992, cortes1995, vapnik}. We build the idea up slowly,
prove the one formula that every textbook quietly assumes, and then
look at the same object from a second angle (closest points of two
convex hulls) that will make the later robustness story intuitive.

\subsection{Hyperplanes and two notions of margin}

We work in $\mathbb{R}^d$. Our data are pairs $(\mathbf{x}_i, y_i)$ with
inputs $\mathbf{x}_i \in \mathbb{R}^d$ and binary labels
$y_i \in \{+1, -1\}$. A \emph{linear classifier} predicts the label of a
new point $\mathbf{x}$ by looking at the sign of an affine function
$\mathbf{w}\cdot\mathbf{x} + b$, where $\mathbf{w}\in\mathbb{R}^d$ is a
\emph{weight vector} and $b\in\mathbb{R}$ is a \emph{bias} (also called an
\emph{intercept} or \emph{threshold}). Here $\mathbf{w}\cdot\mathbf{x}
= \sum_{j=1}^d w_j x_j$ is the ordinary dot product. The prediction is
$+1$ when $\mathbf{w}\cdot\mathbf{x}+b>0$ and $-1$ when it is negative.

\begin{definition}[Hyperplane]
A \emph{hyperplane} in $\mathbb{R}^d$ is the set of points
$\{\mathbf{x} : \mathbf{w}\cdot\mathbf{x} + b = 0\}$ for some nonzero
$\mathbf{w}$ and scalar $b$. It is a flat set of dimension $d-1$: a line
in $\mathbb{R}^2$, an ordinary plane in $\mathbb{R}^3$. The vector
$\mathbf{w}$ is \emph{normal} (perpendicular) to the hyperplane, and the
hyperplane is the \emph{decision boundary} of the classifier: the place
where the score $\mathbf{w}\cdot\mathbf{x}+b$ is exactly zero, so the
predicted label flips.
\end{definition}

Why is $\mathbf{w}$ perpendicular to the hyperplane? Take any two points
$\mathbf{x}$ and $\mathbf{x}'$ that lie on it. Both satisfy the equation,
so $\mathbf{w}\cdot\mathbf{x}+b = 0$ and $\mathbf{w}\cdot\mathbf{x}'+b = 0$.
Subtracting, $\mathbf{w}\cdot(\mathbf{x}-\mathbf{x}') = 0$. The vector
$\mathbf{x}-\mathbf{x}'$ points along the hyperplane (it is the difference
of two of its points), and it is orthogonal to $\mathbf{w}$; since
$\mathbf{x},\mathbf{x}'$ were arbitrary, $\mathbf{w}$ is orthogonal to
every direction inside the hyperplane.

There are two different things one might call ``the margin,'' and keeping
them apart is the single most common source of confusion, so we name them
both \citep{burges1998}.

\begin{definition}[Functional and geometric margin]
For a labelled point $(\mathbf{x}_i, y_i)$ and a classifier
$(\mathbf{w}, b)$, the \emph{functional margin} is
\[
  \hat{\gamma}_i \;=\; y_i\,(\mathbf{w}\cdot\mathbf{x}_i + b).
\]
The \emph{geometric margin} is the functional margin divided by the length
of the weight vector,
\[
  \gamma_i \;=\; \frac{y_i\,(\mathbf{w}\cdot\mathbf{x}_i + b)}{\|\mathbf{w}\|},
\]
where $\|\mathbf{w}\| = \sqrt{\mathbf{w}\cdot\mathbf{w}}$ is the Euclidean
norm.
\end{definition}

The functional margin is positive exactly when the point is classified
correctly (the label $y_i$ and the score share a sign). But it is
\emph{not} a geometric distance: if we multiply both $\mathbf{w}$ and $b$
by $10$, the decision boundary does not move at all (the equation
$10\mathbf{w}\cdot\mathbf{x}+10b=0$ has the same solution set), yet the
functional margin grows tenfold. The geometric margin fixes this: it is
\emph{scale-invariant} (replacing $(\mathbf{w},b)$ by $(c\mathbf{w},cb)$
for $c>0$ leaves it unchanged because the $c$ in the numerator cancels the
$c$ in $\|c\mathbf{w}\|=c\|\mathbf{w}\|$), and as we now prove it equals the
honest perpendicular distance from $\mathbf{x}_i$ to the boundary.

\subsection{The point-to-hyperplane distance, with proof}

Burges' classic tutorial states the distance from a point to a hyperplane
as a fact and moves on \citep{burges1998}. Because this formula is the
backbone of margins, certificates and attacks alike, we supply the proof
in full.

\begin{theorem}[Point-to-hyperplane distance]
\label{thm:dist}
Let $H = \{\mathbf{x} : \mathbf{w}\cdot\mathbf{x} + b = 0\}$ with
$\mathbf{w}\neq \mathbf{0}$, and let $\mathbf{x}_0\in\mathbb{R}^d$ be any
point. The Euclidean distance from $\mathbf{x}_0$ to $H$ is
\[
  \operatorname{dist}(\mathbf{x}_0, H)
  \;=\; \frac{|\mathbf{w}\cdot\mathbf{x}_0 + b|}{\|\mathbf{w}\|}.
\]
\end{theorem}

\begin{proof}
The distance from $\mathbf{x}_0$ to the set $H$ is by definition the
smallest value of $\|\mathbf{x}_0 - \mathbf{x}\|$ over all
$\mathbf{x}\in H$. We will (i) guess the closest point by walking from
$\mathbf{x}_0$ straight toward $H$ along the normal direction, and
(ii) prove no point of $H$ is nearer.

\emph{Step 1 (a unit normal).}
Set $\mathbf{n} = \mathbf{w}/\|\mathbf{w}\|$, the normal $\mathbf{w}$
rescaled to length one, so $\mathbf{n}\cdot\mathbf{n} = 1$.

\emph{Step 2 (walk along the normal).}
Consider the point obtained by moving a signed amount $t\in\mathbb{R}$ from
$\mathbf{x}_0$ along $\mathbf{n}$,
\[
  \mathbf{p}(t) \;=\; \mathbf{x}_0 - t\,\mathbf{n}.
\]
We want the $t$ that lands $\mathbf{p}(t)$ on $H$, i.e.\ that makes
$\mathbf{w}\cdot\mathbf{p}(t)+b = 0$.

\emph{Step 3 (solve for $t$).}
Substitute and expand, using linearity of the dot product in the variable
$\mathbf{p}$,
\begin{align*}
  \mathbf{w}\cdot\mathbf{p}(t) + b
  &= \mathbf{w}\cdot(\mathbf{x}_0 - t\,\mathbf{n}) + b \\
  &= \mathbf{w}\cdot\mathbf{x}_0 \;-\; t\,(\mathbf{w}\cdot\mathbf{n}) \;+\; b.
\end{align*}
Now $\mathbf{w}\cdot\mathbf{n} = \mathbf{w}\cdot(\mathbf{w}/\|\mathbf{w}\|)
= \|\mathbf{w}\|^2/\|\mathbf{w}\| = \|\mathbf{w}\|$. Setting the expression
to zero and solving for $t$,
\[
  \mathbf{w}\cdot\mathbf{x}_0 - t\,\|\mathbf{w}\| + b = 0
  \quad\Longrightarrow\quad
  t^\star \;=\; \frac{\mathbf{w}\cdot\mathbf{x}_0 + b}{\|\mathbf{w}\|}.
\]
So the foot of the perpendicular is $\mathbf{p}(t^\star)\in H$, and the
length of the step we took is
\[
  \|\mathbf{x}_0 - \mathbf{p}(t^\star)\| = \|t^\star\,\mathbf{n}\|
  = |t^\star|\,\|\mathbf{n}\| = |t^\star|
  = \frac{|\mathbf{w}\cdot\mathbf{x}_0 + b|}{\|\mathbf{w}\|}.
\]

\emph{Step 4 (no point of $H$ is nearer).}
Let $\mathbf{x}\in H$ be arbitrary, so $\mathbf{w}\cdot\mathbf{x}+b = 0$.
Write the gap from $\mathbf{x}_0$ as
$\mathbf{x}_0 - \mathbf{x} = \big(\mathbf{x}_0-\mathbf{p}(t^\star)\big)
+ \big(\mathbf{p}(t^\star)-\mathbf{x}\big)$. The first piece equals
$t^\star\mathbf{n}$ (a multiple of the normal). The second piece,
$\mathbf{p}(t^\star)-\mathbf{x}$, is the difference of two points of $H$,
hence (as shown above) orthogonal to $\mathbf{w}$ and so to $\mathbf{n}$.
The two pieces are therefore orthogonal, and Pythagoras gives
\[
  \|\mathbf{x}_0 - \mathbf{x}\|^2
  = \|t^\star\mathbf{n}\|^2 + \|\mathbf{p}(t^\star)-\mathbf{x}\|^2
  \;\ge\; \|t^\star\mathbf{n}\|^2 = |t^\star|^2,
\]
with equality iff $\mathbf{p}(t^\star)=\mathbf{x}$. Hence the minimum
distance is exactly $|t^\star| = |\mathbf{w}\cdot\mathbf{x}_0+b|/\|\mathbf{w}\|$.
\end{proof}

Two corollaries we will reuse constantly. First, the distance from the
hyperplane to the \emph{origin} (set $\mathbf{x}_0 = \mathbf{0}$) is
$|b|/\|\mathbf{w}\|$. Second, comparing with the geometric margin:
$\gamma_i = y_i(\mathbf{w}\cdot\mathbf{x}_i+b)/\|\mathbf{w}\|$ is precisely
the signed distance of $\mathbf{x}_i$ to the boundary, positive when the
point is on its correct side. \emph{The geometric margin is a distance.}

\subsection{Canonical scaling, support vectors, and the hard-margin SVM}

Because the geometric margin does not change under rescaling, we are free
to fix the scale of $(\mathbf{w}, b)$ however is most convenient. The
convention everyone uses is the \emph{canonical} one
\citep{burges1998}: rescale so that the closest correctly classified
points sit at functional margin exactly one. Concretely we require

\begin{equation}
\label{eq:canonical}
  y_i\,(\mathbf{w}\cdot\mathbf{x}_i + b) \;\ge\; 1 \qquad\text{for all }i,
\end{equation}
with equality achieved by the nearest points. Under this convention the
two ``rails'' $\mathbf{w}\cdot\mathbf{x}+b = +1$ and
$\mathbf{w}\cdot\mathbf{x}+b = -1$ are parallel to the boundary and sit
on either side of it. By Theorem~\ref{thm:dist} each rail is at distance
$1/\|\mathbf{w}\|$ from the boundary, so the total width of the
separating slab is

\[
  \text{margin} \;=\; d_+ + d_- \;=\; \frac{1}{\|\mathbf{w}\|} + \frac{1}{\|\mathbf{w}\|}
  \;=\; \frac{2}{\|\mathbf{w}\|},
\]

where $d_+$ ($d_-$) is the distance to the nearest positive (negative)
point \citep{burges1998}.

\begin{definition}[Support vectors]
The \emph{support vectors} are the training points that actually touch a
rail, i.e.\ for which equality holds in \eqref{eq:canonical}:
$y_i(\mathbf{w}\cdot\mathbf{x}_i+b)=1$. They are the points lying exactly
on the slab boundary; the solution is determined entirely by them, and
moving or deleting a non-support point (one strictly inside its own half)
does not change the classifier.
\end{definition}

We want the widest slab, i.e.\ the largest $2/\|\mathbf{w}\|$. Since
$2/\|\mathbf{w}\|$ is a strictly decreasing function of $\|\mathbf{w}\|$,
maximizing the margin is the same as minimizing $\|\mathbf{w}\|$, and
equivalently $\tfrac12\|\mathbf{w}\|^2$ (a smooth, convex stand-in). This
is the \emph{hard-margin SVM} \citep{boser1992, vapnik, burges1998}:

\begin{equation}
\label{eq:svm-primal}
  \min_{\mathbf{w}, b}\ \tfrac{1}{2}\|\mathbf{w}\|^2
  \qquad\text{subject to}\qquad
  y_i\,(\mathbf{w}\cdot\mathbf{x}_i + b) \ge 1 \ \ \forall i.
\end{equation}

\begin{proposition}[Max-margin $\Longleftrightarrow$ minimum norm]
\label{prop:maxmargin}
Among all $(\mathbf{w}, b)$ satisfying the canonical constraints
\eqref{eq:canonical}, the one with the largest geometric margin is exactly
the minimizer of $\tfrac12\|\mathbf{w}\|^2$ in \eqref{eq:svm-primal}.
\end{proposition}
\begin{proof}
The margin under canonical scaling is $2/\|\mathbf{w}\|$. The map
$r \mapsto 2/r$ is strictly decreasing on $r>0$, so its maximum over the
feasible set is attained precisely where $\|\mathbf{w}\|$ (hence
$\tfrac12\|\mathbf{w}\|^2$) is smallest. The two optimizations have the
same minimizer.
\end{proof}

Problem \eqref{eq:svm-primal} is a convex quadratic program with linear
constraints, so it has a unique global solution $\mathbf{w}$; the
\emph{Karush--Kuhn--Tucker} (KKT) conditions characterize it. We pause to
define KKT cleanly because it returns in Sections~\ref{sec:implicit} and is
the language all of these results speak.

\begin{definition}[KKT point]
\label{def:kkt}
Consider $\min_{\mathbf{z}} f(\mathbf{z})$ subject to inequality
constraints $g_i(\mathbf{z})\ge 0$. A point $\mathbf{z}^\star$ is a
\emph{KKT point} if there exist multipliers $\alpha_i \ge 0$ (one per
constraint) such that
\begin{enumerate}[label=(\roman*),nosep]
\item \emph{stationarity:}
  $\nabla f(\mathbf{z}^\star) = \sum_i \alpha_i \nabla g_i(\mathbf{z}^\star)$;
\item \emph{primal feasibility:} $g_i(\mathbf{z}^\star)\ge 0$ for all $i$;
\item \emph{dual feasibility:} $\alpha_i \ge 0$ for all $i$;
\item \emph{complementary slackness:} $\alpha_i\, g_i(\mathbf{z}^\star) = 0$
  for all $i$.
\end{enumerate}
For a convex problem the KKT conditions are necessary \emph{and}
sufficient for a global minimum; for a non-convex problem they are only
necessary (a KKT point need not be a global minimizer). Intuitively, a
multiplier $\alpha_i$ is the ``price'' of constraint $i$, and
complementary slackness says we only pay for constraints that are tight
($g_i = 0$).
\end{definition}

Applying Definition~\ref{def:kkt} to the SVM \eqref{eq:svm-primal} with
$g_i(\mathbf{w},b) = y_i(\mathbf{w}\cdot\mathbf{x}_i+b)-1$, complementary
slackness reads
\[
  \alpha_i\,\big[\,y_i(\mathbf{w}\cdot\mathbf{x}_i+b)-1\,\big] = 0.
\]
So $\alpha_i>0$ forces $y_i(\mathbf{w}\cdot\mathbf{x}_i+b)=1$: the points
with nonzero multiplier are exactly the support vectors. Stationarity in
$\mathbf{w}$ gives $\mathbf{w} = \sum_i \alpha_i y_i \mathbf{x}_i$, so the
optimal weight vector is a nonnegative combination of the support vectors
alone. \emph{Remember this: the SVM solution is built only from the points
that sit on the margin.} Section~\ref{sec:implicit} will show plain
gradient descent rediscovers exactly this same combination on its own.

\subsection{The convex-hull view: closest points of two hulls}

There is a second picture, due to Bennett and Bredensteiner
\citep{bennett2000}, that many people find more intuitive than the
algebra and which previews the robustness story. It needs one definition.

\begin{definition}[Convex hull]
The \emph{convex hull} $\operatorname{conv}(A)$ of a finite set
$A = \{\mathbf{a}_1, \dots, \mathbf{a}_m\}$ is the set of all
\emph{convex combinations} of its points,
\[
  \operatorname{conv}(A) = \Big\{\ \textstyle\sum_{i} u_i \mathbf{a}_i \ :\
  u_i \ge 0,\ \sum_i u_i = 1 \ \Big\}.
\]
It is the smallest convex set containing $A$: in the plane, the rubber
band stretched around the points.
\end{definition}

Suppose the two classes are linearly separable. Let $A$ be the positive
points and $B$ the negative points. Bennett's theorem says the max-margin
hyperplane is determined by the two closest points of the two hulls.

\begin{theorem}[Closest points of the hulls \citep{bennett2000}]
\label{thm:hulls}
For separable classes, let
\[
  (\mathbf{c}^\star, \mathbf{d}^\star) = \arg\min_{\mathbf{c}\in\operatorname{conv}(A),\,
    \mathbf{d}\in\operatorname{conv}(B)} \ \tfrac{1}{2}\|\mathbf{c}-\mathbf{d}\|^2
\]
be the closest pair of points, one in each hull. Then the max-margin
hyperplane is the perpendicular bisector of the segment
$[\mathbf{c}^\star, \mathbf{d}^\star]$, its normal is
$\mathbf{w}\propto \mathbf{c}^\star - \mathbf{d}^\star$, and the margin
(slab width) equals the distance between the hulls,
$\|\mathbf{c}^\star - \mathbf{d}^\star\|$.
\end{theorem}

\begin{proof}[Proof skeleton]
The argument is by contradiction on optimality
\citep{bennett2000}. (a) \emph{The normal must be along
$\mathbf{c}^\star-\mathbf{d}^\star$.} If the optimal separating slab's
normal were \emph{not} parallel to the segment joining the closest hull
points, then the two parallel supporting planes (one touching each hull)
would not be as far apart as possible: tilting the slab to align with the
segment would increase the gap, contradicting maximality. (b) \emph{The
segment is orthogonal to the supporting planes.} If it were not, then the
foot of one endpoint could be slid along its hull to a strictly nearer
point, contradicting that $\mathbf{c}^\star,\mathbf{d}^\star$ are the
closest pair. Conditions (a) and (b) together force the boundary to be the
perpendicular bisector of $[\mathbf{c}^\star,\mathbf{d}^\star]$ with normal
$\mathbf{c}^\star-\mathbf{d}^\star$; the slab width is then the length of
that segment. The full statement is that the ``closest points of the
hulls'' problem is the Wolfe dual of the ``push two parallel supporting
planes apart'' problem, so they have the same solution
\citep{bennett2000}.
\end{proof}

The geometric reading is worth keeping: \emph{maximizing the margin is the
same as finding the two nearest points of the two clouds and bisecting the
segment between them.} (When classes overlap, Bennett shrinks each hull by
capping every weight $u_i\le 1/K$ -- the \emph{reduced convex hull} -- and
shows this softening of the hulls is the dual of allowing margin
violations, the soft-margin SVM of \citet{cortes1995}. We will not need
the soft-margin machinery, but it is good to know the picture extends.)

\subsection{Two worked examples and an exercise}

\begin{example}[Two points, margin $2$]
\label{ex:twopoint}
Take a single positive point $\mathbf{x}_+ = (2,0)$ with $y=+1$ and a
single negative point $\mathbf{x}_- = (0,0)$ with $y=-1$ in $\mathbb{R}^2$.
By symmetry the boundary is the vertical line halfway between them,
$x_1 = 1$. Let us verify both ways.

\emph{Canonical algebra.} The boundary is normal to the $x_1$-axis, so
$\mathbf{w}=(w_1,0)$. The canonical constraints with equality at both
points read $w_1\cdot 2 + b = +1$ and $-(w_1\cdot 0 + b) = +1$, i.e.\
$b=-1$ and $2w_1 - 1 = 1$, giving $w_1 = 1$. Then $\|\mathbf{w}\| = 1$ and
the margin is $2/\|\mathbf{w}\| = 2$. The boundary
$\mathbf{w}\cdot\mathbf{x}+b = x_1 - 1 = 0$ is indeed $x_1=1$.

\emph{Convex-hull view.} Each ``hull'' is a single point:
$\mathbf{c}^\star=(2,0)$, $\mathbf{d}^\star=(0,0)$. The normal is
$\mathbf{w}\propto\mathbf{c}^\star-\mathbf{d}^\star = (2,0)$, the bisector
is $x_1=1$, and the margin is
$\|\mathbf{c}^\star-\mathbf{d}^\star\| = 2$. The two pictures agree.
\end{example}

\begin{example}[Four facing segments, margin $3$]
\label{ex:foursegment}
Let $A = \{(3,1),(3,-1)\}$ ($y=+1$) and $B = \{(0,1),(0,-1)\}$ ($y=-1$).
Each hull is now a vertical segment: $\operatorname{conv}(A)$ is the
segment from $(3,-1)$ to $(3,1)$, and $\operatorname{conv}(B)$ the segment
from $(0,-1)$ to $(0,1)$. The two segments face each other; the closest
distance between them is the horizontal gap, $\|(3,y)-(0,y)\| = 3$, and it
is achieved by \emph{every} horizontally aligned pair (e.g.\ $(3,0)$ and
$(0,0)$, or $(3,1)$ and $(0,1)$). So the margin is $3$. Canonically, the
boundary is $x_1 = 1.5$ with $\mathbf{w}=(w_1,0)$; equality at $x_1=3$
gives $3w_1 + b = 1$ and at $x_1 = 0$ gives $-b = 1$, so $b = -1$ and
$w_1 = 2/3$, hence margin $2/\|\mathbf{w}\| = 2/(2/3) = 3$, matching the
hull distance. Note all four points are support vectors: the closest pair
is non-unique even though the optimal plane is unique.
\end{example}

\begin{problem}[Point versus segment]
\label{ex:pointsegment}
Let $A=\{(2,2)\}$ ($y=+1$) and $B=\{(0,0),(0,2)\}$ ($y=-1$). Sketch
$\operatorname{conv}(B)$, the vertical segment from $(0,0)$ to $(0,2)$.
Project the single positive point $(2,2)$ onto that segment to find the
closest hull point $\mathbf{d}^\star$. Deduce the SVM normal
$\mathbf{w}\propto\mathbf{c}^\star-\mathbf{d}^\star$ and the margin
$\|\mathbf{c}^\star-\mathbf{d}^\star\|$, then solve for the canonical
$(\mathbf{w},b)$ and identify which points are support vectors. (Hint:
the foot of the perpendicular from $(2,2)$ to the segment $x_1=0,\,
0\le x_2\le 2$ is $\mathbf{d}^\star = (0,2)$, so $\mathbf{w}\propto(2,0)$,
margin $2$; $(2,2)$ and $(0,2)$ are support vectors, $(0,0)$ is not.)
\end{problem}

% =====================================================================

\section{What gradient descent actually learns: implicit bias}
\label{sec:implicit}

We have a max-margin classifier as an \emph{optimization problem}
\eqref{eq:svm-primal}. But that is not how neural networks (or even
logistic regression) are trained in practice. In practice we pick a loss
function, write down the average loss over the training set, and run
gradient descent, with \emph{no margin or norm term anywhere in the
objective}. So here is a genuinely surprising fact, and the bridge from
Section~\ref{sec:margins} to the rest of the primer: even though nothing
in the objective mentions the margin, gradient descent on the logistic
loss drifts toward the max-margin solution all by itself. This silent
preference is called \emph{implicit bias} \citep{soudry2018}.

\subsection{The setup and the meaning of ``convergence in direction''}

Fold the labels into the inputs by replacing $\mathbf{x}_n$ with
$y_n\mathbf{x}_n$, so we may assume every label is $+1$ and write the
linear predictor as $\mathbf{w}^\top\mathbf{x}_n$. The data are
\emph{linearly separable} if there is some $\mathbf{w}_*$ with
$\mathbf{w}_*^\top\mathbf{x}_n > 0$ for all $n$ -- a weight vector that
gets every point right \citep{soudry2018}. We minimize the empirical loss
\[
  \mathcal{L}(\mathbf{w}) = \sum_{n=1}^N \ell\!\big(\mathbf{w}^\top\mathbf{x}_n\big),
\]
where $\ell$ is, for instance, the logistic loss
$\ell(u) = \log(1+e^{-u})$ or the exponential loss $\ell(u)=e^{-u}$. Both
are positive, smooth, strictly decreasing, and flatten toward $0$ as
$u\to\infty$. Gradient descent takes steps
\[
  \mathbf{w}(t+1) = \mathbf{w}(t) - \eta\,\nabla\mathcal{L}(\mathbf{w}(t)),
  \qquad
  \nabla\mathcal{L}(\mathbf{w}) = \sum_{n=1}^N \ell'\!\big(\mathbf{w}^\top\mathbf{x}_n\big)\mathbf{x}_n,
\]
with a fixed learning rate $\eta>0$ (a small step-size multiplier).

Here is the catch that makes the question subtle. On separable data, to
push the loss to its infimum $0$ the predictor must send every
$\mathbf{w}^\top\mathbf{x}_n\to+\infty$, which forces
$\|\mathbf{w}(t)\|\to\infty$. There is no finite minimizer to converge to.
But for \emph{classification} only the \emph{direction} of $\mathbf{w}$
matters, since the predicted label is $\operatorname{sign}(\mathbf{w}^\top
\mathbf{x})$ and rescaling $\mathbf{w}$ does not change a sign. So the
right question is about the direction.

\begin{definition}[Convergence in direction]
We say $\mathbf{w}(t)$ \emph{converges in direction} to a unit vector
$\mathbf{u}$ if $\|\mathbf{w}(t)\|\to\infty$ while the normalized iterate
settles, $\mathbf{w}(t)/\|\mathbf{w}(t)\| \to \mathbf{u}$.
\end{definition}

\subsection{The linear separable case, proved}

The centerpiece of this section is the theorem of
\citet{soudry2018}: the direction that gradient descent locks onto is the
$L_2$ max-margin (hard-margin SVM) direction. We state it, then give the
two-part argument it rests on -- a convergence lemma and the
``support-vector gradient-domination'' idea that explains \emph{why} the
SVM direction emerges.

\begin{theorem}[Implicit bias of gradient descent, linear case
\citep{soudry2018}]
\label{thm:soudry}
Let the data be linearly separable and $\ell$ be a smooth, decreasing loss
with an exponential tail (the logistic and exponential losses qualify).
For any small enough learning rate $\eta$ and any starting point,
\[
  \mathbf{w}(t) \;=\; \widehat{\mathbf{w}}\,\log t \;+\; \boldsymbol{\rho}(t),
  \qquad \|\boldsymbol{\rho}(t)\| = O(\log\log t),
\]
where $\widehat{\mathbf{w}}$ is the $L_2$ max-margin vector, the solution
of the hard-margin SVM
$\widehat{\mathbf{w}} = \arg\min \|\mathbf{w}\|^2$ s.t.\
$\mathbf{w}^\top\mathbf{x}_n\ge 1$. Consequently
$\mathbf{w}(t)/\|\mathbf{w}(t)\| \to \widehat{\mathbf{w}}/\|\widehat{\mathbf{w}}\|$.
\end{theorem}

We first record that gradient descent does reach zero loss, with the
norm diverging. This is the ``no finite critical point'' lemma of
\citet{soudry2018}.

\begin{lemma}[The loss goes to zero, the norm diverges \citep{soudry2018}]
\label{lem:diverge}
Under the assumptions of Theorem~\ref{thm:soudry},
$\mathcal{L}(\mathbf{w}(t))\to 0$, $\|\mathbf{w}(t)\|\to\infty$, and
$\mathbf{w}(t)^\top\mathbf{x}_n\to+\infty$ for every $n$.
\end{lemma}
\begin{proof}
Take any separator $\mathbf{w}_*$, so $\mathbf{w}_*^\top\mathbf{x}_n>0$ for
all $n$. Dot the gradient with $\mathbf{w}_*$:
\[
  \mathbf{w}_*^\top \nabla\mathcal{L}(\mathbf{w})
  = \sum_{n=1}^N \ell'\!\big(\mathbf{w}^\top\mathbf{x}_n\big)\,
    \big(\mathbf{w}_*^\top\mathbf{x}_n\big).
\]
Every term is a product of $\ell'(\cdot)<0$ (the loss is strictly
decreasing) and $\mathbf{w}_*^\top\mathbf{x}_n>0$, so every term is
\emph{strictly negative}. Hence at any finite $\mathbf{w}$ the sum is
strictly negative and cannot be $0$, which means \emph{there is no finite
critical point}: $\nabla\mathcal{L}(\mathbf{w})\neq\mathbf{0}$ everywhere.
Gradient descent on a smooth loss with a small enough step size is
guaranteed to drive $\nabla\mathcal{L}(\mathbf{w}(t))\to\mathbf{0}$; since
that can only happen ``at infinity,'' the iterate norm diverges,
$\|\mathbf{w}(t)\|\to\infty$. Because $\ell'(u)\to 0$ only as
$u\to+\infty$, the gradient can vanish only if every margin
$\mathbf{w}(t)^\top\mathbf{x}_n\to+\infty$; this forces
$\ell(\mathbf{w}(t)^\top\mathbf{x}_n)\to 0$ for each $n$, so
$\mathcal{L}(\mathbf{w}(t))\to 0$.
\end{proof}

Now the heart of the matter: why is the limiting \emph{direction} the SVM
direction? The argument is what \citet{soudry2018} call gradient
domination by the support vectors.

\begin{proof}[Why the direction is the max-margin direction (sketch of
Theorem~\ref{thm:soudry})]
Use the exponential loss, $\ell(u)=e^{-u}$, for clarity (the logistic
loss behaves the same way to leading order). Suppose, as
Lemma~\ref{lem:diverge}
permits, that $\mathbf{w}(t)/\|\mathbf{w}(t)\|$ approaches some unit
direction, and write the iterate as $\mathbf{w}(t) = g(t)\,\mathbf{w}_\infty
+ \boldsymbol{\rho}(t)$ with $g(t)\to\infty$ along the limit direction
$\mathbf{w}_\infty$ and $\boldsymbol\rho(t)/g(t)\to 0$. The negative
gradient is

\begin{align*}
  -\nabla\mathcal{L}(\mathbf{w}(t))
  &= \sum_{n=1}^N \exp\!\big(-\mathbf{w}(t)^\top\mathbf{x}_n\big)\,\mathbf{x}_n \\
  &= \sum_{n=1}^N \exp\!\big(-g(t)\,\mathbf{w}_\infty^\top\mathbf{x}_n\big)\,
     \exp\!\big(-\boldsymbol\rho(t)^\top\mathbf{x}_n\big)\,\mathbf{x}_n.
\end{align*}

As $g(t)\to\infty$ the factor $\exp(-g(t)\,\mathbf{w}_\infty^\top
\mathbf{x}_n)$ decays geometrically, and it decays \emph{most slowly} for
the points with the \emph{smallest} value of $\mathbf{w}_\infty^\top
\mathbf{x}_n$ -- precisely the points closest to the boundary, the support
vectors $\arg\min_n \mathbf{w}_\infty^\top\mathbf{x}_n$. All other terms
are exponentially smaller and become negligible. So in the limit the
update direction is dominated by, and is a \emph{nonnegative combination
of}, the support vectors:
\[
  \mathbf{w}_\infty \;\propto\; \widehat{\mathbf{w}}
  \;=\; \sum_{n} \alpha_n \mathbf{x}_n,\qquad
  \alpha_n\ge 0, \ \text{ with } \alpha_n>0 \text{ only for support vectors.}
\]
These are exactly the KKT conditions (Definition~\ref{def:kkt}) of the
hard-margin SVM \eqref{eq:svm-primal}: a nonnegative support-vector
combination, with $\widehat{\mathbf{w}}^\top\mathbf{x}_n=1$ on the
support vectors and $>1$ off them. Hence $\mathbf{w}_\infty$ is
proportional to the SVM solution $\widehat{\mathbf{w}}$. The detailed
proof of \citet{soudry2018} additionally shows $g(t)=\log t$ and bounds
the residual $\boldsymbol\rho(t)$, giving the stated rate.
\end{proof}

Three things to take away, all of which appear later. (1) The convergence
to the SVM direction is \emph{very slow}: the direction error decays like
$O(1/\log t)$, whereas the loss itself decays like $O(1/t)$. So the loss
can be tiny while the margin is still quietly improving -- a reason to keep
training past zero training error, and a warning that loss-based early
stopping can cut off margin (hence robustness) gains \citep{soudry2018}.
(2) The bias is \emph{specific to gradient descent and to exponential-tail
losses}: adaptive methods such as Adam, or polynomial-tail losses, need
not converge to the $L_2$ max-margin direction \citep{soudry2018}. (3) The
same geometric object -- the support vectors, the convex-hull closest
points of Section~\ref{sec:margins} -- is what gradient descent homes in
on.

\subsection{Homogeneous networks: the normalized margin and KKT points}

Linear predictors are a warm-up. Do deep networks, trained the same way,
also drift toward a margin-maximizing solution? \citet{lyuli2020} answer
this for an important class of networks, and we must define two terms
carefully before stating their result.

\begin{definition}[$L$-homogeneous network]
\label{def:homog}
A network $\Phi(\boldsymbol\theta;\mathbf{x})$ with parameters
$\boldsymbol\theta$ is \emph{(positively) $L$-homogeneous} if scaling all
the parameters by a positive constant $c$ scales the output by $c^L$:
\[
  \Phi(c\,\boldsymbol\theta;\mathbf{x}) = c^{L}\,\Phi(\boldsymbol\theta;\mathbf{x})
  \qquad\text{for all } c>0.
\]
Fully-connected and convolutional networks with ReLU (or LeakyReLU)
activations and \emph{no bias terms} are homogeneous, with $L$ equal to
the depth (the number of layers) \citep{lyuli2020}. The reason: ReLU
satisfies $\mathrm{ReLU}(c\,z) = c\,\mathrm{ReLU}(z)$ for $c>0$, and each
linear layer multiplies the output by its weight, so scaling the weights
of all $L$ layers multiplies the final output by $c^L$.
\end{definition}

Just as in the linear case, only the direction of $\boldsymbol\theta$
matters for the predicted class, and by homogeneity the raw margin scales
with $\|\boldsymbol\theta\|^L$. So to compare directions fairly we
normalize.

\begin{definition}[Normalized margin \citep{lyuli2020}]
Let $q_n(\boldsymbol\theta) = y_n\Phi(\boldsymbol\theta;\mathbf{x}_n)$ be
the margin on point $n$, and $q_{\min}(\boldsymbol\theta)
= \min_n q_n(\boldsymbol\theta)$ the smallest margin over the data. The
\emph{normalized margin} divides out the scale,
\[
  \bar\gamma(\boldsymbol\theta)
  = \frac{q_{\min}(\boldsymbol\theta)}{\|\boldsymbol\theta\|^{L}}.
\]
It is invariant to rescaling $\boldsymbol\theta$, just as the geometric
margin was for the linear model.
\end{definition}

\begin{theorem}[Gradient descent maximizes the normalized margin
\citep{lyuli2020}]
\label{thm:lyuli}
Consider an $L$-homogeneous network with an exponential-tail loss
(logistic, exponential, cross-entropy), trained by gradient descent /
gradient flow, and suppose at some time the network reaches $100\%$
training accuracy. Then:
\begin{enumerate}[label=(\alph*),nosep]
\item a smoothed version of the normalized margin $\bar\gamma$ is
  \emph{monotonically non-decreasing} thereafter; and
\item every limit point of the normalized direction
  $\boldsymbol\theta(t)/\|\boldsymbol\theta(t)\|$ lies along a
  \emph{KKT point} of the margin-maximization program
  \[
    \min \ \tfrac12\|\boldsymbol\theta\|^2
    \quad\text{s.t.}\quad q_n(\boldsymbol\theta)\ge 1 \ \ \forall n.
  \]
\end{enumerate}
\end{theorem}

We state Theorem~\ref{thm:lyuli} \emph{without proof}: the argument uses
the Clarke subdifferential (a generalized gradient for the non-smooth ReLU
network), an Euler-type radial/tangential decomposition of the parameter
flow, and a monotonicity argument for the smoothed margin -- machinery
well beyond a first course, and not needed to use the result
\citep{lyuli2020}. \citet{jitelgarsky2019} prove a closely related
statement and, for general (possibly non-separable) data, show gradient
descent follows a unique ray whose direction is the max-margin separator
of the separable part of the data.

\paragraph{An honest caveat: which KKT point?}
The crucial difference from the linear case is the word ``KKT.'' For the
convex SVM, a KKT point \emph{is} the global max-margin solution
(Definition~\ref{def:kkt}). For a non-convex deep network the
margin-maximization program is non-convex, so a KKT point is only a
\emph{first-order stationary point}; it need not be the global
maximum-margin solution, and \citet{lyuli2020} note one can construct
examples where it is not. Theorem~\ref{thm:lyuli} therefore tells us
gradient descent maximizes margin \emph{locally / in a stationarity
sense}, but it does \emph{not} pin down which margin-stationary solution
among possibly many. This honesty matters for the project: ``gradient
descent maximizes the margin'' is exactly true for linear models and only
true ``up to KKT'' for deep nets.

\paragraph{The margin/robustness bridge.}
Why care about margin at all? Because the (normalized) margin
\emph{lower-bounds the distance to the decision boundary}, and that
distance is the robust radius. Concretely \citet{lyuli2020} state that for
a ReLU network the normalized margin on a point, divided by the network's
Lipschitz constant, is a lower bound on the $L_2$ distance from that point
to the boundary (the same Lipschitz--margin certificate that the
robustness section of this primer proves via \citet{tsuzuku2018lipschitz}).
So: plain gradient descent grows the margin, a larger margin means a larger
certified robust radius, and that is the thread connecting ``what training
does'' to ``how robust the model is.''

\begin{example}[2D logistic gradient descent finds the SVM direction]
\label{ex:gd2d}
Take two positive points $\mathbf{x}_1=(1,1)$ and $\mathbf{x}_2=(1,-1)$
(both label $+1$, so no label folding needed) and the exponential loss.
The hard-margin SVM solution is $\widehat{\mathbf{w}}=(1,0)$: it is the
shortest $\mathbf{w}$ with $\mathbf{w}^\top\mathbf{x}_1\ge1$ and
$\mathbf{w}^\top\mathbf{x}_2\ge1$, and by symmetry it points along the
$x_1$-axis. Now run gradient flow from $\mathbf{w}(0)=(0,c)$. The negative
gradient is
\[
  -\nabla\mathcal{L}(\mathbf{w})
  = e^{-(w_1+w_2)}(1,1) + e^{-(w_1-w_2)}(1,-1),
\]
whose second component is $e^{-(w_1+w_2)} - e^{-(w_1-w_2)}$. When $w_2>0$
this is negative (the first exponent is larger, so its term is smaller),
pushing $w_2$ down; when $w_2<0$ it is positive, pushing $w_2$ up. Either
way $w_2\to0$, while the first component is always positive so $w_1\to
+\infty$. Hence $\mathbf{w}(t)/\|\mathbf{w}(t)\|\to(1,0)$: gradient descent
recovers the SVM direction, exactly as Theorem~\ref{thm:soudry} promises.
\end{example}

\begin{problem}[Non-support points do not move the limit]
\label{ex:nonsupport}
Add a third positive point $\mathbf{x}_3=(3,0)$ to
Example~\ref{ex:gd2d}. Check that $\widehat{\mathbf{w}}^\top\mathbf{x}_3
= 3 > 1$, so $\mathbf{x}_3$ is \emph{not} a support vector and its KKT
multiplier is $\alpha_3=0$. Verify the SVM stationarity
$\widehat{\mathbf{w}} = \alpha_1\mathbf{x}_1 + \alpha_2\mathbf{x}_2$, and
argue that although $\mathbf{x}_3$ tugs on the early gradient, the limiting
direction is unchanged and would be the same if $\mathbf{x}_3$ were
deleted. (This is gradient domination by support vectors in action.)
\end{problem}

% =====================================================================

\section{The neural tangent kernel and the lazy-versus-rich dichotomy}
\label{sec:ntk}

The implicit-bias results tell us a network's training drifts toward a
margin-maximizing direction. But \emph{what features} does it use to get
there? Does it learn new features tuned to the task, or does it ride a
fixed set of features baked in at the start? This is the lazy-versus-rich
question, and the tool for making it precise is the \emph{neural tangent
kernel} \citep{jacot2018}. We first build the kernel from a Taylor
expansion, state Jacot's two theorems, give Chizat's criterion for when
training is ``lazy,'' and close with the one definition (feature
averaging) that links all of this back to robustness.

\subsection{Linearization and the neural tangent kernel}

Write a network's scalar output as $f(\boldsymbol\theta, \mathbf{x})$, a
function of parameters $\boldsymbol\theta$ and input $\mathbf{x}$.
Training changes $\boldsymbol\theta$; suppose it stays near its starting
value $\boldsymbol\theta_0$. A first-order Taylor expansion in the
parameters (not the input) gives the \emph{linearization}.

\begin{definition}[Linearization in parameters]
The \emph{linearized model} around $\boldsymbol\theta_0$ is the first-order
Taylor approximation in $\boldsymbol\theta$,
\[
  f(\boldsymbol\theta, \mathbf{x}) \;\approx\;
  f(\boldsymbol\theta_0, \mathbf{x})
  + \nabla_{\boldsymbol\theta} f(\boldsymbol\theta_0, \mathbf{x})
    \cdot (\boldsymbol\theta - \boldsymbol\theta_0).
\]
This is \emph{linear in the parameter shift}
$\boldsymbol\theta-\boldsymbol\theta_0$, with a fixed \emph{feature map}
$\boldsymbol\phi(\mathbf{x}) =
\nabla_{\boldsymbol\theta} f(\boldsymbol\theta_0, \mathbf{x})$ (the
gradient of the output with respect to the parameters, frozen at
$\boldsymbol\theta_0$). It is still nonlinear in $\mathbf{x}$, because
$\boldsymbol\phi(\mathbf{x})$ can be a complicated function of
$\mathbf{x}$.
\end{definition}

\begin{definition}[Neural tangent kernel \citep{jacot2018}]
The \emph{neural tangent kernel} (NTK) is the inner product of two such
feature maps,
\[
  K(\mathbf{x}, \mathbf{x}')
  = \big\langle \nabla_{\boldsymbol\theta} f(\boldsymbol\theta, \mathbf{x}),\,
    \nabla_{\boldsymbol\theta} f(\boldsymbol\theta, \mathbf{x}') \big\rangle
  = \sum_{p} \frac{\partial f}{\partial \theta_p}(\boldsymbol\theta,\mathbf{x})\,
    \frac{\partial f}{\partial \theta_p}(\boldsymbol\theta,\mathbf{x}'),
\]
summing over all parameters $\theta_p$. A \emph{kernel} is just a symmetric
function $K(\mathbf{x},\mathbf{x}')$ that measures similarity between
inputs by an inner product of feature vectors; it lets us do linear
algebra in a feature space without writing the features out. The NTK is the
kernel whose features are the parameter-gradients of the network.
\end{definition}

Under gradient flow the network's \emph{outputs} on the training set move
according to this kernel. If $u_i(t) = f(\boldsymbol\theta(t),\mathbf{x}_i)$
collects the outputs and $\mathbf{y}$ the targets, then
\citep{jacot2018}
\[
  \frac{d\,\mathbf{u}(t)}{dt} = -\,H(t)\,\big(\mathbf{u}(t) - \mathbf{y}\big),
  \qquad H(t)_{ij} = K(\mathbf{x}_i, \mathbf{x}_j; \boldsymbol\theta(t)),
\]
i.e.\ \emph{kernel gradient descent in function space} driven by the NTK
Gram matrix $H$. Two facts make this useful: the cost is \emph{convex in
function space} (even though it is wildly non-convex in
$\boldsymbol\theta$), and the dynamics relax fastest along the top
eigenvectors (principal components) of $H$ -- which is a clean theoretical
motivation for early stopping \citep{jacot2018}.

\subsection{Width, initialization scale, and Jacot's theorems}

Two quantities decide whether the kernel picture is exact: \emph{width}
and \emph{initialization scale}.

\begin{definition}[Width and initialization scale]
The \emph{width} of a hidden layer is its number of neurons; ``infinite
width'' means letting every hidden layer's neuron count grow without
bound. The \emph{initialization scale} is the overall size of the random
parameters at the start, controlled by a factor $\alpha>0$ multiplying the
output (or equivalently the variance of the initial weights). Both are
``knobs'' that, turned up, make the network behave more like its fixed
linearization.
\end{definition}

Jacot's two theorems say that in the infinite-width limit the NTK becomes
deterministic and frozen \citep{jacot2018}.

\begin{theorem}[NTK at initialization \citep{jacot2018}, stated]
\label{thm:jacot1}
As all hidden widths $\to\infty$, the NTK at initialization converges (in
probability) to a deterministic limiting kernel $\Theta_\infty$ that
depends only on the choice of nonlinearity, the depth, and the
initialization variance -- not on the particular random draw of weights.
\end{theorem}

\begin{theorem}[NTK stays constant during training \citep{jacot2018},
stated]
\label{thm:jacot2}
Under mild smoothness, in the same infinite-width limit the NTK stays equal
to $\Theta_\infty$ \emph{throughout} training. Hence training is exactly
kernel gradient descent with one fixed kernel; for a least-squares loss the
network function follows a linear differential equation and the outputs
relax as $\mathbf{u}(t) = \mathbf{y} + e^{-H t}(\mathbf{u}(0)-\mathbf{y})$.
\end{theorem}

We state Theorems~\ref{thm:jacot1}--\ref{thm:jacot2} \emph{without proof}:
they rest on a central-limit-theorem argument over the many neurons and on
controlling how the kernel drifts during training, both of which are
genuinely involved \citep{jacot2018}. The intuition, though, is simple and
worth internalizing: \emph{a wide network has so many neurons that each one
only has to move a tiny amount to fit the data, and if the parameters move
only a little, the first-order Taylor approximation (the linearization)
stays accurate.} One honest subtlety from \citet{jacot2018} (their
Remark~4): each individual neuron's contribution to the kernel does shrink
with width, but the \emph{collective} contribution of the lower layers does
not, which is why even ``lazy'' lower layers still participate.

\subsection{Lazy versus rich, and Chizat's criterion}

We can now name the two regimes. Both terms are now standard, traceable to
\citet{jacot2018} and \citet{chizat2019}.

\begin{definition}[Lazy and rich training]
Training is \emph{lazy} (also called the \emph{kernel} or
\emph{fixed-features} regime) when the NTK barely changes, so the network
behaves like its linearization around $\boldsymbol\theta_0$: it is doing
kernel regression with features \emph{fixed at initialization}, and the
parameters hardly move \citep{chizat2019}. Training is \emph{rich} (the
\emph{feature-learning} regime) when the NTK \emph{evolves} during
training, so the features adapt to the task; this regime lives outside
kernel theory (it is described by mean-field / optimal-transport analyses).
\end{definition}

\citet{chizat2019} pinned down \emph{when} training is lazy with one
dimensionless number. Introduce an explicit scale factor $\alpha$ via the
rescaled objective $F_\alpha(\mathbf{w}) = \tfrac{1}{\alpha^2}
R(\alpha\,h(\mathbf{w}))$. Their criterion (for the square loss $R$) is

\begin{equation}
\label{eq:kappa}
  \kappa \;=\; \frac{\|h(\mathbf{w}_0)-\mathbf{y}^\star\|\;\|D^2 h(\mathbf{w}_0)\|}
                    {\|D h(\mathbf{w}_0)\|^2} \;\ll\; 1
  \quad\Longrightarrow\quad \text{training is lazy},
\end{equation}

where $D h$ and $D^2 h$ are the first and second derivatives of the model
in the parameters. Small $\kappa$ means the model's curvature ($D^2 h$) is
negligible relative to its slope ($Dh$) over the distance training has to
travel, so the linearization holds. The key scaling facts
\citep{chizat2019}:
\begin{itemize}[nosep]
\item \emph{Initialization scale.} Under rescaling by $\alpha$,
  $\kappa_{\alpha h} = \kappa / \alpha$. So \emph{turning up the scale
  $\alpha$ forces laziness}: large $\alpha$ $\Rightarrow$ small $\kappa$
  $\Rightarrow$ lazy. This is the scaling law $\kappa_\alpha = \kappa/\alpha$
  in Chizat et al.\ \citep{chizat2019}.
\item \emph{Width.} For a one-hidden-layer net the gradient norm grows like
  $\sqrt{m}$ (with width $m$) while the curvature stays $O(1)$, so
  $\kappa\sim O(1)/(\sqrt{m})^2 = O(1/m)\to 0$. Width is a second route to
  the same small-$\kappa$ laziness.
\item \emph{Parameters barely move.} In the lazy regime
  $\sup_t\|\mathbf{w}_\alpha(t)-\mathbf{w}_0\| = O(1/\alpha)$, confirming
  ``the parameters hardly change.''
\end{itemize}

\begin{remark}[Chizat's honest counterpoint]
\citet{chizat2019} stress that lazy training is \emph{not} the secret of
deep learning's success. Their experiments show lazy-trained deep CNNs
\emph{underperform} ordinary (rich) training. Laziness gives clean theory
and guaranteed fast fitting, but at the cost of recovering a fixed-kernel
method; the practical wins come from feature learning. We carry this
caveat into the bridge: the lazy/rich distinction is a useful conceptual
axis, not a claim that real networks are lazy.
\end{remark}

\subsection{Worked examples: NTK of simple models}

\begin{example}[NTK of a linear model is constant]
\label{ex:ntklinear}
Let $f(\boldsymbol\theta,\mathbf{x}) = \boldsymbol\theta\cdot\mathbf{x}$
(plain linear model). Then $\nabla_{\boldsymbol\theta} f = \mathbf{x}$ for
\emph{every} $\boldsymbol\theta$, so
\[
  K(\mathbf{x},\mathbf{x}') = \langle \mathbf{x}, \mathbf{x}'\rangle
  = \mathbf{x}\cdot\mathbf{x}',
\]
which does not depend on $\boldsymbol\theta$ or on time at all. A linear
model is ``perfectly lazy'': its NTK is the constant kernel
$\mathbf{x}\cdot\mathbf{x}'$, and training it \emph{is} kernel / linear
regression. (In the scalar case $f=\theta x$ the kernel is just $xx'$.)
\end{example}

\begin{example}[NTK of one hidden unit depends on initialization]
\label{ex:ntkunit}
Let $f(\boldsymbol\theta,x) = b\,\sigma(a x)$ with parameters
$\boldsymbol\theta=(a,b)$ and a smooth activation $\sigma$. Then
$\nabla_{\boldsymbol\theta} f = \big(b\,\sigma'(ax)\,x,\ \sigma(ax)\big)$,
so
\[
  K(x,x') = b^2\,\sigma'(ax)\,\sigma'(ax')\,x x' \;+\; \sigma(ax)\,\sigma(ax').
\]
Unlike the linear case, this \emph{depends on $(a,b)$} -- so at random
initialization the kernel is random, and at finite width it drifts as
$(a,b)$ move during training. What infinite width buys is exactly the
freezing of this dependence into a deterministic $\Theta_\infty$
(Theorem~\ref{thm:jacot1}).
\end{example}

\begin{problem}[Rescaled-model laziness]
\label{ex:rescaled}
Consider a rescaled model $f_\alpha = \alpha\, g$ with the model arranged so
that $g(\mathbf{w}_0,\mathbf{x})=0$ at initialization. Using the criterion
\eqref{eq:kappa}, show $\kappa_{\alpha g} = \tfrac{1}{\alpha}\,\kappa_g$,
and conclude that the linearization error vanishes as $\alpha\to\infty$, so
the model trains lazily. (Extension: verify $\kappa\sim 1/m$ for a
one-hidden-layer net using the $\sqrt{m}$-versus-$O(1)$ scaling of the
gradient and curvature norms.)
\end{problem}

\subsection{Feature averaging and orbit-structured data}

The last definition is the mechanism that ties lazy/rich to the project's
real question: \emph{why a max-margin solution can fail to be robust}.

\begin{definition}[Feature averaging, informally; orbit-structured data]
\label{def:featureavg}
Many natural inputs are built from \emph{shifted (or otherwise
transformed) copies} of a few underlying patterns. Calling the set of all
shifts of a pattern its \emph{orbit} (the group-action term made precise in
the symmetry section of this primer), we say data are
\emph{orbit-structured} when each class is a union of such orbits: a cat is
a cat wherever it sits in the image. \emph{Feature averaging} is the
phenomenon where a network, instead of detecting the discriminative
pattern wherever it appears, learns a weight that effectively sums or
averages the signal across all positions of the orbit. Such an averaged
feature can separate the training data with a large margin, yet a small,
carefully placed perturbation that disturbs the averaging can flip the
prediction -- so the large-margin solution is \emph{not} robust
\citep{li2025feature}.
\end{definition}

We keep Definition~\ref{def:featureavg} at the level of motivation,
because none of the five learning-theory sources read here
(\citealp{soudry2018, lyuli2020, jitelgarsky2019, jacot2018, chizat2019})
proves a direct theorem of the form ``lazy/rich $\Leftrightarrow$
shift-invariant/robust.'' What they \emph{do} give us is the scaffolding:
a lazy network cannot learn a new invariance that is not already in its
initialization kernel (its features are frozen), whereas a rich network
can adapt its features and \emph{learn} shift-invariance rather than
inherit it. Whether the network's margin-maximizing solution is a robust,
invariance-respecting one or a brittle, feature-averaging one therefore
depends on which regime it trained in -- the conceptual hinge the project
investigates.

% =====================================================================

\subsection*{Closing bridge: how the Part-1 thread fits together}

\begin{quote}
\textbf{The one-paragraph story this whole part has been building.}
\emph{Shift-invariance} is, at bottom, a \emph{group-averaging} operation:
averaging a signal over all the shifts in its orbit is the Reynolds
projection onto the invariant subspace (the symmetry section of this
primer makes this an honest orthogonal projection when the shifts act by
orthogonal matrices). The classical \emph{descriptors} of that invariant
content are the \emph{power spectrum} and \emph{bispectrum} (the
shift-invariants section): a shift only rotates the phase of each Fourier
coefficient, leaving $|X|^2$ untouched, and the bispectrum keeps enough
relative phase to tell apart signals the power spectrum confuses.
\emph{Robustness} of a classifier is governed by two numbers we met here
and in the Lipschitz section: the \emph{margin} (Section~\ref{sec:margins})
and the \emph{Lipschitz constant}, with the certified robust radius
behaving like margin over Lipschitz constant. The reason margin is the
right currency is the implicit-bias result
(Section~\ref{sec:implicit}): \emph{plain gradient descent, with no
robustness term in the objective, already drifts toward the max-margin
solution} -- exactly the SVM solution / convex-hull closest points for
linear models, and a KKT point of the margin program for homogeneous
networks. Finally, whether the network \emph{learns} a shift-invariant,
robust representation or merely \emph{inherits} a fixed set of features
(and possibly a brittle, feature-averaging margin) depends on the
\emph{lazy-versus-rich} regime (Section~\ref{sec:ntk}): a lazy,
wide-or-large-scale network is stuck with its initialization kernel and
cannot learn a new invariance, whereas a rich network can. That is the
chain -- \emph{group-averaging $\rightarrow$ spectral descriptors
$\rightarrow$ margin and Lipschitz $\rightarrow$ implicit bias of gradient
descent $\rightarrow$ lazy vs.\ rich} -- on which the rest of the project
is built.
\end{quote}

% =====================================================================
% \citep keys used in this file (to be provided by the host bibliography):
%   vapnik           -- Vapnik 1995/1998, Statistical Learning Theory (max-margin / SVM origin)
%   boser1992        -- Boser, Guyon, Vapnik 1992 (original SVM)
%   cortes1995       -- Cortes, Vapnik 1995 (soft-margin SVM)
%   burges1998       -- Burges 1998 tutorial (canonical margin build-up; asserts distance formula)
%   bennett2000      -- Bennett, Bredensteiner 2000 (duality/geometry, convex-hull closest points)
%   soudry2018       -- Soudry et al. 2018 JMLR (implicit bias, linear separable max-margin)
%   lyuli2020        -- Lyu, Li 2020 ICLR (homogeneous nets, normalized margin -> KKT)
%   jitelgarsky2019  -- Ji, Telgarsky 2019 (non-separable ray / risk-margin)
%   jacot2018        -- Jacot, Gabriel, Hongler 2018 NeurIPS (NTK; Thms 1-2)   [bib key in repo: jacot2018neural]
%   chizat2019       -- Chizat, Oyallon, Bach 2019 NeurIPS (lazy training, kappa criterion)
%   li2025feature    -- feature-averaging non-robustness (used only for the motivation in Def. of feature averaging)  [bib key in repo: li2025feature]
%   tsuzuku2018lipschitz -- Tsuzuku et al. 2018 (Lipschitz-margin certificate; referenced in the bridge)  [bib key in repo]
% Note: jacot2018, li2025feature and tsuzuku2018 already exist in paper/report/references.bib
%       as jacot2018neural, li2025feature, tsuzuku2018lipschitz. If this fragment is
%       compiled against that bib, rename jacot2018 -> jacot2018neural. The SVM/implicit-bias
%       keys (vapnik, boser1992, cortes1995, burges1998, bennett2000, soudry2018, lyuli2020,
%       jitelgarsky2019, chizat2019) are NOT yet in references.bib and must be added.
% =====================================================================


%==============================================================================
\part*{Part II: Finding the questions}
\addcontentsline{toc}{section}{Part II: Finding the questions}
Part~I gave you the tools. This part gives you the situation that the project started from and asks you to generate
the research questions yourself. Each problem is presented with candidate directions and their trade-offs, and no
answer: commit to a direction before you read Part~III.

% =====================================================================
%  PRIMER PART II -- FINDING THE QUESTIONS
%  Comes AFTER Part I (background), BEFORE Part III (results/proofs).
%  Purpose: take a reader who has the background and teach them to
%  GENERATE and CHOOSE the project's research questions themselves,
%  WITHOUT revealing the answers (those live in Part III).
%
%  This is a FRAGMENT: no \documentclass, no \begin{document}.
%  The assembler supplies the preamble, which already defines:
%  definition, theorem, lemma, proposition, corollary, example
%  (printed as "Worked example"), remark, problem, and the macros
%  \R \Z \E \sign \Lip \one \conv \dist \Orb \fdc \Vinv \Vac \Sep
%  \Piinv etc.  "Find-the-question" prompts use the {problem}
%  environment plus \paragraph.
%
%  Citation keys used appear in the comment block at the END of file.
%  NEW keys (not yet in primer.bib) are flagged there.
% =====================================================================

% NOTE: the assembler (primer_v2.tex) already emits the \part* heading and its
% table-of-contents line for Part II, so this fragment must NOT repeat them.

Part~I handed you a toolbox: a way to measure how robust a classifier is, the Lipschitz--margin
inequality (a margin divided by a Lipschitz constant lower-bounds the certified radius) that turns
a smoothness budget and a margin into a certified radius,
the Fourier/projector machinery that says what a symmetry preserves, and the implicit-bias and
NTK lenses that say where gradient descent actually lands. This part does something different.
It does \emph{not} give you results. It puts you in the chair of the person who had to decide
\emph{which questions to ask in the first place}, hands you the same situation they faced, and asks
you to form your own hypotheses before Part~III tells you what happened.

That is a deliberate choice. The fastest way to misread a research paper is to learn its answers
and then reverse-engineer questions that make those answers look inevitable. The answers almost
never were inevitable. They were one defensible bet among several, taken at a fork where reasonable
people would have gone different ways. If you can stand at that fork yourself and feel the pull of
each branch, you will read Part~III as a sequence of decisions rather than a sequence of
revelations, and you will be able to do the same thing on a problem nobody has solved yet.

So: read this part with a pencil. Every research question below is presented as a genuine open
problem with two to four candidate directions and their trade-offs. At each one you will be asked to
commit, on paper, to the direction \emph{you} would take and why, before turning to Part~III. The
goal is not to guess the ``right'' answer. It is to develop the instinct that lets you generate good
questions and triage them.

%======================================================================
\section{The situation: a clean story that broke}
\label{sec:p2-situation}

Research questions do not come from the sky. They come from a concrete irritant: something a
previous effort could measure but not explain. Here is the irritant, told as it actually unfolded,
because the shape of the chain is itself a lesson in where questions live.

\subsection{Act I --- a tidy hypothesis, reproduced}
\label{sub:p2-act1}

The starting point was a clean and quotable claim from the literature: \emph{shift invariance can
reduce adversarial robustness} \citep{ge2021shift}. The mechanism, for the linear case, is exactly
the kind of thing Part~I prepared you to read. A classifier that is invariant to all circular
shifts of its input must give the same answer to a signal and to every shift of it; for a
\emph{linear} classifier this forces the weight vector into the fixed subspace of the shift group,
which (Part~I, the Fourier section) is spanned by the constant vector $\one$. The model can then see
the data only through its mean, the DC (zero-frequency) component, and discards all the rest. On a
synthetic ``one black pixel versus one white pixel'' problem this collapses the achievable max-margin
from a constant down to $2/\sqrt{d}$ in the image dimension $d$ \citep{ge2021shift}. More invariance,
smaller margin, less robustness. A satisfying monotone story.

A team reproduced it. On small image families (MNIST, Fashion-MNIST) trained from a shared template,
they measured a \emph{shift-consistency} score --- the fraction of inputs whose predicted label is
unchanged when the input is shifted, an empirical stand-in for ``how invariant did this model turn
out to be'' --- and they measured adversarial robustness with a projected-gradient (PGD) attack.
Across the families the headline ordering held: \emph{higher consistency went with lower robustness},
just as predicted.

\subsection{Act II --- the deviations it could not name}
\label{sub:p2-act2}

Then the tidy story sprang leaks, and the leaks are the interesting part.

First, the ordering was only \emph{mostly} right. Specific architectures sat in the wrong place:
a couple of convolutional variants were more robust than their consistency score said they should
be, or less. The reproduction could record the deviation but had no second axis to explain it.
The hypothesis ``consistency orders robustness'' was being falsified at the margins by its own data,
and nobody could say what the better axis was.

Second, to turn the picture into something actionable --- ``give me a model that is \emph{both}
shift-invariant and robust'' --- the team drew a two-dimensional scatter (consistency on one axis,
robust accuracy on the other), split it into four quadrants using a fully-connected baseline as the
origin, and aimed for the desirable corner. It was a useful screening picture. But it was a picture,
not a principle: nothing in it said \emph{why} a given architecture lands in a given quadrant, or
whether the quadrants would survive a stronger attack, or whether they would persist as models grow.
The diagnostic worked and was unjustified at the same time, which is the most uncomfortable place a
result can sit.

Third, and most honestly, the write-up's own appendix for theory read, in effect, ``to do.''
Everything was empirical, on tiny images, with a single attack family. There was no certificate, no
worst-case attack, no measurement of the very margin whose collapse the base theory was about.

\paragraph{The lesson of Act II.} Notice where the questions are being born. They are not born in the
parts that worked. They are born at the \emph{deviations} (the models in the wrong quadrant), at the
\emph{unjustified tool} (the quadrant diagram), and at the \emph{explicit hole} (the empty theory
appendix). A reproduction that merely confirms a result generates no questions. A reproduction whose
seams show generates all of them. Train yourself to read for the seams.

\subsection{Act III --- the literature contradicts itself}
\label{sub:p2-act3}

If the base story were the whole story, the fix would just be ``measure more carefully.'' It is not,
because the supervising literature does not agree with itself. Four papers, read precisely, point in
different directions while using the same word, ``invariance.''

\begin{itemize}[leftmargin=1.4em]
  \item \textbf{Invariance hurts --- conditionally.} \citet{ge2021shift} is the base result above.
        But read past the headline and Ge themselves add a crucial nuance: on a dataset whose
        classes live in orthogonal frequencies, where shifting does \emph{not} inflate the effective
        dimension, the convolutional network is \emph{more} robust than the fully-connected one. Their
        own conclusion is that invariance hurts only when it shrinks the margin, and that
        ``the linear margin predicts robustness better than the dimension of the data''
        \citep{ge2021shift}. The base paper's own fine print already suspects the headline axis is
        wrong.
  \item \textbf{Invariance trades off --- provably.} \citet{kamath2020} show empirically that training
        for invariance to larger rotations shrinks the average distance to the decision boundary and
        hurts $\ell_\infty$ robustness; \citet{kamath2021} prove a trade-off on a synthetic
        distribution built from a cyclic error-correcting code, and offer a curriculum that buys a
        better Pareto point without beating the single-objective adversarial baseline. Their
        ``rate of invariance'' $\Pr(f(TX)=f(X))$ is structurally the same quantity as the
        consistency score from Act~I --- yet neither Kamath paper cites Ge. Two independent lines,
        one seam between them.
  \item \textbf{A specific architecture choice causes vulnerability.} \citet{galloway2019} argue that
        batch normalization, by fixing each neuron's first two moments, destroys input-space
        distances and drops PGD robustness by amounts (roughly $8$--$20\%$ across datasets)
        comparable to the architecture effects Act~I reported. This is not a contradiction so much as
        a warning: any cross-architecture robustness comparison is confoundable as a normalization
        effect unless normalization is controlled.
  \item \textbf{Invariance helps.} The anti-aliasing line says the opposite of Ge. Inserting a
        low-pass blur before downsampling restores shift-equivariance, and \citet{zhang2019blurpool}
        reports it \emph{improves} shift-consistency and clean accuracy and average-case corruption
        robustness at once. Crucially, Zhang makes no worst-case $\ell_p$ adversarial claim; that is
        left as future work. The natural ``invariance helps'' counter-voice is about
        \emph{average-case} robustness, not worst-case.
\end{itemize}

State the contradiction plainly. Same word, opposite verdicts: one camp says (conditional)
invariance hurts worst-case $\ell_p$ robustness; another reports it helps; a third proves a
trade-off on a cyclic-shift construction; and two of the three do not cite the base result. A
contradiction in the literature is not a nuisance to be smoothed over in related work. It is a
flashing sign that some shared assumption is wrong, and the assumption is almost always a word doing
too much work. Here the overloaded word is ``invariance.''

\subsection{Act IV --- the gap the wider field leaves open}
\label{sub:p2-act4}

Widen the search and the contradiction sharpens into a precise hole. There is a theory camp that
shows invariance or implicit bias can hurt robustness --- but never about shift. \citet{tramer2020}
prove that defending against $\ell_p$ sensitivity forces \emph{excessive invariance}, opening a
distinct invariance-based attack that changes the true label while the prediction stays put; but
their ``invariance'' is over-invariance to an $\ell_p$ ball, not a group symmetry.
\citet{frei2023} prove that gradient descent's max-margin implicit bias on two-layer ReLU networks
reaches solutions that generalize yet are \emph{provably non-robust}, even when robust solutions
fitting the same data exist; \citet{li2025feature} pin the mechanism on feature averaging and
recover robustness with finer supervision; \citet{minvidal2024} flip the bias from non-robust to
robust by swapping the activation. All three are about the optimizer's bias on synthetic, near-
orthogonal clusters \emph{with no group symmetry}, under $\ell_2$. Their lever on robustness is
activation or supervision, not translation.

And there is an empirical/linear camp that says symmetry helps --- but not in a way that closes the
gap. \citet{sahagokhale2024} state verbatim that ``better shift invariance is generally correlated
with better adversarial robustness,'' and show the effect is architectural rather than a capacity
artifact --- but it is a pure correlation, with PGD/FGSM only (so possibly gradient-masked, see
Part~I), one backbone, no causal isolation, and no theorem linking shift to any radius.
\citet{wang2025bridging} report that rotation/scale equivariance helps without adversarial training,
but that is the \emph{wrong group} (CNNs are already translation-equivariant), it is enforced rather
than emergent, and again no $\ell_p$ theorem. \citet{chenzhu2023} and \citet{lawrence2021} prove
that linear equivariant networks reach a wider margin or a characterized implicit bias --- but they
are \emph{linear only} and stop at margin/generalization, never an $\ell_p$ radius.

\paragraph{The gap, in one sentence.} Across every paper on both sides, the same thing is missing:
\emph{no result says what discriminative signal survives a given translation-invariance, turns that
into an $\ell_p$ robustness certificate, and validates it under a standardized strong attack at
matched capacity.} The ``helps'' side is empirical, or linear, or the wrong group. The ``hurts''
side is about $\ell_p$-ball over-invariance or optimizer bias on symmetry-free clusters. The seam
between them is exactly where the Act~I reproduction's monotone story tore.

This is the situation. A reproduced result whose deviations have no second axis; a working diagnostic
with no principle; a four-paper contradiction; and a field-wide gap with a precise shape. Everything
that follows in this part is a question you could have asked while standing here. Now you get to ask
them.

%======================================================================
\section{Five open problems, and the forks inside them}
\label{sec:p2-problems}

For each question I give you the \emph{situation} (what you know and what is unsettled), the
\emph{candidate directions} you could take with their honest trade-offs, and a prompt asking you to
commit before reading on. I will not name the quantity, the theorem, or the experiment that resolves
it. Resist the urge to flip ahead. The value is entirely in deciding under uncertainty.

A note on the format. ``Direction'' means a research strategy, not a guaranteed win. Each has a cost.
Part of taste is reading those costs accurately, because a direction that cannot be wrong is usually
a direction that cannot teach you anything either.

%----------------------------------------------------------------------
\subsection{Q1 --- What single quantity actually orders robustness?}
\label{sub:p2-q1}

\paragraph{Situation.} Act~I gave you a metric, shift-consistency, that orders robustness
\emph{most} of the time and is falsified at the margins by its own outliers. Act~III gave you the
base paper quietly conceding that ``the linear margin predicts robustness better than the dimension
of the data.'' You suspect consistency is a symptom, not a cause. The deeper question is whether a
\emph{single measurable scalar} even exists that you can attach to a trained model and that ranks
models by their robust radius --- and if one does, which family of quantities it belongs to. You are
genuinely choosing among candidate \emph{families} here, not polishing a known winner: a
\emph{consistency} score, a bare \emph{margin}, a \emph{sensitivity} measure (a Lipschitz or
gradient-norm proxy), or some \emph{ratio} of a margin to a sensitivity. Each is a different bet
about what governs the robust radius, and when your chosen scalar disagrees with consistency, you
want it to be the one that is \emph{right}.

\paragraph{Candidate directions.}
\begin{itemize}[leftmargin=1.4em]
  \item \emph{Refine consistency.} Keep the consistency axis but measure it better --- average over
        more shifts, weight by frequency, use a soft (probability) version instead of a hard
        label-flip count. \textbf{Trade-off:} cheap, continuous with the prior work, and it might
        clean up the outliers. But it doubles down on an axis the base paper's own nuance already
        suspects. If consistency is the wrong \emph{kind} of quantity, no amount of polishing fixes
        the outliers; you will have built a better thermometer for the wrong temperature.
  \item \emph{Bet on a single geometric quantity.} Drop consistency and pick \emph{one} geometric
        scalar instead: a bare margin (distance to the boundary), a bare sensitivity (a Lipschitz or
        gradient-norm estimate), or a \emph{margin-to-sensitivity ratio} that trades the two against
        each other. Part~I gives you reason to take the ratio family seriously --- the
        Lipschitz--margin certificate lower-bounds the robust radius by a margin over a Lipschitz
        constant --- but the certificate is background, not a verdict: it does not tell you that a
        ratio \emph{orders} models better than a bare margin or a bare sensitivity does, nor in which
        feature space to compute any of them. \textbf{Trade-off:} principled, and each candidate
        ties to geometry you can compute. But a margin and a Lipschitz constant are both nontrivial
        to measure on a real network; the certificate inequality is one-directional (a lower bound,
        not an equality --- you will revisit this in Q3); and you must decide \emph{which} feature
        space (raw pixels? the space the invariance leaves behind?) to measure in. Picking the wrong
        family, or the wrong space, is the way this direction fails.
  \item \emph{Let the data choose by regression.} Throw a basket of candidate scalars (consistency,
        margin, Lipschitz, gradient norms, boundary distance) at a panel of models and report which
        correlates best with measured robustness. \textbf{Trade-off:} maximally empirical and
        honest about what predicts what; it cannot be accused of theory-bias. But correlation on a
        model zoo confounds everything with everything (capacity, normalization, training), and a
        leaderboard of correlations is not an explanation --- it tells you \emph{that} something
        orders robustness, never \emph{why}.
\end{itemize}

\begin{problem}[commit before Part III]
First decide whether you even believe a single scalar can order robustness across architectures; if
not, say what minimal set of scalars you would need instead. If you do, pick the family you would bet
on --- consistency, a margin, a sensitivity, or a margin-to-sensitivity ratio --- and name the single
most likely reason your bet fails. If you chose a ratio or any geometric scalar, write down which
feature space you would measure it in and why raw pixels might be the wrong choice. If you chose
regression, write down the one confound you would control first.
\end{problem}

%----------------------------------------------------------------------
\subsection{Q2 --- What discriminative signal survives an imposed invariance?}
\label{sub:p2-q2}

\paragraph{Situation.} The linear story is brutally clean: force shift-invariance on a linear model
and it sees only the mean. But real networks are not linear classifiers on raw pixels; they push
data through nonlinear feature maps first. If a convolution followed by a squaring nonlinearity and
a pooling step is invariant to shifts, it is \emph{not} reduced to the mean --- it keeps something
richer. You want to know \emph{exactly what}. Not ``it keeps more''; the precise functional. Because
whatever survives the invariance is the only material left to build a margin out of, and Q1 told you
margin is the thing that matters.

\paragraph{Candidate directions.}
\begin{itemize}[leftmargin=1.4em]
  \item \emph{Quotient and diagonalize.} Treat the invariance as a projection: the model can only
        depend on the part of the input that the group leaves fixed, so write that projector down and
        diagonalize the group action (Fourier, for cyclic shifts) to read off the surviving
        coordinates in closed form. \textbf{Trade-off:} this is exactly the machinery Part~I built,
        it gives an \emph{exact} answer for the linear and quadratic cases, and it makes the margin
        computable. But it is cleanest for a single, idealized invariant layer; pushing it through a
        deep stack of real layers is where the rigor frays.
  \item \emph{Climb the ladder of classical invariants.} There is a known hierarchy: the mean
        (degree~1), the power spectrum / autocorrelation (degree~2, what a square-and-pool computes),
        the bispectrum (degree~3), and beyond. Ask which rung a given architecture sits on and what
        it can therefore separate. \textbf{Trade-off:} it connects the architecture to a beautiful,
        well-understood body of signal-processing theory, and it predicts \emph{failure} cases ---
        classes that one rung cannot tell apart but the next can. But higher rungs cost more
        (a larger Lipschitz constant; recall that any margin-to-sensitivity ratio has a sensitivity
        in its denominator), and you must show the
        architecture \emph{actually} realizes the rung, not just that the rung exists.
  \item \emph{Probe empirically.} Skip the analysis: train the invariant model and read out which
        input perturbations it is blind to, mapping the surviving signal by experiment.
        \textbf{Trade-off:} works for any architecture, no idealization. But it gives you a point
        cloud, not a formula; it will not tell you the \emph{margin} the surviving signal supports,
        and it cannot certify anything.
\end{itemize}

\paragraph{A pointed sub-question.} Recall the base paper's worst case: the ``one hot pixel''
example, $\pm e_j$, where the class is encoded by \emph{where} the spike sits. Two such inputs that
differ only by a shift have the \emph{same} power spectrum --- the degree-2 invariant is blind to
location. So if your architecture stops at degree~2, it cannot separate them at all, and the forced
trade-off bites hardest exactly there. Does a richer invariant rescue this corner, and at what cost
to the Lipschitz denominator? Hold that thought; it is a live fork.

\begin{problem}[commit before Part III]
For a convolution--square--pool layer, write your best guess at the surviving signal in one line.
Then decide: would you try to \emph{defeat} the hot-pixel worst case with a higher-order invariant,
or would you accept it as a genuine limit of low-degree invariance? State the price of your choice in
terms of the margin and the sensitivity (Lipschitz) it would move.
\end{problem}

%----------------------------------------------------------------------
\subsection{Q3 --- Does the surviving signal certify an $\ell_p$ radius?}
\label{sub:p2-q3}

\paragraph{Situation.} Suppose Q2 hands you an exact margin $\eta$ for the surviving signal, and
suppose you are pursuing the margin-to-sensitivity ratio of Q1's second direction, so the object on
the table is the Lipschitz--margin certificate $r_2 \ge \eta/L$ from Part~I. It is a \emph{lower}
bound: it guarantees no adversary inside radius $\eta/L$, but it says nothing about how much
\emph{further} the true safe radius might extend. For this to be the missing link the whole field
lacks --- a result tying translation-invariance to an actual $\ell_p$ radius --- you have to decide
what kind of statement you are even after. (Whether this ratio actually \emph{orders} models better
than its rivals is still Q1's open business; here you are only asking what the certificate, taken on
its own terms, can and cannot promise.)

\paragraph{Candidate directions.}
\begin{itemize}[leftmargin=1.4em]
  \item \emph{Settle for a one-sided certificate (and predictor).} Accept $r_2 \ge \eta/L$ as a
        guarantee and an empirical predictor, and do not claim it pins the radius exactly.
        \textbf{Trade-off:} honest and immediately usable; it is enough to \emph{order} models and to
        \emph{guarantee} safety up to the bound. But a lower bound alone cannot tell you whether a
        model with a small certificate is actually fragile or merely has a loose bound --- and a
        referee will ask whether your quantity \emph{characterizes} robustness or merely lower-bounds
        it.
  \item \emph{Chase a converse.} Try to prove an upper bound too, so the ratio brackets the true
        radius from both sides. \textbf{Trade-off:} a two-sided bracket is far stronger and would
        make the quantity a genuine measure of robustness, not just a safety floor. But you should
        first ask whether a \emph{universal} converse can even exist --- look for a degenerate example
        where the ratio is tiny yet the true radius is large. If you can build one, an unconditional
        converse is dead, and the honest move is to find the extra \emph{condition} under which a
        converse holds. (This is the search for what has no converse, a research move in its own
        right; see Section~\ref{sec:p2-taste}.)
  \item \emph{Add a second, invariance-specific radius.} The certificate above is the
        \emph{additive-noise} radius. But Tramer's excessive-invariance attack changes the true label
        \emph{along the orbit} while the prediction stays fixed --- a different failure the additive
        radius does not see. You could certify a separate ``orbit-flip'' radius for that mode.
        \textbf{Trade-off:} it captures a real and distinct threat and makes the ``invariance can
        also \emph{hurt}'' direction precise. But it complicates the story --- now there are two
        radii to report, and you must say when each one is the binding constraint.
\end{itemize}

\begin{problem}[commit before Part III]
Decide whether you would chase a converse or settle for a one-sided certificate. If you would chase
it, sketch the degenerate example you would build \emph{first} to check whether an unconditional
converse is even possible --- a function whose certificate is loose by an unbounded factor. If you
would not, say what claim you are giving up.
\end{problem}

%----------------------------------------------------------------------
\subsection{Q4 --- Does any of this survive a strong attack at matched capacity?}
\label{sub:p2-q4}

\paragraph{Situation.} Acts~I--IV are littered with the same two methodological landmines. First,
PGD/FGSM attacks can \emph{appear} to show robustness that vanishes under a stronger evaluation
because the gradients were masked, not the model genuinely robust (Part~I). The ``invariance helps''
evidence is almost entirely PGD/FGSM. Second, the architectures being compared differ in
\emph{capacity} (and in normalization, per \citet{galloway2019}), so any robustness gap might be a
capacity or normalization gap wearing an invariance costume. You want to know whether any geometric
story you settle on in Q1--Q3 is \emph{real} or an artifact of weak attacks and unmatched models.

\paragraph{Candidate directions.}
\begin{itemize}[leftmargin=1.4em]
  \item \emph{Trust a stronger attack ensemble.} Re-evaluate every model with a standardized,
        parameter-free, ensemble attack designed precisely to catch gradient masking, and treat that
        number as ground truth. \textbf{Trade-off:} this is the community standard and the only way to
        kill the gradient-masking objection. But it is expensive, and a single robust-accuracy number
        at one perturbation budget hides the \emph{radius} you actually care about.
  \item \emph{Measure the radius directly with a min-norm attack.} Instead of ``what fraction
        survives at budget $\epsilon$,'' find, per input, the \emph{smallest} perturbation that flips
        it --- the empirical robust radius. \textbf{Trade-off:} this is the right quantity to compare
        against an $\eta/L$ \emph{radius} certificate, and it gives a distribution, not one number.
        But min-norm attacks are themselves attacks and can be fooled; you should cross-check them
        against the ensemble above.
  \item \emph{Control capacity by construction.} Build a family of architectures --- a plain one, an
        anti-aliased one, an exactly-cyclic-invariant one, a shift-augmented one --- with
        \emph{identical} parameter counts, the normalization confound pinned down, and the
        \emph{same} training. Then any robustness difference is attributable to the invariance, not
        to size. \textbf{Trade-off:} this is the strongest possible causal isolation and lets you ask
        not ``does invariance help?'' but ``which term, margin or Lipschitz, does \emph{this}
        invariance move?'' But matching capacity across structurally different models is fiddly, and
        you must verify the matched models really are comparable in every nuisance dimension.
  \item \emph{Stay on synthetic data where the truth is known.} Run everything on the constructions
        where the margin and radius are computable in closed form, sidestepping attack noise
        entirely. \textbf{Trade-off:} clean and unimpeachable, and it isolates mechanism. But a
        referee --- and the gap from Act~IV --- will say the open question is precisely about
        \emph{real} data under a \emph{strong} attack; synthetic-only re-opens the very gap you set
        out to close.
\end{itemize}

\begin{problem}[commit before Part III]
Design your headline experiment in three lines: which models you compare, what you hold fixed to
make the comparison causal, and which attack(s) you trust for the final number. Name the one
confound that, if a reviewer raised it, would most damage your conclusion --- and say how your design
already answers it.
\end{problem}

%----------------------------------------------------------------------
\subsection{Q5 --- Why do trained models land where the diagnostic put them?}
\label{sub:p2-q5}

\paragraph{Situation.} Q1--Q4 are about what a \emph{given} model's geometry implies. But the
Act~II quadrant diagram was about \emph{trained} models: gradient descent, not a fixed weight
vector, decides where each architecture ends up on the consistency-versus-robustness plane. So even
a perfect geometric theory leaves a gap: \emph{why does optimization land an architecture in the
quadrant it lands in?} Answer this and the heuristic quadrant diagram becomes a principle; leave it
and you have explained the map but not the territory. Part~I's implicit-bias and NTK lenses are the
tools here.

\paragraph{Candidate directions.}
\begin{itemize}[leftmargin=1.4em]
  \item \emph{Reduce to a regime where the bias is known.} In the lazy/NTK regime the trained
        network behaves like a fixed kernel method, where ``what gradient descent selects'' has a
        clean answer. Ask what enforcing invariance does to that kernel's solution. \textbf{Trade-
        off:} tractable, and it can turn an empirical trend into a theorem about the selected
        solution. But the lazy regime is an idealization (wide, near-linear training), the activation
        matters (a smooth activation behaves very differently from bare ReLU under this lens), and
        what holds lazily may not hold in the feature-learning regime real training lives in.
  \item \emph{Split by data geometry.} The Act~III/IV papers suggest the answer is conditional. On
        data whose classes are \emph{closed under the symmetry} (every shift of a class member is in
        the class), enforcing invariance may cost nothing --- gradient descent might symmetrize for
        free. On near-orthogonal \emph{cluster} geometry (the Frei/Li/Min-Vidal setting) invariance
        may genuinely help the ratio. Prove the easy case, conjecture the hard one.
        \textbf{Trade-off:} this matches the contradiction's actual shape --- different geometries,
        different verdicts --- and gives you a provable positive result plus an honest open problem.
        But it splits the question into pieces, one of which you will have to leave open, and you must
        be candid that the cluster case resists a full trajectory-level proof.
  \item \emph{Stay empirical and characterize the landing pattern.} Train the matched families from
        Q4 across seeds and data regimes and report \emph{where} they land and how tightly,
        upgrading the quadrant diagram into a measured phenomenon without proving it.
        \textbf{Trade-off:} robust and honest, and it documents exactly the regime-dependence the
        theory should later explain. But ``we observed where models land'' is a description, not the
        principle the Act~II diagram was missing; a reviewer asks \emph{why}.
\end{itemize}

\begin{problem}[commit before Part III]
Decide whether you would attack this in the lazy regime, split it by data geometry, or characterize
it empirically --- and say which part of the question you are willing to leave \emph{open}. (Honest
answer: at least one piece of Q5 is still open in the project. Predict which piece, and why it is the
hard one.)
\end{problem}

%======================================================================
\section{Research taste: how to find questions like these}
\label{sec:p2-taste}

The five questions above did not arrive labelled. Someone had to manufacture them from a tangle of
contradictory papers and an unexplained scatter plot. This section is about that manufacturing.
It is not a checklist --- taste is not a checklist --- but there are recurring \emph{moves} that
generate good questions, and there are real \emph{criteria} for deciding which questions are worth
the months they cost. I illustrate each move with a situation from this project, without giving the
answer, so you can practice the move rather than memorize the result.

\subsection{Recurring moves}
\label{sub:p2-moves}

\paragraph{Quotient out the nuisance.} When a symmetry or a degree of freedom is making the problem
look complicated, do not fight it --- divide by it. Ask what is left once the nuisance is collapsed.
\emph{In this project:} ``the model is shift-invariant'' looks like a statement about infinitely many
shifted inputs. Quotient by the shift group and it becomes a statement about a single object --- the
projection of the data onto the part the group fixes. The infinitely-many-shifts problem becomes a
one-projector problem. (This is the engine behind Q2.)

\paragraph{Diagonalize the symmetry.} A symmetry that commutes with your operator can be diagonalized,
and in the eigenbasis the operator is just multiplication --- everything decouples. \emph{In this
project:} circular shifts are diagonalized by the Fourier basis, so an invariant's action on the data
becomes coordinate-wise on frequencies, and ``what survives invariance'' becomes ``which frequency
coordinates the invariant can still read.'' The hard algebra of convolutions turns into reading off a
diagonal. Whenever you see a group acting linearly, ask for its eigenbasis first.

\paragraph{Bracket what has no converse.} A one-sided inequality is a confession that you do not yet
know the other side. Before assuming both sides hold, deliberately hunt for a counterexample to the
converse. If you find one, you have learned that the clean two-sided statement is \emph{false in
general} --- which is itself a result --- and the real question becomes ``under what extra condition
does a converse hold?'' \emph{In this project:} the certificate $r_2 \ge \eta/L$ is one-sided; the
right reflex is to ask whether a universal converse can exist before trying to prove one. (This is
the heart of Q3, and the degenerate-example reflex below.)

\paragraph{Find the degenerate case.} The fastest way to test a claim or kill a converse is to push
it to an extreme where one term blows up or collapses, and see if the claim survives. \emph{In this
project:} to probe whether $\eta/L$ pins the radius, you reach for a function whose curvature can be
cranked arbitrarily high so the certificate becomes arbitrarily loose while the true radius does not
move. A single such example settles a general question instantly. Degenerate cases are cheap and
decisive; build them \emph{before} you invest in a hard proof.

\paragraph{Reduce to the per-unit atom.} A network is a composition of many small pieces. Instead of
reasoning about the whole, find the smallest piece that carries the phenomenon and solve \emph{that}
exactly, then ask how the pieces compose. \emph{In this project:} rather than analyze a deep
invariant network, isolate one convolution--nonlinearity--pool unit and compute exactly what it keeps
and what margin it supports. The atom is solvable; the composition is the next question. If you
cannot solve the atom, you certainly cannot solve the network --- and if you can, you have a building
block and a sanity check for everything larger.

\paragraph{Transport a finite lemma to function space.} A clean fact about finite-dimensional vectors
(projections are non-expansive; averaging cannot increase a norm) often has an exact analogue in a
function space (a reproducing-kernel Hilbert space) that turns an empirical trend about trained
networks into a theorem. \emph{In this project:} ``averaging over the orbit cannot increase the norm''
is a one-line fact for vectors; transported to the kernel that describes a wide network's training, it
becomes a statement about the solution gradient descent selects. The move is: prove it where it is
trivial, then carry it across the bridge that your regime (e.g.\ the lazy/NTK lens from Part~I)
provides. The catch is always whether the bridge holds --- which is exactly why Q5's lazy direction
carries an activation caveat.

\subsection{Which questions are worth asking}
\label{sub:p2-worth}

Moves generate candidate questions faster than you can answer them. Triage is the other half of
taste. Three criteria did the real work here.

\paragraph{1. There is a small, testable instance.} A question you can probe on a two-frequency toy
or a closed-form linear case \emph{this afternoon} is worth a hundred you can only settle with a
six-month training run. The small instance tells you fast whether the big claim is even plausible,
and it doubles as a sanity check on the eventual general proof. Every one of Q1--Q5 has a synthetic
miniature where the answer is computable; that is not a coincidence, it is the selection filter.
If you cannot construct a small instance, you usually do not yet understand the question well enough
to ask it.

\paragraph{2. It sits exactly at a literature contradiction.} The most fertile questions live on the
seam where two credible papers disagree, because the disagreement guarantees that \emph{someone}
is wrong about \emph{something specific}, and resolving it cannot be a no-op. Q1 sits on the seam
between ``consistency orders robustness'' and Ge's own ``margin predicts better than dimension.''
Q5 sits on the seam between ``invariance hurts'' (Frei/Li) and ``invariance helps'' (Saha--Gokhale,
Chen--Zhu). A question that no published paper would dispute is rarely worth asking; a question that
forces two published papers to reconcile almost always is.

\paragraph{3. The answer changes a decision.} Ask what someone would \emph{do} differently once they
know the answer. Q1's answer changes which scalar you compute to screen architectures. Q4's answer
changes whether you trust a PGD number at all. Q3's answer changes whether you can advertise a
\emph{certificate} or only a \emph{predictor}. If both possible answers lead to the same action, the
question is decorative; drop it.

\paragraph{Be honest: some good-looking questions did not pan out.} Taste includes being wrong
gracefully. Several attractive questions in this project bought less than they promised. Polishing the
consistency metric (Q1's first direction) was tempting and continuous with prior work, but if the
axis is the wrong \emph{kind} of quantity, a better thermometer does not help --- a lesson you only
learn by trying it or by reasoning hard about why not. Chasing an \emph{unconditional} converse for
the certificate (Q3) is a dead end the degenerate-case move kills in an afternoon; the salvageable
question is the \emph{conditional} one, which is strictly humbler. And the cleanest part of Q5 ---
the optimization story on symmetric data --- yields a real theorem, while the part everyone wants ---
the cluster-geometry case in the realistic feature-learning regime --- stays a conjecture supported
by experiments, not a proof. A primer that pretended otherwise would be teaching you the wrong
lesson. Knowing which of your questions is provable, which is merely true, and which is only plausible
is not a footnote to the research. It \emph{is} the research.

\medskip

You now have the five questions, the moves that generate them, and the criteria that select them.
Before you read Part~III, go back to your five committed answers from the \textit{commit before
Part~III} prompts. Write, for each, the single observation that would most change your mind. Then,
and only then, turn the page.

% =====================================================================
%  CITATION KEYS USED IN THIS FRAGMENT
% =====================================================================
%  REUSED (already in primer.bib):
%    li2025feature   -- Li et al. "Feature averaging..." (the Li 2024 paper;
%                       primer.bib already keys it as li2025feature)
%
%  NEW KEYS -- must be added to primer.bib (real, verified references):
%    ge2021shift      -- Ge, Kuang & Zou (or as published), "Shift Invariance
%                        Can Reduce Adversarial Robustness", NeurIPS 2021.
%    kamath2020       -- Kamath, Deshpande, Subrahmanyam (et al.),
%                        "Invariance vs. Robustness of Neural Networks",
%                        ICLR 2020 / OpenReview 2020.
%    kamath2021       -- Kamath et al., "Can we have it all? On the trade-off
%                        between spatial and adversarial robustness ...",
%                        NeurIPS 2021 (introduces CuSP).
%    galloway2019     -- Galloway, Golubeva, Tanay, Moussa, Taylor,
%                        "Batch Normalization is a Cause of Adversarial
%                        Vulnerability", ICML 2019 workshop / arXiv:1905.02161.
%    zhang2019blurpool-- R. Zhang, "Making Convolutional Networks
%                        Shift-Invariant Again", ICML 2019 (BlurPool/anti-alias).
%    tramer2020       -- Tramer, Behrmann, Carlini, Papernot, Jacobsen,
%                        "Fundamental Tradeoffs between Invariance and
%                        Sensitivity to Adversarial Perturbations", ICML 2020.
%    frei2023         -- Frei, Vardi, Bartlett, Srebro, "The double-edged
%                        sword of implicit bias: generalization vs.
%                        robustness in ReLU networks", NeurIPS 2023.
%    minvidal2024     -- Min & Vidal, "Can Implicit Bias Imply Adversarial
%                        Robustness?", ICML 2024 (polynomial-ReLU flips bias).
%    sahagokhale2024  -- Saha & Gokhale, "Translation-Invariant Polyphase
%                        Sampling (TIPS)", WACV 2024.  NOTE: the local file is
%                        mislabeled "rajput2024_..."; correct authorship is
%                        Saha & Gokhale.
%    wang2025bridging -- (Wang et al.) "Bridging Symmetry and Robustness:
%                        ... Equivariance ...", 2025 (rotation/scale equivariance).
%    chenzhu2023      -- Chen & Zhu, "Implicit Bias of (Linear) Equivariant
%                        Steerable Networks", 2023.
%    lawrence2021     -- Lawrence et al., "Implicit Bias of Linear Equivariant
%                        Networks", 2021.
% =====================================================================


%==============================================================================
\part*{Part III: Results and proofs}
\addcontentsline{toc}{section}{Part III: Results and proofs}
Here are the answers to the Part~II problems, each motivated, labelled honestly (classical, clean rephrasing, or new
contribution), and proved in full. Compare them against the directions you committed to.

% =====================================================================
%  PRIMER PART 3 — RESULTS WITH FULL PEDAGOGICAL PROOFS
%  The answers to the questions raised in Part II, proved in detail
%  for a strong undergraduate who has done Part I (robustness,
%  Lipschitz/eta-L certificates, margins/convex hulls, implicit bias,
%  NTK/lazy training, Fourier/power-spectrum/bispectrum invariants).
%
%  This is a FRAGMENT: no \documentclass, no \begin{document}.
%  The assembler (primer.tex) supplies the preamble, which already
%  defines: definition, theorem, lemma, proposition, corollary,
%  conjecture, example, remark; and the macros
%  \R \Z \E \sign \Lip \one \conv \dist \Orb \Vinv \Vac \Piinv
%  \fdc \Sep .
%
%  Proofs are EXPANDED line-by-line for a strong undergraduate.
%  Citation keys used appear in the comment block at the END of file
%  (new keys are flagged).
% =====================================================================

\section{Results: what survives an invariance, and what it costs}
\label{sec:results}

Part~II asked five questions. Each came from a real friction point: the old project's monotone
``more consistency, less robustness'' story broke and could not be explained; the base papers
openly contradict each other on whether invariance helps or hurts; and no prior result ties a
\emph{translation} invariance to an $\ell_2$ robust radius. This part answers those questions with
theorems and full proofs. Two honesty rules run through the whole part, and we state them once,
plainly, so no proof reads as a bigger claim than it is.

\begin{remark}[Calibration: what is new here and what is borrowed]
\label{rem:calibration}
Most of the \emph{machinery} below is classical. Group averaging, fixed subspaces, the
Lipschitz--margin certificate, power-spectrum and bispectrum invariants, the implicit-bias
characterization of gradient flow, and the lazy/NTK picture are all established tools, taught in
Part~I and cited there. What this part contributes is narrow and specific:
(1)~an \emph{exact projection-margin accounting} of what discriminative signal a given invariance
keeps (Theorem~\ref{thm:fg-margin});
(2)~a \emph{two-sided bracket} on the robust radius (a pair of bounds, lower and upper, that trap the
true radius from both sides) that completes the missing upper arm of the known certificate chain
(Theorem~\ref{thm:sandwich});
(3)~a \emph{negative} optimization result, that on \emph{orbit-closed} data (data where every shift
of a training point is itself a training point of the same label) imposing invariance is margin-free
(Theorem~\ref{thm:selfsym});
(4)~a \emph{positive} optimization result in the \emph{lazy} (wide-network, fixed-kernel) regime,
that invariance lowers the Lipschitz constant at equal margin (Theorem~\ref{thm:lazy}).
We never call $\eta/L$ a ``new robustness quantity'': it is a standard certificate, and our angle
is only to compute it \emph{after} identifying the invariant feature space and then to track which
of its two parts (the margin $\eta$ or the Lipschitz term $L$) a given architecture moves.
\end{remark}

\begin{remark}[Two plain-language terms we will need]
\label{rem:terms}
A \emph{converse} to an inequality $A \le B$ is a bound in the other direction, $B \le C\,A$ for
some constant $C$: it says $A$ is not just a lower estimate of $B$ but pins $B$ down to within a
factor. A certificate $r_2 \ge \eta/L$ becomes \emph{useful as a predictor} only if it has a
converse, otherwise the true radius could be arbitrarily larger than the certificate.
A function $g$ is \emph{co-Lipschitz} (or lower-Lipschitz) with constant $\alpha$ along a path if it
\emph{cannot} change too slowly: $|g(z)-g(x)| \ge \alpha\,\|z-x\|_2$. Ordinary Lipschitz continuity
($|g(z)-g(x)| \le L\,\|z-x\|_2$) caps how fast $g$ changes; the co-Lipschitz bound puts a floor
under it. The certificate's missing converse is exactly a co-Lipschitz statement, as we will see.
\end{remark}

\begin{remark}[Notation: one margin, several names]
\label{rem:margin-notation}
The same conceptual quantity --- the \emph{margin}, ``how decisively is a point classified'' ---
wears different symbols depending on the setting, and we flag this once so it never confuses.
In Part~I's robustness section the \emph{logit} margin (top logit minus best competitor) was
$M_{F,x}$; in Part~I's learning section the \emph{geometric} (SVM) margin, a distance to the
boundary, was $\gamma$. In this part we write $\eta$ for the \emph{feature margin} of an invariant
score (the gap the surviving signal supports) and $\mu = y\,g(x)$ for the signed margin used inside
the radius bracket; for a max-margin separator these coincide, $\mu = \eta$. Likewise $r_2$ here is
the $\ell_2$ robust radius written $r(x)$ in Part~I. Finally, Part~III's Fourier proofs use the
\emph{unitary} DFT $\hat x_k = d^{-1/2}\sum_j x_j e^{-2\pi i jk/d}$ (so Parseval is an exact isometry),
whereas Part~I fixed the \emph{unnormalised} forward DFT (so Parseval carried a factor of $N$); both
use the same minus-sign root of unity $\omega = e^{-2\pi i/d}$.
\end{remark}


% =====================================================================
\subsection{Q2 answered: the invariant linear margin is a projected convex-hull distance}
\label{sub:fg-margin}

\paragraph{The question this answers.}
A shift-invariant linear classifier must give every circular shift of a signal the same score. So
it cannot use any feature that a shift can change. \emph{Which} discriminative signal is left? Part~I
showed (for cyclic shifts) that the fixed subspace is the one-dimensional line spanned by the
all-ones vector, so an invariant linear score sees only the mean (DC) component. The clean general
statement is not special to cyclic shifts or to the DC line: for \emph{any} finite group of
orthogonal symmetries, the surviving signal is the data \emph{projected onto the fixed subspace},
and the best achievable invariant margin is exactly half the distance between the two projected
convex hulls. This is the structural anchor for everything that follows.

\paragraph{Status.}
This is a clean generalization, not a new phenomenon. \citet{ge2021shift} prove the cyclic/DC special
case (an invariant linear classifier is governed by the DC gap), and group averaging onto a fixed
subspace is standard representation theory. The modest contribution is the general finite-group
statement with the convex-hull normalization of the margin; the cyclic case recovers Ge, including
the $2/\sqrt{d}$ collapse, as a corollary. We do not claim the cyclic effect as new.

We use one fact from Part~I about two finite point sets $A,B$ in $\R^m$: the largest margin of a
hyperplane that separates them, where the margin is the smallest distance from any point to the
hyperplane, equals $\tfrac12\dist(\conv A,\conv B)$, the half-distance between their convex hulls.
This is the geometric (convex-hull) form of the hard-margin support vector machine
\citep{bennett2000,burges1998}. We will reuse it verbatim, so we restate it as the lemma we lean on.

\begin{lemma}[Hard-margin = half convex-hull distance]
\label{lem:hull}
For finite sets $A,B \subset \R^m$ with $\conv A \cap \conv B = \varnothing$, the maximum over unit
vectors $n$ and biases $b$ of the separation margin
\[
  \min\big(\,\min_{a\in A}(n^\top a + b),\ \min_{b'\in B}(-n^\top b' - b)\,\big)
\]
equals $\tfrac12\dist(\conv A,\conv B)$, and is attained by the unit normal pointing along the
segment joining the two closest hull points, with the hyperplane placed at its midpoint.
\end{lemma}

\begin{theorem}[Invariant linear margin for a finite orthogonal group]
\label{thm:fg-margin}
Let a finite group $G$ act on $\R^d$ by orthogonal matrices $\{R_g\}_{g\in G}$. Define the fixed
subspace and the group-average (Reynolds) projector
\[
V^G := \{v : R_g v = v \ \text{for all } g\in G\}, \qquad
P_G := \frac{1}{|G|}\sum_{g\in G} R_g .
\]
A linear score $h(x) = w^\top x + b$ is exactly invariant, $h(R_g x) = h(x)$ for all $x$ and $g$, if
and only if $w \in V^G$. For finite labelled sets $X_+, X_- \subset \R^d$, the largest $\ell_2$
margin achievable by an exactly invariant linear classifier is
\[
\gamma_{\mathrm{inv}}
= \tfrac12 \dist\big(\conv(P_G X_+),\ \conv(P_G X_-)\big).
\]
\end{theorem}

\begin{proof}
\emph{Step 1 (which normals are invariant).}
Fix $g$. The identity $h(R_g x) = h(x)$ for \emph{all} $x$ means $w^\top R_g x + b = w^\top x + b$
for all $x$, i.e. $w^\top R_g x = w^\top x$ for all $x$. Subtracting, $(R_g^\top w - w)^\top x = 0$
for all $x \in \R^d$. A vector orthogonal to all of $\R^d$ is zero, so $R_g^\top w = w$.
\emph{(orthogonality of the action)} Because $R_g$ is orthogonal, $R_g^\top = R_g^{-1} = R_{g^{-1}}$,
and as $g$ ranges over the group so does $g^{-1}$; hence $R_g^\top w = w$ for all $g$ is the same
condition as $R_g w = w$ for all $g$, that is, $w \in V^G$. The converse is the same computation read
backwards: if $w \in V^G$ then $R_g^\top w = w$ gives $w^\top R_g x = w^\top x$.

\emph{Step 2 ($P_G$ is the orthogonal projector onto $V^G$).}
First $P_G$ is idempotent: \emph{(group closure)}
\[
P_G^2 = \frac{1}{|G|^2}\sum_{g,h\in G} R_g R_h
       = \frac{1}{|G|^2}\sum_{g,h} R_{gh}
       = \frac{1}{|G|^2}\sum_{h}\Big(\sum_g R_{gh}\Big)
       = \frac{1}{|G|^2}\sum_{h} |G|\,P_G = P_G,
\]
where for each fixed $h$ the inner sum over $g$ runs over all group elements $gh$, hence equals
$|G|\,P_G$. Next $P_G$ is symmetric: \emph{(orthogonality again)}
$P_G^\top = \frac{1}{|G|}\sum_g R_g^\top = \frac{1}{|G|}\sum_g R_{g^{-1}} = P_G$. A symmetric
idempotent is an orthogonal projector. Its range is $V^G$: if $v\in V^G$ then $P_G v = \frac1{|G|}
\sum_g R_g v = \frac1{|G|}\sum_g v = v$, so $V^G \subseteq \mathrm{range}(P_G)$; and for any $x$,
$R_h(P_G x) = \frac1{|G|}\sum_g R_{hg} x = P_G x$ by the same reindexing, so
$\mathrm{range}(P_G)\subseteq V^G$.

\emph{Step 3 (invariant classification = classification of the projection).}
Take any invariant $w \in V^G$. By Step~2, $w = P_G w$, and since $P_G$ is symmetric,
\[
w^\top x = (P_G w)^\top x = w^\top P_G x .
\]
So evaluating the invariant score at $x$ is the same as evaluating the \emph{same} normal $w$ at the
projected point $P_G x$, with the same $\|w\|_2$. Therefore an invariant linear classifier on
$X_+ \cup X_-$ is exactly an ordinary linear classifier on the projected sets $P_G X_+$ and
$P_G X_-$, and the signed margins match point by point. Conversely any linear classifier in the
projected space lifts (take its normal, which already lies in $V^G$) to an invariant one.

\emph{Step 4 (apply the hard-margin lemma).}
By Step~3, maximizing the invariant margin over original-space data is the same problem as
maximizing the ordinary linear margin over the projected data $P_G X_+$, $P_G X_-$.
Lemma~\ref{lem:hull} with $A = P_G X_+$ and $B = P_G X_-$ gives that this maximum equals
$\tfrac12\dist(\conv(P_G X_+),\conv(P_G X_-))$, attained by the closest-points normal and the
midpoint bias. This is the claim.
\end{proof}

\begin{remark}[Why the right object is a distance, not a one-sided gap]
A natural but wrong guess writes the invariant margin as a labelled gap such as
$\tfrac12(\min_{X_+}\fdc - \max_{X_-}\fdc)_+$, which is only correct when the positive class happens
to sit ``above'' the negative one along the fixed direction. The margin is symmetric in the labels:
it is the \emph{distance} between the two projected convex hulls, which is orientation-free and is
zero exactly when they overlap. This is the symmetric form proved above.
\end{remark}

\paragraph{Recovering Ge's $2/\sqrt d$ collapse.}
For cyclic shifts on $\R^d$ the fixed subspace is the single line $\mathrm{span}\{\one\}$ and
$P_G = \Piinv = ee^\top$ with $e = d^{-1/2}\one$ (Part~I). Then $P_G x = (e^\top x)e = \fdc(x)\,e$,
so the projected data lie on one line and are identified isometrically with the scalar DC values
$\fdc(x)$. The two projected hulls become two intervals $I_+, I_-$ on that line, and
Theorem~\ref{thm:fg-margin} reduces to
\[
\gamma_{\mathrm{inv}} = \tfrac12\dist(I_+, I_-).
\]
Take Ge's black/white dot example: $x_+ = e_j$ and $x_- = -e_j$ (a single bright pixel versus a
single dark pixel). Their DC values are $\fdc(\pm e_j) = \pm d^{-1/2}$, so the interval distance is
$2/\sqrt d$ and the invariant margin is $\gamma_{\mathrm{inv}} = 1/\sqrt d$. Yet \emph{without} the
invariance constraint a linear classifier may use the coordinate $x_j$ directly, separating
$\pm e_j$ with margin $1$. The invariant margin therefore decays like $1/\sqrt d$ in the image size
$d$ while the unconstrained margin stays constant: this is exactly the $2/\sqrt d$ collapse of
\citet{ge2021shift}, recovered as the cyclic specialization of the projection identity. The general
theorem says \emph{why} it happens: the invariance projects away every coordinate except the mean,
and a localized pattern almost entirely lives in the coordinates that get projected away.


% =====================================================================
\subsection{Q3 answered: the $\eta/L$ certificate, its missing converse, and the two-sided bracket}
\label{sub:sandwich}

\paragraph{The question this answers.}
Part~I gave the Lipschitz--margin certificate: if an invariant score $g$ has feature margin $\eta$
and Lipschitz constant $L$, then the $\ell_2$ robust radius satisfies $r_2(g,x) \ge \eta/L$. This is
the link from an invariance to an $\ell_2$ radius that no prior paper supplied. But a one-sided
certificate is a weak predictor unless we also know how much it can \emph{under}estimate the true
radius. So Q3 splits into three sub-questions. First, is the certificate ever loose by an unbounded
amount (no universal converse)? Second, under what condition can we recover an upper bound, turning
the certificate into a two-sided bracket that pins $r_2$ down? Third, when is the bracket exact?

\paragraph{Status.}
The lower bound is the textbook Lipschitz--margin certificate \citep{hein2017formal,tsuzuku2018lipschitz};
we do not claim it. The genuine contribution here is the \emph{upper arm}. \citet{tsuzuku2018lipschitz}
state a certificate chain whose top inequality (certificate $\le$ true radius) they say cannot be
guaranteed, and measure the resulting gap empirically as $1.8$--$2.4\times$; prior upper bounds on the
true radius were only empirical, read off by running a minimal-distortion attack
\citep{croce2020reliable}. Theorem~\ref{thm:sandwich} supplies that missing arm under an explicit,
checkable reachability condition (a co-Lipschitz bound along one path to the boundary). The linear
case (Corollary~\ref{cor:linear}) is the classical point-to-hyperplane distance
\citep{hein2017formal,cortes1995}, included as a tightness sanity check.

We first show the certificate has \emph{no} universal converse, so we know an upper arm must require
an extra hypothesis.

\begin{proposition}[No universal converse for $\eta/L$]
\label{prop:noconverse}
There is no universal constant $C$ such that every invariant classifier obeying the
Lipschitz--margin certificate also obeys $r_2(g,x) \le C\,\eta/L$. The failure occurs even when the
invariant feature depends only on the DC coordinate.
\end{proposition}

\begin{proof}
We exhibit a one-parameter family whose ratio $r_2/(\eta/L)$ grows without bound.
\emph{Step 1 (a DC-only invariant feature).}
Write $t = e^\top x = \fdc(x)$, the DC coordinate, which is shift-invariant. For an integer
$n \ge 1$ set the invariant scalar feature and classifier
\[
\Psi_n(x) = t^{2n+1}, \qquad g_n(x) = \Psi_n(x).
\]
Take the two data points $x_+ = e$ (label $+1$) and $x_- = -e$ (label $-1$), where $e = d^{-1/2}\one$
so that $t(x_\pm) = \pm 1$.

\emph{Step 2 (the feature margin).}
Then $g_n(e) = 1^{2n+1} = 1$ and $g_n(-e) = (-1)^{2n+1} = -1$. In the one-dimensional feature space
the two classes sit at $+1$ and $-1$, so the signed feature margin is $\eta = 1$.

\emph{Step 3 (the Lipschitz constant on the relevant segment).}
On the DC segment $\{t e : t \in [-1,1]\}$, the derivative of $\Psi_n$ in $t$ is $(2n+1)t^{2n}$, of
magnitude at most $2n+1$, attained at $t = \pm 1$. \emph{(maximum over $t$ of $|(2n+1)t^{2n}|$ on
$[-1,1]$)} Hence the Lipschitz constant is $L_n = 2n+1$, and the certificate gives
$\eta/L_n = 1/(2n+1)$.

\emph{Step 4 (the true radius).}
The decision boundary is $g_n = 0$, i.e. $t = 0$. The closest point to $x_+ = e$ with $t = 0$ is the
origin, at Euclidean distance $\|e - 0\|_2 = 1$; perturbing inside the anti-invariant subspace $\Vac$
(where $\one^\top\delta = 0$) does not change $t$ at all, so it cannot move the score. Therefore the
true robust radius is $r_2(g_n, e) = 1$.

\emph{Step 5 (the ratio is unbounded).}
Combining Steps 3 and 4,
\[
\frac{r_2(g_n, e)}{\eta/L_n} = \frac{1}{1/(2n+1)} = 2n+1 \xrightarrow[n\to\infty]{} \infty.
\]
No single constant $C$ can dominate $2n+1$ for all $n$, so no universal converse exists. Every
feature here depends only on $t$, proving the last sentence.
\end{proof}

The reason the certificate is loose in Proposition~\ref{prop:noconverse} is now visible: the feature
$t^{2n+1}$ is steep near the data ($L_n$ large) but \emph{flat} near the boundary $t=0$, so the
upper Lipschitz constant overstates how fast the score actually falls toward the boundary. Putting a
\emph{floor} under that fall, a co-Lipschitz bound (Remark~\ref{rem:terms}), is exactly what closes
the gap.

\begin{theorem}[Two-sided robust-radius bracket: a conditional converse]
\label{thm:sandwich}
Let $g$ be continuous and $L$-Lipschitz in $\ell_2$ on a convex set $S$ containing $x$, its nearest
boundary point, and the path $\gamma$ below; let $g$ be correct at $x$ with margin
$\mu = y\,g(x) > 0$, and write $\mathcal B = \{z : g(z) = 0\}$ for the decision boundary.
\begin{enumerate}[leftmargin=1.7em,label=\textup{(\roman*)}]
\item \emph{Lower arm:} $r_2(g,x) \ge \mu/L$.
\item \emph{Upper arm:} if there is a continuous path $\gamma:[0,T]\to S$ with $\gamma(0)=x$,
$g(\gamma(T))=0$, and a co-Lipschitz bound $|g(\gamma(t)) - g(x)| \ge \alpha\,\|\gamma(t)-x\|_2$ up to
the first boundary hit, then $r_2(g,x) \le \mu/\alpha$, and necessarily $\alpha \le L$.
\end{enumerate}
Hence, with the local bi-Lipschitz condition number $\kappa := L/\alpha \ge 1$,
\[
\frac{\mu}{L} \ \le\ r_2(g,x) \ \le\ \frac{\mu}{\alpha} = \kappa\cdot\frac{\mu}{L},
\]
so the certificate determines the robust radius up to the factor $\kappa$, and is exact
($r_2 = \mu/L$) if and only if $\kappa = 1$.
\end{theorem}

\begin{proof}
\emph{Lower arm.}
For any admissible $\delta$, the Lipschitz bound gives
$|g(x+\delta) - g(x)| \le L\,\|\delta\|_2$. If $\|\delta\|_2 < \mu/L$ then
$|g(x+\delta) - g(x)| < \mu = y\,g(x)$, so $y\,g(x+\delta) > 0$ and the sign cannot flip. Hence no
adversarial example exists inside radius $\mu/L$, i.e. $r_2(g,x) \ge \mu/L$. \emph{(this is the
Part~I certificate; restated for completeness)}

\emph{Upper arm.}
Evaluate the co-Lipschitz hypothesis at the first boundary hit $t = T$, where $g(\gamma(T)) = 0$:
\[
\mu = |g(x)| = |g(x) - g(\gamma(T))| \ \ge\ \alpha\,\|\gamma(T) - x\|_2,
\]
using $g(\gamma(T))=0$ in the middle and the co-Lipschitz floor on the right. Rearranging \emph{(divide
by $\alpha>0$)},
\[
\|\gamma(T) - x\|_2 \ \le\ \frac{\mu}{\alpha}.
\]
Now $\gamma(T) \in \mathcal B$ is a point where the score is zero, so an arbitrarily small further push
across $\mathcal B$ flips the prediction; the robust radius, being the \emph{distance to the boundary
set}, is at most the distance to this one boundary point: \emph{(infimum over boundary points is at
most one particular boundary point)}
\[
r_2(g,x) = \dist(x,\mathcal B) \ \le\ \|\gamma(T) - x\|_2 \ \le\ \frac{\mu}{\alpha}.
\]

\emph{The constants compare correctly.}
Along $\gamma$ the same secant $|g(\gamma(t)) - g(x)|$ is floored by $\alpha\|\gamma(t)-x\|_2$ and
capped by $L\|\gamma(t)-x\|_2$, so $\alpha\,\|\gamma(t)-x\|_2 \le L\,\|\gamma(t)-x\|_2$, hence
$\alpha \le L$ and $\kappa = L/\alpha \ge 1$. Chaining the two arms gives the bracket, and equality
$r_2 = \mu/L$ forces $\mu/\alpha = \mu/L$, i.e. $\kappa = 1$.
\end{proof}

\begin{remark}[When is the upper arm usable?]
The lower arm $r_2\ge\mu/L$ is free (it needs only the Lipschitz constant). The upper arm needs an
\emph{existence} witness: one path to the boundary with a positive co-Lipschitz slope $\alpha$. This is
not automatic on an arbitrary trained network, so the bracket is a \emph{conditional} converse, and
making it usable means producing such a path. Two cases where it is concrete: (i) for an explicit
invariant feature, such as the power-spectrum quadratic surrogate below, $\alpha$ is computable in
closed form on a ball, giving a numerical bracket; (ii) for a general differentiable score one can
descend along $-\nabla g$ from $x$ to the boundary and take $\alpha$ to be the smallest gradient norm
met along that path, which is a valid (if loose) co-Lipschitz floor. The honest gap between ``a
conditional converse exists'' and ``a tight bracket on a trained CNN'' is exactly the difficulty of
certifying a good $\alpha$, just as the lower arm's looseness is the difficulty of certifying a tight
$L$.
\end{remark}

\begin{corollary}[Linear invariant features: the bracket is exact]
\label{cor:linear}
If $g(x) = w^\top x$ with $w \in V^G$, then along the straight path $\gamma(t) = x - t\,w/\|w\|_2$ one
has $L = \alpha = \|w\|_2$, so $\kappa = 1$ and
\[
r_2(g,x) = \frac{\mu}{\|w\|_2} = \frac{\eta}{\|w\|_2},
\]
the exact point-to-hyperplane distance \citep{hein2017formal,cortes1995}.
\end{corollary}

\begin{proof}
For linear $g$, $|g(x) - g(z)| = |w^\top(x-z)| \le \|w\|_2\,\|x-z\|_2$ by Cauchy--Schwarz, with
equality exactly when $x-z$ is parallel to $w$. So the upper Lipschitz constant is $L = \|w\|_2$, and
along the straight path in direction $-w$ the inequality is an equality, giving co-Lipschitz constant
$\alpha = \|w\|_2$ as well. Thus $\kappa = 1$, and the path reaches the boundary at distance
$\mu/\|w\|_2$; Theorem~\ref{thm:sandwich} pins this as the exact radius. The margin $\mu$ equals the
feature margin $\eta$ for the max-margin separator, giving the stated form.
\end{proof}

\begin{remark}[Reconciling the no-converse counterexample with the bracket]
\label{rem:reconcile}
The bracket is not contradicted by Proposition~\ref{prop:noconverse}; it explains it. For
$\Psi_n(x) = t^{2n+1}$ at $x = e$, the straight DC path to the origin has the feature drop from $1$ to
$0$ over Euclidean distance $1$, so its secant co-Lipschitz constant is $\alpha = 1$, while the upper
Lipschitz on the segment is $L = 2n+1$. The bracket reads $[\mu/L,\,\mu/\alpha] = [1/(2n+1),\,1]$ and
correctly contains $r_2 = 1$, sitting at the $\mu/\alpha$ end. The failure of a \emph{universal}
converse is precisely that $\kappa = 2n+1$ is unbounded \emph{across the model family}, not that any
single model lacks a bracket. Note also the difference between the \emph{pointwise} lower-Lipschitz
$\inf_t|\Psi_n'|$, which vanishes at the boundary (a metric-subregularity failure
\citep{dontchev2009}) and so kills any derivative-based converse, and the \emph{secant}
co-Lipschitz above, which stays positive and yields the finite bracket.
\end{remark}

\begin{remark}[A computable instance]
For the power-spectrum surrogate of Section~\ref{sub:ps}, both arms are explicit: the upper Lipschitz
satisfies $L \le 2B$ on the ball $\|x\|_2 \le B$, and along the gradient-descent path the
co-Lipschitz constant is the minimum gradient norm $\alpha = \inf_{z\in\gamma}\|\nabla g(z)\|_2$,
strictly positive as long as the path avoids the single critical point of the quadratic. This gives
the first two-sided radius statement for an invariant-feature model, with measured condition number
$\kappa \approx 1.6$--$2.4$ and the attack-found radius lying strictly inside the bracket for every
seed.
\end{remark}


% =====================================================================
\subsection{Q2, deeper: a nonlinear invariant the linear one cannot match}
\label{sub:ps}

\paragraph{The question this answers.}
Theorem~\ref{thm:fg-margin} shows linear invariance is weak: for cyclic shifts it keeps only the
mean. The natural next step is a \emph{nonlinear} invariant feature that keeps more, while still
being computable by a standard layer. The classical choice is the power spectrum $\{|\hat x_k|^2\}_k$,
which is shift-invariant because a shift only rotates Fourier phases. The point of this subsection is
twofold: first, a single convolution--square--global-pool layer literally \emph{realizes} power-spectrum
coordinates, so this is not an exotic feature but what such a layer already computes; second, the power
spectrum is still blind to \emph{phase}, and a degree-three invariant, the bispectrum, separates
phase-coded classes that the power spectrum cannot, including DC-matched pairs that no shift-invariant linear feature separates either.

\paragraph{Status.}
Both facts are known-and-reproved. Power spectra, autocorrelation, and scattering invariants are
classical \citep{brunamallat2012}; the realization lemma is a low-level architectural identity (with
a corrected DFT normalization and real-filter handling). The bispectrum and its completeness-up-to-shift
are classical \citep{invariants_kakarala}. The only fresh angle is the \emph{application}: using a
higher-order invariant to defeat exactly the localized, phase-coded corner where Ge's forced trade-off
bites, and quantifying its cost as a larger Lipschitz constant, hence a smaller $\eta/L$ at equal margin.

\begin{lemma}[A conv--square--pool layer realizes the power spectrum]
\label{lem:ps}
Use the unitary DFT $\hat x_k = d^{-1/2}\sum_j x_j e^{-2\pi i jk/d}$. Define the unnormalised circular
cross-correlation $(w \star x)_s = \sum_j w_j x_{j+s \bmod d}$ and the feature
\[
\Phi_w(x) := \frac1d \sum_{s=0}^{d-1} (w\star x)_s^2 ,
\]
i.e. correlate $x$ with a filter $w$, square the response, and global-average-pool. Then for real
$w,x$,
\[
\Phi_w(x) = \sum_{k=0}^{d-1} |\hat w_k|^2\,|\hat x_k|^2 ,
\]
a nonnegative-weighted combination of the power-spectrum coordinates $|\hat x_k|^2$. Real filters can
isolate every nonredundant real power coordinate ($|\hat x_0|^2$, $|\hat x_{d/2}|^2$ when $d$ is even,
and each conjugate-pair sum $|\hat x_k|^2 + |\hat x_{d-k}|^2$).
\end{lemma}

\begin{proof}
\emph{Step 1 (DFT of a correlation).}
By the convolution/correlation theorem for the unitary DFT, the transform of $w\star x$ is
$\widehat{w\star x}_k = \sqrt{d}\,\overline{\hat w_k}\,\hat x_k$ \emph{(the $\sqrt d$ is the unitary
normalization; the conjugate is because correlation flips one argument)}.

\emph{Step 2 (Parseval turns the pooled square into a frequency sum).}
The pooling sum is a squared $\ell_2$ norm in space, and the unitary DFT preserves $\ell_2$ norms
(Parseval):
\[
\Phi_w(x) = \frac1d\sum_s |(w\star x)_s|^2 = \frac1d\sum_k |\widehat{w\star x}_k|^2 .
\]

\emph{Step 3 (substitute Step 1).}
Insert $\widehat{w\star x}_k = \sqrt d\,\overline{\hat w_k}\hat x_k$ and take squared modulus:
\[
\Phi_w(x) = \frac1d\sum_k \big|\sqrt d\,\overline{\hat w_k}\hat x_k\big|^2
          = \frac1d\sum_k d\,|\hat w_k|^2|\hat x_k|^2
          = \sum_k |\hat w_k|^2\,|\hat x_k|^2 ,
\]
the $\sqrt d$ squaring to $d$ and cancelling the $1/d$. This is the claimed identity.

\emph{Step 4 (which coordinates real filters reach).}
For real $w$, conjugate symmetry $\hat w_{d-k} = \overline{\hat w_k}$ forces $|\hat w_{d-k}|^2 =
|\hat w_k|^2$, so a real filter weights $|\hat x_k|^2$ and $|\hat x_{d-k}|^2$ \emph{equally}; it cannot
separate a conjugate pair. But it \emph{can} isolate a pair: choose $\hat w$ supported only on
$\{k, d-k\}$ with conjugate values (which makes $w$ real), giving exactly the pair sum
$|\hat x_k|^2 + |\hat x_{d-k}|^2$; the self-conjugate coordinates $k=0$ and (for even $d$) $k=d/2$ are
isolated singly. Linear combinations of these give the stated real span.
\end{proof}

So one standard layer escapes the one-dimensional invariant linear family of
Theorem~\ref{thm:fg-margin} and sees the whole power spectrum. But the power spectrum still discards
phase, and that is exactly what Ge's localized example exploits.

\begin{proposition}[The bispectrum separates phase-coded classes the power spectrum cannot]
\label{prop:bispectrum}
The bispectrum $B_{k_1,k_2}(x) = \hat x_{k_1}\hat x_{k_2}\overline{\hat x_{k_1+k_2}}$ is
shift-invariant and is realizable as a degree-three convolution--product--pool feature. There exist
classes with \emph{identical} power spectra that a single bispectrum coordinate separates: for Ge's
dot, $x_+ = e_j$ and $x_- = -e_j$ have $|\hat x_k|^2 = 1/d$ for all $k$, so every power-spectrum
feature is blind to the difference, yet
\[
B_{0,0}(e_j) = +d^{-3/2}, \qquad B_{0,0}(-e_j) = -d^{-3/2},
\]
separating them with margin $d^{-3/2}$. (The linear DC invariant also separates this particular pair,
through $\fdc(\pm e_j) = \pm d^{-1/2}$; the sharper claim, proved in Step~3, is that the bispectrum
separates \emph{any} two non-shift-equivalent orbits, including pairs with identical DC and identical
power spectrum, which no shift-invariant linear feature can do.) This separability comes at the cost
of a larger Lipschitz constant, $O(B^2)$ versus $O(B)$ for the power spectrum on the ball
$\|x\|_2 \le B$, i.e. a smaller certified $\eta/L$ at equal margin.
\end{proposition}

\begin{proof}
\emph{Step 1 (shift-invariance of the bispectrum).}
Under $x \mapsto S_s x$ each Fourier coordinate picks up a unit phase, $\hat x_k \mapsto \omega^{sk}\hat
x_k$ with $\omega = e^{-2\pi i/d}$ (the forward DFT convention of Part~I; the shift theorem
$\mathrm{DFT}(R_\tau x)_k = \omega^{k\tau}\hat x_k$). The three factors of $B_{k_1,k_2}$ acquire phases $\omega^{s k_1}$,
$\omega^{s k_2}$, and (from the conjugate) $\omega^{-s(k_1+k_2)}$, whose product is
$\omega^{s(k_1 + k_2 - k_1 - k_2)} = \omega^0 = 1$. The phases cancel, so $B_{k_1,k_2}$ is invariant.
It is a cubic form in $x$, hence realizable by circular correlations, a pointwise product, and global
average pooling (a degree-three analogue of Lemma~\ref{lem:ps}).

\emph{Step 2 (the dot has a flat power spectrum).}
For $x = e_j$ with the unitary DFT, $\hat x_k = d^{-1/2}\omega^{jk}$ where $\omega = e^{-2\pi i/d}$, so
$|\hat x_k|^2 = 1/d$ for every $k$ (the phase has modulus one). The same holds for $-e_j$. The two
power spectra are identical, so by Lemma~\ref{lem:ps} no power-spectrum feature distinguishes them.
A shift-invariant \emph{linear} feature sees only the DC value (Theorem~\ref{thm:fg-margin}); for this
particular pair the DC value $\fdc(\pm e_j) = \pm d^{-1/2}$ does differ, so the linear invariant
separates $\pm e_j$, but only with the weak $1/\sqrt d$ margin of the Ge collapse, and it would be
\emph{completely} blind to any DC-matched pair such as two shifts of a phase-coded signal.

\emph{Step 3 (a bispectrum coordinate separates them).}
Take $k_1 = k_2 = 0$. Then $B_{0,0}(x) = \hat x_0 \hat x_0 \overline{\hat x_0} = \hat x_0\,|\hat x_0|^2$.
For $e_j$, $\hat x_0 = d^{-1/2}$ and $|\hat x_0|^2 = 1/d$, so $B_{0,0}(e_j) = d^{-1/2}\cdot d^{-1} =
d^{-3/2}$. For $-e_j$ the sign of $\hat x_0$ flips, giving $-d^{-3/2}$. The single coordinate
$B_{0,0}$ takes opposite signs on the two classes, separating them with margin $d^{-3/2}$. More
generally, when all Fourier coefficients are nonzero the bispectrum determines the signal up to a
global shift \citep{invariants_kakarala}, so it separates any two shift-orbits that are not
shift-equivalent.

\emph{Step 4 (the Lipschitz cost).}
A cubic form has gradient of order $\|x\|_2^2$, so on the ball $\|x\|_2 \le B$ the bispectrum is
locally Lipschitz with constant $O(B^2)$, versus $O(B)$ for the quadratic power spectrum. At equal
margin a larger $L$ means a smaller certified $\eta/L$, which is the stated trade: the higher-order
invariant buys separability of the phase-coded corner at the price of a smaller certificate.
\end{proof}


% =====================================================================
\subsection{Q5 answered, part one: on orbit-closed data invariance is margin-free}
\label{sub:selfsym}

\paragraph{The question this answers.}
The reproduction's quadrant diagnostic (Part~II, Act~II) plotted models on a
consistency-versus-robustness plane and could not say \emph{why} a model lands where it lands, nor
whether building in invariance helps. Q5 asks for
the optimization story. The first half of the answer is a surprising \emph{negative} result: when the
data is orbit-closed (every shift of a training point is also a training point with the same label),
imposing exact shift-invariance does not change the max-margin value at all. Gradient descent on the
unconstrained network already self-symmetrizes for free, so there is no margin to be gained by tying
weights. This explains the experiment in which unconstrained and shift-tied networks reach essentially
the same robustness on generic orbit data.

\paragraph{Status.}
A medium-strength novel candidate. It is the nonlinear-ReLU, orbit-closed analogue of the linear
steerable/augmentation equivalence of \citet{chenzhu2023}; the linear case is theirs, the ReLU
min-norm accounting below is the new content. We use one architectural identity from Part~I: a
convolution-plus-global-pool unit is the same function as a two-layer net whose hidden weights are
tied across the shift orbit of one filter (the ``orbit bundle''). We also use the standard fact that a
positively homogeneous ReLU unit $a\,\sigma(w^\top x)$ can be rescaled $(a,w)\mapsto(a/c,\,cw)$ without
changing its function, so its minimum-norm representation costs $|a|\,\|w\|$ (taking $c^2 = |a|/\|w\|$
balances $a^2$ and $\|w\|^2$).

\begin{theorem}[Invariance is margin-free on orbit-closed data]
\label{thm:selfsym}
Let the data be orbit-closed with labels constant on orbits. For two-layer ReLU networks
$f(x) = \sum_r a_r \sigma(w_r^\top x)$ (positively homogeneous, no bias), the minimum parameter norm
$\tfrac12\sum_r(a_r^2 + \|w_r\|^2)$ achieving unit margin is the \emph{same} for the unconstrained
class and for the shift-invariant (tied orbit-bundle) class. Equivalently, imposing exact
shift-invariance changes neither the max-margin value nor the achievable margin.
\end{theorem}

\begin{proof}
Write $C^\star_{\mathrm{unc}}$ and $C^\star_{\mathrm{inv}}$ for the two minimum costs. We show
$C^\star_{\mathrm{inv}} \ge C^\star_{\mathrm{unc}}$ and $C^\star_{\mathrm{inv}} \le
C^\star_{\mathrm{unc}}$.

\emph{Step 1 ($\ge$, the easy direction).}
The tied (invariant) class is a \emph{subset} of the unconstrained class: every invariant network is
in particular some network. Minimizing the same cost over a smaller feasible set cannot give a smaller
value, so $C^\star_{\mathrm{inv}} \ge C^\star_{\mathrm{unc}}$.

\emph{Step 2 (the symmetrization map).}
Let $f^\star = \sum_r a_r\sigma(w_r^\top x)$ be any unconstrained network with margin $\ge 1$. Define
its orbit average
\[
f_{\mathrm{sym}}(x) = \frac1d\sum_{s=0}^{d-1} f^\star(S_s x).
\]
Reindexing the pooling sum shows $f_{\mathrm{sym}}$ is shift-invariant, and by the orbit-bundle
identity it lies in the tied class.

\emph{Step 3 (symmetrization preserves the margin).}
Fix a training point $(x_i, y_i)$. Because the data is orbit-closed with orbit-constant labels, each
shifted point $S_s x_i$ is itself a training point with the same label $y_i$, so $y_i f^\star(S_s x_i)
\ge 1$ for every $s$. Averaging over $s$ \emph{(average of numbers each $\ge 1$ is $\ge 1$)},
\[
y_i f_{\mathrm{sym}}(x_i) = \frac1d\sum_s y_i f^\star(S_s x_i) \ge \frac1d\sum_s 1 = 1 .
\]
So $f_{\mathrm{sym}}$ still has unit margin.

\emph{Step 4 (the min-norm cost of one unit).}
For a single ReLU unit $a\,\sigma(w^\top x)$, positive homogeneity of $\sigma$ means
$a\,\sigma(w^\top x) = (a/c)\,\sigma((cw)^\top x)$ for any $c>0$, so we may rescale freely. The cost
$\tfrac12((a/c)^2 + \|cw\|^2)$ is minimized over $c$ by the arithmetic--geometric-mean balance at
$c^2 = |a|/\|w\|$, giving minimal cost $|a|\,\|w\|$. \emph{(this is the standard min-norm value of one
homogeneous unit)}

\emph{Step 5 (the bundle realizing $f_{\mathrm{sym}}$ costs no more than $f^\star$).}
Expanding the average, $f_{\mathrm{sym}}(x) = \frac1d\sum_r\sum_s a_r\,\sigma(w_r^\top S_s x)
= \frac1d\sum_r\sum_s a_r\,\sigma((S_{-s}w_r)^\top x)$, using that $S_s$ is orthogonal so
$w_r^\top S_s x = (S_s^\top w_r)^\top x = (S_{-s}w_r)^\top x$. This is a network whose units are
$\{(a_r/d,\ S_{-s}w_r)\}_{r,s}$. By Step~4 its min-norm cost is
\[
\sum_r\sum_s \frac{|a_r|}{d}\,\|S_{-s}w_r\|
= \sum_r\sum_s \frac{|a_r|}{d}\,\|w_r\|
= \sum_r |a_r|\,\|w_r\| ,
\]
where the shift $S_{-s}$ preserves the norm (orthogonality) and the inner sum over the $d$ shifts of
the factor $1/d$ collapses. But $\sum_r |a_r|\,\|w_r\|$ is exactly the min-norm cost of the original
$f^\star$ by Step~4. Hence $f_{\mathrm{sym}}$ is realized in the tied class at cost no larger than the
cost of $f^\star$.

\emph{Step 6 (conclude).}
Taking $f^\star$ to be a unconstrained \emph{minimizer}, Step~5 produces a feasible invariant network
of the same cost, so $C^\star_{\mathrm{inv}} \le C^\star_{\mathrm{unc}}$. With Step~1, the two are
equal. There is therefore an invariant max-margin solution of the same norm: invariance is
margin-free here.
\end{proof}

\begin{remark}[What it does and does not give]
\label{rem:selfsym}
Theorem~\ref{thm:selfsym} equalizes the \emph{margin}; it does not say which max-margin solution
gradient descent selects, and the certified radius $\eta/L$ depends on the realized Lipschitz term
$L$, not the margin alone. So invariance can still change $\eta/L$ by changing $L$ even when it cannot
change $\eta$. That residual selection question, on data where the orbit-closed premise fails, is the
content of the next subsection and its open conjecture.
\end{remark}


% =====================================================================
\subsection{Q5 answered, part two: in the lazy regime invariance lowers the Lipschitz constant}
\label{sub:lazy}

\paragraph{The question this answers.}
Theorem~\ref{thm:selfsym} says invariance cannot help \emph{the margin} on orbit-closed data. But on
data with near-orthogonal cluster geometry, experiment shows it \emph{does} help $\eta/L$, and the
gap is largest in the lazy (NTK) training regime. The remaining question is whether that advantage is
provable. The answer is yes in the lazy, smooth-activation regime: there the trained network is the
minimum-RKHS-norm interpolant of a fixed kernel, orbit-averaging is an orthogonal projection in that
RKHS, and a projection cannot increase the norm, hence cannot increase the Lipschitz constant, at
equal margin. So invariance buys robustness through a smaller $L$, exactly the lever
Remark~\ref{rem:selfsym} left open.

\paragraph{Status.}
A medium-strength novel candidate, with an honestly limited scope. It proves the cluster conjecture
\emph{only} in the lazy, smooth-activation regime, by combining the lazy/NTK fixed-kernel picture
\citep{jacot2018,chizat2019} with the RKHS orbit-averaging projection of \citet{elesedy2021} and
the invariant-kernel symmetrization of \citet{haasdonk2007}. The smooth-activation caveat is
real and we state it: the bare two-layer ReLU NTK feature map is \emph{not} Lipschitz
\citep{biettimairal2019}, so the Lipschitz step needs a smooth activation; the rich (small-init)
and bare-ReLU trajectories stay open. From Part~I we use the lazy regime (a wide network behaves like
a fixed kernel), the RKHS norm and its Lipschitz bound $|f(x)-f(y)| \le \|f\|_K\,L_\Phi\,\|x-y\|_2$,
and orthogonal projection (a projection is non-expansive and gives a Pythagorean norm split).

\begin{theorem}[Invariance lowers the Lipschitz constant in the lazy, smooth-activation regime]
\label{thm:lazy}
Work in the lazy/NTK regime, where each trained network equals the minimum-RKHS-norm interpolant of
its fixed neural tangent kernel $K$ \citep{jacot2018,chizat2019}, with a smooth activation so that the
tangent feature map $\Phi$ is $L_\Phi$-Lipschitz on the data ball \citep{biettimairal2019}. Let the
data be orbit-closed with orbit-constant labels and let $K$ be $G$-invariant, so the Reynolds average
$\Pi_G f = \frac1{|G|}\sum_{g\in G} f(g\cdot\,)$ is defined on the RKHS $\mathcal H_K$ and the
shift-tied bundle realizes exactly the invariant subspace $\mathcal H^G = \Pi_G\mathcal H_K$ with the
restricted norm \citep{haasdonk2007}. Let $f_{\mathrm{unc}}$ be the unconstrained min-norm
interpolant and $f_{\mathrm{inv}}$ the min-norm interpolant in $\mathcal H^G$. Then
\begin{enumerate}[leftmargin=1.7em,label=\textup{(\roman*)}]
\item $\Pi_G$ is the orthogonal projection of $\mathcal H_K$ onto $\mathcal H^G$, non-expansive, with
$\|f\|_K^2 = \|\Pi_G f\|_K^2 + \|(I-\Pi_G)f\|_K^2$ \citep{elesedy2021};
\item $\Pi_G f_{\mathrm{unc}}$ still interpolates with the same unit margin, so
$\|f_{\mathrm{inv}}\|_K \le \|\Pi_G f_{\mathrm{unc}}\|_K \le \|f_{\mathrm{unc}}\|_K$;
\item hence $L_{\mathrm{inv}} \le L_\Phi\,\|f_{\mathrm{inv}}\|_K \le L_\Phi\,\|f_{\mathrm{unc}}\|_K$,
and at the common unit margin the invariant model has the weakly larger certified ratio $\eta/L$. The
gap is the discarded non-invariant energy $\|(I-\Pi_G)f_{\mathrm{unc}}\|_K^2$ (the orbit-variance),
strict unless $f_{\mathrm{unc}}$ is already invariant.
\end{enumerate}
\end{theorem}

\begin{proof}
\emph{(i) $\Pi_G$ is an orthogonal projection.}
Because $K$ is $G$-invariant, the averaging operator $\Pi_G$ maps $\mathcal H_K$ into itself, and it
is idempotent ($\Pi_G^2 = \Pi_G$, the same group-reindexing as in Theorem~\ref{thm:fg-margin}, Step~2)
and self-adjoint on $\mathcal H_K$. An idempotent self-adjoint operator has spectrum in $\{0,1\}$ and
is the orthogonal projection onto its range, which is the invariant subspace $\mathcal H^G$. For an
orthogonal projection, every $f$ splits as $f = \Pi_G f + (I-\Pi_G)f$ into orthogonal pieces, so by
the Pythagorean theorem $\|f\|_K^2 = \|\Pi_G f\|_K^2 + \|(I-\Pi_G)f\|_K^2$
\citep{elesedy2021}; dropping the second nonnegative term gives non-expansiveness
$\|\Pi_G f\|_K \le \|f\|_K$.

\emph{(ii) the projected interpolant is still feasible, with no larger norm.}
Orbit-constant labels mean $y_i f_{\mathrm{unc}}(g\cdot x_i) \ge 1$ for every $g$ (each $g\cdot x_i$
is a training point of the same label, interpolated by $f_{\mathrm{unc}}$). Averaging over $g$,
\[
y_i\,(\Pi_G f_{\mathrm{unc}})(x_i) = \frac1{|G|}\sum_g y_i f_{\mathrm{unc}}(g\cdot x_i) \ge 1,
\]
so $\Pi_G f_{\mathrm{unc}} \in \mathcal H^G$ is feasible for the invariant min-norm problem. Since
$f_{\mathrm{inv}}$ is the \emph{minimizer} of that problem, $\|f_{\mathrm{inv}}\|_K \le
\|\Pi_G f_{\mathrm{unc}}\|_K$, and by (i) this is $\le \|f_{\mathrm{unc}}\|_K$.

\emph{(iii) smaller norm gives smaller Lipschitz, hence larger $\eta/L$.}
For any RKHS function, the reproducing property and Cauchy--Schwarz give
\[
|f(x) - f(y)| = |\langle f,\ \Phi(x) - \Phi(y)\rangle_K|
\le \|f\|_K\,\|\Phi(x) - \Phi(y)\|_K
\le \|f\|_K\,L_\Phi\,\|x - y\|_2,
\]
the last step using that $\Phi$ is $L_\Phi$-Lipschitz on the data ball (smooth activation)
\citep{biettimairal2019}. So a valid Lipschitz constant is $L \le L_\Phi\,\|f\|_K$. Applying this to
each interpolant and inserting (ii), $L_{\mathrm{inv}} \le L_\Phi\,\|f_{\mathrm{inv}}\|_K \le
L_\Phi\,\|f_{\mathrm{unc}}\|_K$. The two models share the same unit margin $\eta$, so the invariant
model's certified $\eta/L_{\mathrm{inv}}$ is at least the unconstrained model's. The size of the gap is
controlled by the discarded energy $\|(I-\Pi_G)f_{\mathrm{unc}}\|_K^2$ from the split in (i): it is
zero exactly when $f_{\mathrm{unc}}$ is already invariant, and strictly positive otherwise.
\end{proof}

\begin{remark}[The ReLU caveat, stated honestly]
\label{rem:relu}
Steps (i)--(ii) are activation-independent: they are pure RKHS geometry. Only step (iii) uses
smoothness. The bare two-layer ReLU NTK feature map is \emph{not} Lipschitz; its RKHS contains
unit-norm functions of unbounded Lipschitz constant \citep{biettimairal2019}, so for ReLU one must
either take a smooth activation or carry $\sup\|\nabla\Phi\|$ on the data ball explicitly. Within that
scope the theorem proves the cluster conjecture, exactly in the regime where the project's
experiments report the largest gap (orbit-variance $\approx 0.3$, tied-over-unconstrained $\eta/L$ up
to $4.6\times$). The rich (small-init) regime and the bare-ReLU trajectory remain open.
\end{remark}


% =====================================================================
\subsection{The open frontier: the rich-regime cluster conjecture and how to attack it}
\label{sub:conjecture}

Theorem~\ref{thm:lazy} closes the lazy half. The rich (feature-learning, small-init) ReLU half is
genuinely open. We state it as a conjecture, then, in the spirit of a primer, list the candidate
proof paths and the Part~I tools each one needs, so a collaborator can pick one up.

\begin{conjecture}[Invariance helps on cluster geometry]
\label{conj:cluster}
On near-orthogonal cluster-geometry orbit data, the tied-invariant two-layer ReLU network's gradient
flow reaches a strictly larger $\eta/L$ than the unconstrained network's, with the gap growing as
initialization moves toward the lazy regime. Equivalently, weight sharing plus pooling re-points the
implicit bias toward the robust invariant solution.
\end{conjecture}

Why it is hard, in one line each: the orbit data is rank-structured, \emph{not} the near-orthogonal
$\Omega(d)$-cluster model whose non-robust selection is known, so that conclusion does not transfer;
the tied program has its own constrained KKT conditions that must be written, not inherited;
homogeneity gives convergence to \emph{some} KKT point but not \emph{which} one; and $\eta/L$ depends on
the realized Lipschitz term, so one must track the input-gradient norm along the trajectory, not just
the functional margin.

\paragraph{Candidate proof paths and the tools each needs.}
\begin{enumerate}[leftmargin=1.6em]
\item \textbf{Frequency-domain reduction (one orbit per class).} Replace the ReLU head by the
conv--square--pool surrogate so the features are power-spectrum coordinates (Lemma~\ref{lem:ps}); the
invariant margin program then becomes a linear/quadratic max-margin in Fourier-magnitude space,
possibly solvable in closed form, and one compares its $\eta/L$ to the unconstrained KKT point on the
same data, then swaps the surrogate for ReLU.
\emph{Tools:} the \emph{Fourier/power-spectrum invariants} and the \emph{convex-hull max-margin}
geometry, both introduced in the invariance and learning clusters of Part~I.
\item \textbf{Explicit KKT in Fourier coordinates.} Use the circulant/DFT diagonalization of the shift
action so that frequencies decouple, write the unconstrained and constrained KKT systems per
frequency, and solve frequency by frequency.
\emph{Tools:} the \emph{implicit-bias/KKT machinery} for homogeneous networks
\citep{lyuli2020,jitelgarsky2019,soudry2018} from the learning cluster of Part~I, plus the
\emph{per-irreducible (DFT) decomposition} of the shift group from the invariance cluster.
\item \textbf{Rich versus lazy contrast.} Run the same construction at small init (rich,
feature-learning) and at NTK scale (lazy), and show the selected $\eta/L$ differs, with the invariant
constraint mattering most in the rich regime.
\emph{Tools:} the \emph{NTK/lazy-versus-rich} distinction from the learning cluster of Part~I; the
lazy end is already pinned by Theorem~\ref{thm:lazy}.
\item \textbf{Bound the two ends separately.} Lower-bound the invariant network's achieved margin by
the projected convex-hull margin of Theorem~\ref{thm:fg-margin} in feature space, upper-bound the
unconstrained network's $\eta/L$ via a feature-averaging direction, then compare.
\emph{Tools:} the \emph{projection-margin accounting} of Section~\ref{sub:fg-margin} and the
\emph{feature-averaging implicit-bias} results \citep{li2025feature} from the learning cluster.
\end{enumerate}
Each route also needs the supporting setup lemmas, the bare $\eta/L$ lower bound, the orbit-rank and
orbit-bundle identities, and the KKT scaffolding; we have stated here only the load-bearing results
and leave the routine scaffolding to the reader who picks up a route.


% =====================================================================
\subsection{Tying theory to the benchmark: the capacity-matched $\eta/L$ decomposition}
\label{sub:empirical}

\paragraph{The question this answers.}
Q1 and Q4: if shift-consistency does not order robustness, what single measurable quantity does, and
does the answer survive a strong attack at matched capacity? The empirical headline is that $\eta/L$,
computed in the invariant feature space, is that quantity, and the theory above predicts \emph{which}
of its two parts a given architecture moves.

\paragraph{Status.}
This is the best empirical candidate, and its framing is the honest one of
Remark~\ref{rem:calibration}: not ``does invariance help?'' but ``when capacity is controlled, which
term, margin $\eta$ or Lipschitz $L$, does a given invariance move?'' Robustness is measured by a
standardized strong attack (AutoAttack) at matched parameter counts \citep{croce2020reliable,robustbench2020},
with capacity controlled as in architecture-robustness studies \citep{robarch2023}.

\paragraph{The findings, briefly.}
Across capacity-matched arms (standard, anti-aliased, exact-cyclic, shift-augmented, identical
parameter counts), $\eta/L$ versus the $\ell_2$ robust radius tracks at Pearson $0.998$, while
shift-consistency \emph{anti}-predicts robustness. The decomposition attributes mechanism exactly as
the theory says it should: anti-aliasing raises $\eta/L$ by lowering the Lipschitz term $L$
(Theorem~\ref{thm:lazy}'s lever), and exact cyclic invariance lowers $\eta/L$ by shrinking the margin
$\eta$ (the real-data form of the Ge collapse of Section~\ref{sub:fg-margin}). Under adversarial
training, where AutoAttack becomes informative with no gradient masking, shift-consistency still
anti-predicts robustness (Pearson $-0.88$), and the \emph{threat-matched} ratio
$\eta/\|\nabla M\|_1$, using the $\ell_\infty$ dual norm, predicts AutoAttack robust accuracy at
Pearson $0.90$, whereas the mismatched $\ell_2$ ratio does not. The two laws replicate on MNIST and
Fashion-MNIST and under an $\ell_2$ threat. This is the validated margin/Lipschitz decomposition that
turns the contradictory ``helps or hurts'' literature into a single sharp question with a measurable
answer.


% =====================================================================
% CITATION KEYS USED (all already present in primer.bib unless flagged NEW):
%   bennett2000          (existing) — convex-hull / SVM duality, Lemma hull
%   burges1998           (existing) — SVM tutorial, Lemma hull
%   cortes1995           (existing) — SVM / point-to-hyperplane, Cor linear
%   hein2017formal       (existing) — Lipschitz-margin / exact linear radius
%   tsuzuku2018lipschitz (existing) — certificate chain whose upper arm is completed
%   croce2020reliable    (existing) — empirical (attack) upper bound replaced
%   soudry2018           (existing) — implicit bias, linear separable
%   lyuli2020            (existing) — homogeneous-net KKT margin
%   jitelgarsky2019      (existing) — implicit bias / KKT
%   jacot2018            (existing) — NTK / lazy
%   chizat2019           (existing) — lazy training
%   li2025feature        (existing) — feature-averaging implicit bias
%   invariants_kakarala  (existing) — bispectrum completeness
%
%   ---- NEW KEYS (NOT yet in primer.bib — must be added) ----
%   ge2021shift               NEW — Ge et al., Shift Invariance Can Reduce Adversarial Robustness, NeurIPS 2021
%                              (the base result; cyclic 2/sqrt(d) collapse, projected-margin special case)
%   brunamallat2012      NEW — Bruna & Mallat, Invariant Scattering Convolution Networks, 2012/2013
%                              (classical translation-invariant nonlinear features)
%   chenzhu2023          NEW — Chen & Zhu, implicit bias of equivariant steerable networks, 2023
%                              (linear steerable/augmentation equivalence extended by thm:selfsym)
%   elesedy2021    NEW — Elesedy, kernel/RKHS orbit-averaging projection lemma, 2021
%                              (orthogonal-projection + Pythagorean split for thm:lazy)
%   haasdonk2007 NEW — Haasdonk & Burkhardt, invariant kernels by symmetrization, 2007
%                              (invariant tangent kernel = symmetrized kernel)
%   biettimairal2019  NEW — Bietti & Mairal, group invariance / smoothness of (convolutional) NTK, 2019
%                              (smooth-activation Lipschitz; bare-ReLU NTK non-Lipschitz caveat)
%   jacot2018, chizat2019 are the existing lazy/NTK keys reused here; the TR uses Chizat2018 but the
%   primer scheme is chizat2019 (matched to primer.bib).
%   dontchev2009 NEW — Dontchev & Rockafellar, metric (sub)regularity, 2009
%                              (cited once in rem:reconcile for the subregularity-failure aside; OPTIONAL —
%                               drop the parenthetical if the key is not added)
%   croce2020reliable       NEW — Croce & Hein, AutoAttack, 2020 (standardized strong attack)
%                              NOTE: distinct from existing croce2020reliable; if they are the same entry
%                              in primer.bib, replace croce2020reliable with croce2020reliable.
%   robustbench2020      NEW — Croce et al., RobustBench, 2020/2021 (standardized robustness benchmark)
%   robarch2023          NEW — Peng et al. (RobArch), capacity-controlled architecture robustness, 2023
% =====================================================================


%==============================================================================
\bibliographystyle{plainnat}
\bibliography{primer}

\end{document}
